% Linear Transformation — Further Mathematics Upper Sixth — Cours signé OnBuch+
% Official MINESEC syllabus. Compiler avec : tectonic FMA13-linear-transformation.tex
\newcommand{\DOCMATIERE}{Further Mathematics}
\newcommand{\DOCNIVEAU}{Upper Sixth}
\newcommand{\DOCMODULE}{Upper Sixth — Further mathematics}
\newcommand{\DOCLECON}{13}
\newcommand{\DOCTITRE}{Linear Transformation}
% OnBuch+ — préambule « cours signés » ENRICHI (compilé avec tectonic / XeLaTeX)
% Système visuel « Pop » d'OnBuch+ : fond crème (jamais blanc), bordures encre
% épaisses, ombres « dures » décalées, titres Archivo Black en capitales.
% Version MATHÉMATIQUES : graphes (pgfplots/tikz), boîtes pédagogiques
% (définition, propriété, méthode, attention, le savais-tu, exercice résolu,
% exercice, TP, bilan), page de garde et signature OnBuch+ ancrée en bas.
%
% Macros attendues dans main.tex AVANT \input{preamble} :
%   \DOCMATIERE  \DOCNIVEAU  \DOCMODULE  \DOCLECON (numéro)  \DOCTITRE
\documentclass[11pt]{article}

\usepackage{fontspec}
\usepackage[english]{babel}
\usepackage[a4paper,margin=2cm,top=3.4cm,bottom=2.8cm,headheight=1.4cm,footskip=1.3cm]{geometry}
\usepackage{amsmath,amssymb,amsfonts}
\usepackage{mathtools} % \xmapsto, \coloneqq... souvent utilisés par les modèles
\usepackage{cancel} % simplifications barrées (bilans de forces)
% \abs{z} (module, valeur absolue) et \norm{u} (norme), utilisés par les modèles
\providecommand{\abs}{}\let\abs\relax\DeclarePairedDelimiter{\abs}{\lvert}{\rvert}
\providecommand{\norm}{}\let\norm\relax\DeclarePairedDelimiter{\norm}{\lVert}{\rVert}
\usepackage[table]{xcolor} % [table] : fournit aussi \rowcolors (alternance de lignes)
\usepackage{pagecolor}
\usepackage{tikz}
\usetikzlibrary{arrows.meta,positioning,calc,shapes.geometric,shapes.misc,shapes.arrows,decorations.pathmorphing,decorations.pathreplacing,decorations.markings,patterns,shadows,mindmap,trees,fit,backgrounds,matrix,intersections,angles,quotes}
\usepackage{pgfplots}
\usepackage[version=4]{mhchem} % équations nucléaires \ce{^{235}_{92}U}
\usepackage[european,siunitx]{circuitikz} % schémas électriques
\pgfplotsset{compat=1.18}
\usepgfplotslibrary{fillbetween,groupplots} % aires entre courbes (calcul intégral)
\usepackage{enumitem}
\usepackage{fancyhdr}
% [most] charge toutes les bibliothèques tcolorbox (listings, minted, raster,
% documentation...) et gonfle inutilement la mémoire TeX sur un document long
% (~300 boîtes) : on ne charge que ce qui est réellement utilisé.
\usepackage[skins,breakable]{tcolorbox}
\usepackage{graphicx}
\usepackage{colortbl}
\usepackage{booktabs}
\usepackage{tabularx}
\usepackage{diagbox} % case d'en-tête diagonale des tableaux à double entrée
% Colonnes extensibles centrée / gauche / droite (C, L, R), souvent utilisées par les modèles
\newcolumntype{C}{>{\centering\arraybackslash}X}
\newcolumntype{L}{>{\raggedright\arraybackslash}X}
\newcolumntype{R}{>{\raggedleft\arraybackslash}X}
\usepackage{array}
\usepackage{multicol}
\usepackage{siunitx}
% parse-numbers=false : accepte aussi \SI{2,5 \times 10^{-3}}{...}
\sisetup{locale=UK,output-decimal-marker={.},group-minimum-digits=4,parse-numbers=false}
\DeclareSIUnit\ohm{\ensuremath{\symup{\Omega}}} % U+2126 absent de la police
\DeclareSIUnit\division{div} % réglages d'oscilloscope (ms/div, V/div)
\DeclareSIUnit\tr{tr} % tours (vitesse de rotation, ex. \SI{1500}{\tr\per\minute})
\usepackage{qrcode}
% \degree : commande fréquemment utilisée par les modèles pour un angle ou une
% température en dehors de \SI{}{} (non fournie par les paquets chargés).
\providecommand{\degree}{^{\circ}}
% \qed (fin de démonstration), utilisé par les modèles sans amsthm
\providecommand{\qed}{\hfill$\square$}
% \barre : double barre des valeurs interdites dans un tableau de variations (tkz-tab)
\providecommand{\barre}{\ensuremath{\|}}
% environnement proof (amsthm non chargé) utilisé par les modèles
\ifdefined\proof\else\newenvironment{proof}[1][Proof]{\par\noindent\textit{#1.}\ }{\qed\par}\fi
\usepackage{lastpage}
\usepackage{hyperref}
\hypersetup{hidelinks}

% ============================================================
% Palette « Pop » OnBuch+ — mêmes hex que la classe P (plus_ui.dart)
% ============================================================
\definecolor{poporange}{HTML}{F59321}
\definecolor{popdark}{HTML}{E07A0C}
\definecolor{popcream}{HTML}{FFF6E8}
\definecolor{popink}{HTML}{1C1714}
\definecolor{poppurple}{HTML}{7A5AE0}
\definecolor{popgreen}{HTML}{2E9E6A}
\definecolor{poppink}{HTML}{E0457B}
\definecolor{popblue}{HTML}{2F7DE1}
\definecolor{popgold}{HTML}{C98A12}
\definecolor{popmuted}{HTML}{8A7F76}
\definecolor{popline}{HTML}{EADFD2}
\definecolor{popgray}{HTML}{8A7F76}
\colorlet{popgrey}{popgray}
% Alias tolérants (le modèle nomme parfois les couleurs différemment)
\colorlet{popwhite}{white}
\colorlet{popblack}{popink}
\colorlet{popred}{poppink}
\colorlet{popyellow}{popgold}
% Teintes claires pour fonds de tableaux / zones
\colorlet{popblueL}{popblue!10!white}
\colorlet{poporangeL}{poporange!14!white}
\colorlet{popgreenL}{popgreen!12!white}
\colorlet{poppurpleL}{poppurple!12!white}
\colorlet{poppinkL}{poppink!12!white}
\colorlet{popgoldL}{popgold!14!white}

\pagecolor{popcream}

% ============================================================
% Typographie
% ============================================================
\defaultfontfeatures{Ligatures=TeX}
\setmainfont{PlusJakartaSans-Medium.ttf}[Path=../../../fonts/,BoldFont=PlusJakartaSans-Bold.ttf]
% Maths en sans-serif (Fira Math) avec chiffres et lettres droites de Plus
% Jakarta, pour que les formules se fondent dans le texte.
\usepackage[math-style=ISO,bold-style=ISO]{unicode-math}
\setmathfont{FiraMath-Regular.otf}[Path=../../../fonts/]
\setmathfont{PlusJakartaSans-Medium.ttf}[Path=../../../fonts/,range={up/num,up/{Latin,latin},\mathup}]
% (icomma retiré : décimales avec point en anglais)
% Caractères mathématiques Unicode absents des polices de texte (Plus Jakarta,
% Archivo Black) : rendus par leur équivalent mathématique, en texte comme en
% formule — sinon ils s'affichent en carré vide (ex. « Arithmétique dans ℤ »).
\usepackage{newunicodechar}
\newunicodechar{ℝ}{\ensuremath{\mathbb{R}}}
\newunicodechar{ℤ}{\ensuremath{\mathbb{Z}}}
\newunicodechar{ℂ}{\ensuremath{\mathbb{C}}}
\newunicodechar{ℕ}{\ensuremath{\mathbb{N}}}
\newunicodechar{ℚ}{\ensuremath{\mathbb{Q}}}
\newunicodechar{𝔻}{\ensuremath{\mathbb{D}}}
\newunicodechar{α}{\ensuremath{\alpha}}
\newunicodechar{β}{\ensuremath{\beta}}
\newunicodechar{γ}{\ensuremath{\gamma}}
\newunicodechar{δ}{\ensuremath{\delta}}
\newunicodechar{ε}{\ensuremath{\varepsilon}}
\newunicodechar{θ}{\ensuremath{\theta}}
\newunicodechar{λ}{\ensuremath{\lambda}}
\newunicodechar{μ}{\ensuremath{\mu}}
\newunicodechar{σ}{\ensuremath{\sigma}}
\newunicodechar{τ}{\ensuremath{\tau}}
\newunicodechar{φ}{\ensuremath{\varphi}}
\newunicodechar{ψ}{\ensuremath{\psi}}
\newunicodechar{ω}{\ensuremath{\omega}}
\newunicodechar{Σ}{\ensuremath{\Sigma}}
% majuscules produites par \MakeUppercase dans les titres (α-aminés -> Α)
\newunicodechar{Α}{\ensuremath{\alpha}}
\newunicodechar{Β}{\ensuremath{\beta}}
\newunicodechar{Γ}{\ensuremath{\Gamma}}
\newunicodechar{Δ}{\ensuremath{\Delta}}
\newunicodechar{Ε}{\ensuremath{\varepsilon}}
\newunicodechar{Θ}{\ensuremath{\Theta}}
\newunicodechar{Λ}{\ensuremath{\Lambda}}
\newunicodechar{Μ}{\ensuremath{\mu}}
\newunicodechar{Τ}{\ensuremath{\tau}}
\newunicodechar{Φ}{\ensuremath{\Phi}}
\newunicodechar{Ψ}{\ensuremath{\Psi}}
\newunicodechar{Ω}{\ensuremath{\Omega}}
\newunicodechar{↦}{\ensuremath{\mapsto}}
\newunicodechar{∈}{\ensuremath{\in}}
\newunicodechar{∉}{\ensuremath{\notin}}
\newunicodechar{∧}{\ensuremath{\wedge}}
\newunicodechar{⇔}{\ensuremath{\Leftrightarrow}}
\newunicodechar{⟺}{\ensuremath{\Longleftrightarrow}}
\newunicodechar{⇒}{\ensuremath{\Rightarrow}}
\newunicodechar{⟹}{\ensuremath{\Longrightarrow}}
\newunicodechar{∘}{\ensuremath{\circ}}
\newunicodechar{∪}{\ensuremath{\cup}}
\newunicodechar{∩}{\ensuremath{\cap}}
\newunicodechar{≡}{\ensuremath{\equiv}}
\newunicodechar{⊂}{\ensuremath{\subset}}
\newunicodechar{⊥}{\ensuremath{\perp}}
\newunicodechar{∃}{\ensuremath{\exists}}
\newunicodechar{∀}{\ensuremath{\forall}}
\newunicodechar{∛}{\ensuremath{\sqrt[3]{\,}}}
\newunicodechar{∜}{\ensuremath{\sqrt[4]{\,}}}
\newunicodechar{⃗}{\ensuremath{{}^{\to}}}
\newunicodechar{ⁿ}{\ifmmode^{n}\else\textsuperscript{\ensuremath{n}}\fi}
\newunicodechar{ˣ}{\ifmmode^{x}\else\textsuperscript{\ensuremath{x}}\fi}
\newunicodechar{ᵇ}{\ifmmode^{b}\else\textsuperscript{\ensuremath{b}}\fi}
\newunicodechar{ʳ}{\ifmmode^{r}\else\textsuperscript{\ensuremath{r}}\fi}
\newunicodechar{ᵘ}{\ifmmode^{u}\else\textsuperscript{\ensuremath{u}}\fi}
\newunicodechar{ʸ}{\ifmmode^{y}\else\textsuperscript{\ensuremath{y}}\fi}
\newunicodechar{ᵃ}{\ifmmode^{a}\else\textsuperscript{\ensuremath{a}}\fi}
\newunicodechar{ᵉ}{\ifmmode^{e}\else\textsuperscript{\ensuremath{e}}\fi}
\newunicodechar{ᵏ}{\ifmmode^{k}\else\textsuperscript{\ensuremath{k}}\fi}
\newunicodechar{ᵖ}{\ifmmode^{p}\else\textsuperscript{\ensuremath{p}}\fi}
\newunicodechar{ᴺ}{\ifmmode^{N}\else\textsuperscript{\ensuremath{N}}\fi}
\newunicodechar{ᶜ}{\ifmmode^{c}\else\textsuperscript{\ensuremath{c}}\fi}
\newunicodechar{ʰ}{\ifmmode^{h}\else\textsuperscript{\ensuremath{h}}\fi}
\newunicodechar{ᵠ}{\ifmmode^{q}\else\textsuperscript{\ensuremath{q}}\fi}
\newunicodechar{ₙ}{\ifmmode_{n}\else\textsubscript{\ensuremath{n}}\fi}
\newunicodechar{ₐ}{\ifmmode_{a}\else\textsubscript{\ensuremath{a}}\fi}
\newunicodechar{ᵢ}{\ifmmode_{i}\else\textsubscript{\ensuremath{i}}\fi}
\newunicodechar{ᵣ}{\ifmmode_{r}\else\textsubscript{\ensuremath{r}}\fi}
\newunicodechar{ₖ}{\ifmmode_{k}\else\textsubscript{\ensuremath{k}}\fi}
\newunicodechar{ᵦ}{\ifmmode_{\beta}\else\textsubscript{\ensuremath{\beta}}\fi}
\newunicodechar{ₓ}{\ifmmode_{x}\else\textsubscript{\ensuremath{x}}\fi}
\newfontfamily\popdisplay{ArchivoBlack-Regular.ttf}[Path=../../../fonts/]
\newfontfamily\poplogo{SpaceGrotesk-Bold.ttf}[Path=../../../fonts/]
\newfontfamily\popsemi{PlusJakartaSans-SemiBold.ttf}[Path=../../../fonts/]
\newfontfamily\popxbold{PlusJakartaSans-ExtraBold.ttf}[Path=../../../fonts/]
\newfontfamily\popxboldls{PlusJakartaSans-ExtraBold.ttf}[Path=../../../fonts/,LetterSpace=3.0]

\setlist[enumerate]{leftmargin=1.7em,itemsep=5pt,topsep=4pt}
\setlist[itemize]{leftmargin=1.7em,itemsep=4pt,topsep=4pt,label={\color{poporange}\textbullet}}
\setlength{\parindent}{0pt}
\setlength{\parskip}{5pt}
\renewcommand{\arraystretch}{1.25}

% pgfplots : style Pop par défaut
\pgfplotsset{
  every axis/.append style={axis line style={popink,line width=1pt},tick style={popink},
    grid style={popline,line width=0.5pt},label style={font=\small},tick label style={font=\footnotesize}},
  popcurve/.style={poporange,line width=1.8pt,no markers,smooth},
}
% Même style disponible dans un \draw TikZ simple (hors pgfplots)
\tikzset{popcurve/.style={poporange,line width=1.8pt}}
\tikzset{popgrid/.style={popline,very thin}}
\colorlet{popcurve}{poporange} % le nom du style est parfois utilisé comme couleur

% ============================================================
% Pill / squiggle
% ============================================================
\newcommand{\poppill}[3]{%
  \tikz[baseline=-0.55ex]\node[draw=popink,line width=1.2pt,rounded corners=2.6pt,%
    fill=#2,inner xsep=6.5pt,inner ysep=3pt]{%
    \popxboldls\fontsize{7}{7}\selectfont\color{#3}\MakeUppercase{#1}};%
}
\newcommand{\popsquiggle}[1]{%
  \tikz[baseline=-0.4ex]{
    \draw[#1,line width=2.6pt,line cap=round]
      plot [smooth] coordinates {(0,0.09) (0.34,0.01) (0.68,0.21) (1.02,0.01) (1.36,0.10)};
  }%
}
% Mot-clé surligné (marqueur orange sous le texte)
\newcommand{\cle}[1]{{\popxbold\color{popdark}#1}}

% ============================================================
% En-tête / pied
% ============================================================
\pagestyle{fancy}
\fancyhf{}
\renewcommand{\headrulewidth}{0pt}
\renewcommand{\footrulewidth}{0pt}
\lhead{\raisebox{-0.15ex}{\poplogo\fontsize{13}{13}\selectfont\color{popink}OnBuch\color{poporange}+}}
\rhead{\poppill{\DOCMATIERE\ \textbullet\ \DOCNIVEAU\ \textbullet\ Chapter \DOCLECON}{white}{popink}}
\lfoot{\popxbold\fontsize{7.6}{7.6}\selectfont\color{popmuted}\MakeUppercase{Signed course OnBuch+}\ \textbullet\ {\popsemi\DOCTITRE}}
\rfoot{\tikz[baseline=-0.6ex]\node[rounded corners=8.5pt,fill=popink,minimum height=17pt,minimum width=17pt,inner xsep=5pt,inner sep=0]{\popxbold\fontsize{8}{8}\selectfont\color{popcream}\thepage\,/\,\pageref*{LastPage}};}

% ============================================================
% Titres
% ============================================================
\newcommand{\chaptitre}[2]{%
  \begin{tcolorbox}[enhanced,colback=poporange,colframe=popink,boxrule=2.6pt,arc=11pt,
      drop shadow=popink,left=15pt,right=15pt,top=12pt,bottom=11pt,before skip=3pt,after skip=18pt]
    \poppill{#2}{popink}{popcream}\par\vspace{9pt}
    {\raggedright\hyphenpenalty=10000\exhyphenpenalty=10000\popdisplay\fontsize{22}{24}\selectfont\color{popcream}\MakeUppercase{#1}\par}
  \end{tcolorbox}
}
% Section numérotée : pastille orange avec le numéro + titre Archivo
\newcounter{popsec}
\newcommand{\coursec}[1]{%
  \stepcounter{popsec}\setcounter{popsub}{0}%
  \par\vspace{16pt}\noindent
  \begin{minipage}{\linewidth}
  \tikz[baseline=-0.7ex]\node[draw=popink,line width=1.6pt,fill=poporange,rounded corners=4pt,minimum size=22pt,inner sep=0]{\popdisplay\fontsize{12}{12}\selectfont\color{popcream}\thepopsec};%
  \hspace{8pt}{\popdisplay\fontsize{15}{17}\selectfont\color{popink}\MakeUppercase{#1}}\par
  \vspace{4pt}\noindent{\color{poporange}\rule{36pt}{3.6pt}}
  \end{minipage}\par\nopagebreak\vspace{8pt}
}
\newcounter{popsub}
\newcommand{\courssub}[1]{%
  \stepcounter{popsub}%
  \par\vspace{9pt}\noindent{\popxbold\fontsize{11.5}{13}\selectfont\color{popink}{\color{poporange}\thepopsec.\thepopsub}\hspace{6pt}#1}\par\nopagebreak\vspace{4pt}
}

% ============================================================
% Boîtes pédagogiques — même recette : fond blanc, bordure épaisse colorée,
% ombre dure encre, badge plein. Titre optionnel [#1].
% ============================================================
\tcbset{popbase/.style={breakable,enhanced,colback=white,boxrule=2.3pt,arc=9pt,
  shadow={3pt}{-3pt}{0pt}{popink},left=14pt,right=14pt,top=11pt,bottom=11pt,before skip=11pt,after skip=15pt,
  fontupper=\popsemi\fontsize{10.6}{14.5}\selectfont\color{popink},
  fontlower=\popsemi\fontsize{10.6}{14.5}\selectfont\color{popink}}}
% Coche dessinée (le glyphe ✓ n'existe pas dans les polices du document)
\newcommand{\popcheck}[1]{\tikz[baseline=-0.5ex]\draw[#1,line width=1.6pt,line cap=round,line join=round](-0.13,0.0)--(-0.04,-0.09)--(0.13,0.1);}
\providecommand{\coche}{\popcheck{popgreen}}
\providecommand{\resultat}[1]{\par\smallskip\noindent\fcolorbox{popgreen}{popgreenL}{\parbox{\dimexpr\linewidth-2\fboxsep-2\fboxrule\relax}{\textbf{Result.} #1}}\par\smallskip}
\newcommand{\popboxhead}[3]{\poppill{#1}{#2}{white}\ifx\relax#3\relax\else\hspace{6pt}{\popxbold\fontsize{10.6}{12}\selectfont\color{popink}#3}\fi\par\vspace{7pt}}

\newtcolorbox{aretenir}[1][]{popbase,colframe=popgreen,before upper={\popboxhead{Key points}{popgreen}{#1}}}
\newtcolorbox{exemplebox}[1][]{popbase,colframe=poporange,before upper={\popboxhead{Exemple}{poporange}{#1}}}
\newtcolorbox{definition}[1][]{popbase,colframe=popblue,before upper={\popboxhead{Definition}{popblue}{#1}}}
\newtcolorbox{propriete}[1][]{popbase,colframe=poppurple,before upper={\popboxhead{Property}{poppurple}{#1}}}
\newtcolorbox{methode}[1][]{popbase,colframe=popgold,before upper={\popboxhead{Method}{popgold}{#1}}}
\newtcolorbox{attention}[1][]{popbase,colframe=poppink,before upper={\popboxhead{Warning}{poppink}{#1}}}
\newtcolorbox{experience}[1][]{popbase,colframe=popblue,colback=popblueL,before upper={\popboxhead{Experiment}{popblue}{#1}}}
% « Le savais-tu ? » : carte violette pleine (contexte, culture, Cameroun)
\newtcolorbox{savaistu}[1][]{popbase,colback=poppurple,colframe=popink,
  fontupper=\popsemi\fontsize{10.4}{14.2}\selectfont\color{white},
  before upper={\poppill{Did you know?}{white}{poppurple}\ifx\relax#1\relax\else\hspace{6pt}{\popxbold\color{white}#1}\fi\par\vspace{7pt}}}
% Exercice résolu : énoncé en haut, \tcblower puis la solution sur fond crème
\newcounter{popexo}\newcounter{popexr}
\newtcolorbox{exoresolu}[1][]{popbase,colframe=poporange,colbacklower=popcream,
  segmentation style={popink,line width=1.2pt,dashed},
  before upper={\stepcounter{popexr}\popboxhead{Worked example \thepopexr}{poporange}{#1}},
  before lower={\poppill{Solution}{popink}{popcream}\par\vspace{6pt}}}
% Exercice (non résolu en place) — la correction est dans la partie Corrigés
\newtcolorbox{exercice}[1][]{popbase,colframe=popink,
  before upper={\stepcounter{popexo}\popboxhead{Exercise \thepopexo}{popink}{#1}}}
% Corrigé
\newtcolorbox{corrige}[1][]{popbase,colframe=popgreen,colback=popgreenL,
  before upper={\popboxhead{Solution}{popgreen}{#1}}}
% Objectifs / prérequis / bilan
\newtcolorbox{objectifs}{popbase,colframe=popink,colback=poporangeL,
  before upper={\popboxhead{Chapter objectives}{popink}{}}}
\newtcolorbox{prerequis}{popbase,colframe=popink,colback=popblueL,
  before upper={\popboxhead{Prerequisites}{popblue}{}}}
\newtcolorbox{bilan}{popbase,colframe=popink,colback=white,boxrule=2.8pt,
  before upper={\popboxhead{Fiche bilan}{poporange}{L'essentiel à connaître pour l'examen}}}
% Figure encadrée avec légende
\newtcolorbox{popfigure}[1][]{enhanced,breakable=false,colback=white,colframe=popline,boxrule=1.4pt,arc=8pt,
  left=10pt,right=10pt,top=10pt,bottom=8pt,before skip=10pt,after skip=12pt,halign=center,
  lower separated=false,halign lower=center,
  fontlower=\popsemi\fontsize{9}{11.5}\selectfont\color{popmuted},
  #1}
\newcommand{\legende}[1]{\par\vspace{6pt}{\popsemi\fontsize{9}{11.5}\selectfont\color{popmuted}#1}}

% ============================================================
% Signature OnBuch+
% ============================================================
\newcommand{\poplogoinline}[1]{{\poplogo\fontsize{#1}{#1}\selectfont\color{popink}OnBuch\color{poporange}+}}
\newcommand{\signatureonbuch}{%
  \begin{tcolorbox}[enhanced,colback=popink,colframe=popink,boxrule=2.2pt,arc=10pt,
      drop shadow=poporange,left=14pt,right=14pt,top=10pt,bottom=10pt,before skip=0pt,after skip=0pt]
    \tikz[baseline=-0.7ex]\node[circle,fill=popgreen,draw=popcream,line width=1.3pt,minimum size=22pt,inner sep=0]{\popcheck{white}};%
    \hspace{9pt}\begin{minipage}[c]{0.62\linewidth}
      {\popxbold\fontsize{12}{13}\selectfont\color{popcream}Signed\ \textbullet\ {\poplogo OnBuch\color{poporange}+}}\par\vspace{2pt}
      {\raggedright\popsemi\fontsize{8}{10.5}\selectfont\color{popline}\MakeUppercase{OnBuch+ teaching team}\\\MakeUppercase{Follows the official MINESEC syllabus}\par}
    \end{minipage}\hfill
    {\poplogo\fontsize{22}{22}\selectfont\color{popcream}OnBuch\color{poporange}+}
  \end{tcolorbox}%
}

% ============================================================
% Page de garde
% #1 titre · #2 matière · #3 niveau · #4 module · #5 n° leçon · #6 durée ·
% #7 description / objectifs (texte ou itemize)
% ============================================================
\newcommand{\pagegarde}[7]{%
  \thispagestyle{empty}
  \newgeometry{a4paper,margin=1.7cm,top=1.6cm,bottom=1.4cm}%
  \begin{tikzpicture}[remember picture,overlay]
    \fill[poporange!18] ([xshift=-3.2cm,yshift=-3.6cm]current page.north east) circle (4.6cm);
    \fill[poppurple!14] ([xshift=2.4cm,yshift=7.6cm]current page.south west) circle (3.8cm);
  \end{tikzpicture}%
  {\poplogo\fontsize{30}{30}\selectfont\color{popink}OnBuch\color{poporange}+}\hfill
  \raisebox{0.6ex}{\poppill{Signed course \textbullet\ Premium}{popink}{popcream}}\par
  \vspace{1pt}\popsquiggle{poppurple}\par
  \vspace{8pt}
  \begin{tcolorbox}[enhanced,colback=poporange,colframe=popink,boxrule=2.8pt,arc=13pt,
      drop shadow=popink,left=18pt,right=18pt,top=12pt,bottom=13pt,after skip=10pt]
    \poppill{#2 \textbullet\ #3}{popink}{popcream}\hspace{5pt}\poppill{#4}{white}{popink}\par\vspace{8pt}
    {\popdisplay\fontsize{14}{14}\selectfont\color{popink}CHAPTER #5}\par\vspace{4pt}
    {\raggedright\hyphenpenalty=10000\exhyphenpenalty=10000\popdisplay\fontsize{25}{27}\selectfont\color{popcream}\MakeUppercase{#1}\par}
    \vspace{8pt}
    {\popxbold\fontsize{9.5}{9.5}\selectfont\color{popink}\MakeUppercase{Suggested time: #6}}
  \end{tcolorbox}
  \begin{tcolorbox}[enhanced,colback=white,colframe=popink,boxrule=2.2pt,arc=10pt,
      drop shadow=popink,left=16pt,right=16pt,top=10pt,bottom=10pt,after skip=8pt]
    \poppill{In this chapter}{poppurple}{white}\par\vspace{5pt}
    \popsemi\fontsize{9.3}{12.6}\selectfont\color{popink}#7
  \end{tcolorbox}
  \noindent\begin{minipage}[t]{0.487\linewidth}
    \begin{tcolorbox}[enhanced,colback=popgreen,colframe=popink,boxrule=2.2pt,arc=10pt,
        drop shadow=popink,left=11pt,right=11pt,top=10pt,bottom=10pt,height=3.1cm,valign=center]
      \tcbox[colback=white,colframe=popink,boxrule=1.3pt,arc=5pt,boxsep=0pt,nobeforeafter,
        left=3pt,right=3pt,top=3pt,bottom=3pt,valign=center]{\qrcode[height=1.9cm]{https://play.google.com/store/apps/details?id=com.luvvix.onbuch&hl=en}}%
      \hspace{8pt}\begin{minipage}[c]{0.5\linewidth}
        {\popdisplay\fontsize{11}{12.5}\selectfont\color{white}\MakeUppercase{Download OnBuch}}\par\vspace{4pt}
        {\raggedright\popsemi\fontsize{8.4}{11}\selectfont\color{white}Free \textbullet\ Google Play\\AI tutor, past papers, results\par}
      \end{minipage}
    \end{tcolorbox}
  \end{minipage}\hfill
  \begin{minipage}[t]{0.487\linewidth}
    \begin{tcolorbox}[enhanced,colback=poppurple,colframe=popink,boxrule=2.2pt,arc=10pt,
        drop shadow=popink,left=11pt,right=11pt,top=10pt,bottom=10pt,height=3.1cm,valign=center]
      \tikz[baseline=-0.7ex]\node[circle,fill=white,draw=popink,line width=1.3pt,minimum size=1.6cm,inner sep=0]{\popdisplay\fontsize{24}{24}\selectfont\color{poppurple}+};%
      \hspace{8pt}\begin{minipage}[c]{0.6\linewidth}
        {\popdisplay\fontsize{11}{12.5}\selectfont\color{white}\MakeUppercase{Go OnBuch+}}\par\vspace{4pt}
        {\raggedright\popsemi\fontsize{8.4}{11}\selectfont\color{white}All signed courses, unlimited AI tutor, PDF sheets — OnBuch+ tab in the app\par}
      \end{minipage}
    \end{tcolorbox}
  \end{minipage}
  \vspace{8pt}
  \popcontenu
  \vfill
  \signatureonbuch
  \restoregeometry
  \newpage
}

% Rangée « Ce cours contient » (page de garde)
\newcommand{\poptile}[3]{% #1 couleur #2 gros texte #3 légende
  \begin{tcolorbox}[enhanced,colback=#1,colframe=popink,boxrule=2pt,arc=9pt,shadow={2.5pt}{-2.5pt}{0pt}{popink},
      left=6pt,right=6pt,top=9pt,bottom=9pt,halign=center,valign=center,height=1.75cm,width=\linewidth,nobeforeafter]
    {\popdisplay\fontsize{12.5}{13}\selectfont\color{white}\MakeUppercase{#2}}\par\vspace{3pt}
    {\centering\popsemi\fontsize{7.8}{9.5}\selectfont\color{white}#3\par}
  \end{tcolorbox}}
\newcommand{\popcontenu}{%
  \par\noindent{\popxboldls\fontsize{7.5}{7.5}\selectfont\color{popmuted}THIS SIGNED COURSE CONTAINS}\par\vspace{7pt}
  \noindent\begin{minipage}[t]{0.235\linewidth}\poptile{popblue}{Course}{complete, illustrated, step by step}\end{minipage}\hfill
  \begin{minipage}[t]{0.235\linewidth}\poptile{poporange}{Methods}{and worked examples}\end{minipage}\hfill
  \begin{minipage}[t]{0.235\linewidth}\poptile{poppink}{Exercises}{exam style + solutions}\end{minipage}\hfill
  \begin{minipage}[t]{0.235\linewidth}\poptile{popgreen}{Summary sheet}{the essentials to remember}\end{minipage}\par}

% Sommaire
\newtcolorbox{sommaire}{enhanced,colback=white,colframe=popink,boxrule=2.2pt,arc=10pt,
  drop shadow=popink,left=18pt,right=18pt,top=13pt,bottom=13pt,before skip=6pt,after skip=20pt,
  before upper={\poppill{Contents}{poppurple}{white}\par\vspace{9pt}\popsemi\fontsize{10.6}{14.5}\selectfont\color{popink}}}

% Signature de fin de cours — poussée tout en bas de la dernière page
\newcommand{\findecours}{%
  \par\vfill\nopagebreak
  \begin{center}\popsquiggle{poporange}\end{center}\vspace{4pt}
  \signatureonbuch
}
\AtBeginDocument{\def\llbracket{\mathopen{[\mkern-3.5mu[}}\def\rrbracket{\mathclose{]\mkern-3.5mu]}}} % intervalles d'entiers (glyphe ⟦ absent de la police)
% Glyphes absents de Fira Math : redéfinis à partir de glyphes disponibles
\AtBeginDocument{%
  \def\setminus{\mathbin{\backslash}}\let\smallsetminus\setminus
  \def\vdots{\mathord{\vbox{\baselineskip3.2pt\lineskiplimit0pt\kern4pt\hbox{.}\hbox{.}\hbox{.}}}}%
  \def\ddots{\mathinner{\mkern1mu\raise6.5pt\hbox{.}\mkern2mu\raise3.6pt\hbox{.}\mkern2mu\raise0.7pt\hbox{.}\mkern1mu}}%
  \def\triangle{\mathord{\scalebox{0.9}{$\symup{\Delta}$}}}%
  \def\nabla{\mathord{\raisebox{\depth}{\scalebox{1}[-1]{$\symup{\Delta}$}}}}%
  \def\checkmark{\popcheck{popdark}}%
  \def\star{\mathbin{\ast}}\def\bigstar{\mathord{\ast}}%
  \def\diamond{\mathbin{\scalebox{0.6}{$\Diamond$}}}%
  \def\bigcup{\mathop{\mathchoice{\raisebox{-0.3em}{\scalebox{1.7}{$\cup$}}}{\scalebox{1.25}{$\cup$}}{\cup}{\cup}}\displaylimits}%
  \def\bigcap{\mathop{\mathchoice{\raisebox{-0.3em}{\scalebox{1.7}{$\cap$}}}{\scalebox{1.25}{$\cap$}}{\cap}{\cap}}\displaylimits}%
  \def\lesseqqgtr{\mathrel{\lessgtr}}\def\bigoplus{\mathop{\oplus}\displaylimits}%
}
\providecommand{\ssi}{if and only if} % abréviation fréquente
\AtBeginDocument{\providecommand{\wideparen}{\overparen}} % arc de cercle

% Fonctions usuelles que les modèles utilisent sans que amsmath les fournisse
\providecommand{\arsinh}{\operatorname{arsinh}}\providecommand{\arcosh}{\operatorname{arcosh}}
\providecommand{\artanh}{\operatorname{artanh}}\providecommand{\arsech}{\operatorname{arsech}}
\providecommand{\arcsch}{\operatorname{arcsch}}\providecommand{\arcoth}{\operatorname{arcoth}}
\providecommand{\arcsinh}{\operatorname{arsinh}}\providecommand{\arccosh}{\operatorname{arcosh}}
\providecommand{\arctanh}{\operatorname{artanh}}\providecommand{\sech}{\operatorname{sech}}
\providecommand{\csch}{\operatorname{csch}}\providecommand{\arcsec}{\operatorname{arcsec}}
\providecommand{\arccsc}{\operatorname{arccsc}}\providecommand{\arccot}{\operatorname{arccot}}
\providecommand{\sgn}{\operatorname{sgn}}\providecommand{\lcm}{\operatorname{lcm}}
\providecommand{\Var}{\operatorname{Var}}\providecommand{\Cov}{\operatorname{Cov}}

%% FIGFIX-BEGIN v5 — correctifs globaux des figures (docs/onbuch-generation/tools/figfix)
\usepackage{adjustbox}\usepackage{etoolbox}\tcbuselibrary{raster}
% toute figure TikZ/pgfplots plus large que la ligne est réduite à la largeur disponible
\BeforeBeginEnvironment{tikzpicture}{\begin{adjustbox}{max width=\linewidth}}
\AfterEndEnvironment{tikzpicture}{\end{adjustbox}}
% légendes pgfplots lisibles par-dessus les courbes
\pgfplotsset{every axis/.append style={legend style={fill=white,fill opacity=0.92,text opacity=1,draw=black!15,inner sep=3pt}}}
% étiquettes d'arêtes (syntaxe edge["..."]) avec fond blanc
\tikzset{every edge quotes/.append style={fill=white,inner sep=1.2pt}}
%% FIGFIX-END
\begin{document}

\pagegarde{Linear Transformation}{Further Mathematics}{Upper Sixth}{Upper Sixth — Further mathematics}{13}{6 periods}{This chapter studies linear transformations of the plane and of space: their algebraic properties, their kernels and images, their inverses and invariant points, the way they act on lines and planes, and the eigenvalues and eigenvectors that reveal their hidden structure. It provides the essential bridge between matrix algebra and geometric reasoning required for the GCE Advanced Level examination.
\vspace{4pt}
\begin{multicols}{2}\small
\begin{itemize}
\item By the end of this chapter I can define a linear transformation and verify whether a given mapping is linear
\item By the end of this chapter I can state and use the properties of linear transformations (preservation of addition, scalar multiplication, the zero vector and linear combinations)
\item By the end of this chapter I can determine the kernel and the image of a linear transformation and interpret them geometrically
\item By the end of this chapter I can decide whether a linear transformation is injective, surjective or bijective using the kernel and image
\item By the end of this chapter I can find the inverse of an invertible linear transformation and determine its invariant points and invariant lines
\item By the end of this chapter I can find the images of lines and planes under a linear transformation
\end{itemize}\end{multicols}}

\begin{sommaire}
\begin{multicols}{2}
\begin{enumerate}
\item Linear Transformations and Their Properties
\item Kernel and Image of a Linear Transformation
\item Inverse of a Linear Transformation and Invariant Points
\item Images of Lines and Planes
\item Eigenvalues and Eigenvectors
\item Integration activity
\item Methods and worked examples
\item Exercises
\item Solutions to the exercises
\item Summary sheet
\end{enumerate}
\end{multicols}
\end{sommaire}

\begin{objectifs}
\begin{itemize}
\item By the end of this chapter I can define a linear transformation and verify whether a given mapping is linear
\item By the end of this chapter I can state and use the properties of linear transformations (preservation of addition, scalar multiplication, the zero vector and linear combinations)
\item By the end of this chapter I can determine the kernel and the image of a linear transformation and interpret them geometrically
\item By the end of this chapter I can decide whether a linear transformation is injective, surjective or bijective using the kernel and image
\item By the end of this chapter I can find the inverse of an invertible linear transformation and determine its invariant points and invariant lines
\item By the end of this chapter I can find the images of lines and planes under a linear transformation
\item By the end of this chapter I can compute the eigenvalues and eigenvectors of a linear transformation or of its associated matrix
\item By the end of this chapter I can use eigenvectors to describe the geometric effect of a transformation (stretching along fixed directions)
\item By the end of this chapter I can solve GCE-style problems combining matrices, linear transformations and geometry
\end{itemize}
\end{objectifs}

\begin{prerequis}
\begin{itemize}
\item Vectors in the plane and in space: coordinates, addition, scalar multiplication, collinearity (Lower Sixth)
\item Matrices: operations, determinant of a 2x2 and 3x3 matrix, inverse of a matrix (Lower Sixth / Term 1)
\item Solving systems of linear equations by substitution, elimination and row reduction
\item Equations of lines and planes in 2D and 3D (vector, parametric and Cartesian forms)
\item Mappings and functions: domain, image set, injectivity, surjectivity, composition (Lower Sixth)
\end{itemize}
\end{prerequis}

\newpage

\coursec{Linear Transformations and Their Properties}

When a graphic designer in Douala enlarges a photo of the Reunification Monument, or when a surveyor in Bamenda projects a hillside onto a flat map, every point of the original figure is moved to a new position by a precise rule. Some rules preserve straight lines, parallelism and the origin; others do not. How can we describe such rules algebraically, predict which points remain fixed, and find the special directions that are merely stretched or compressed? These questions lead us to the powerful notion of a \cle{linear transformation}, a tool that unifies geometry, matrices and systems of equations.

\courssub{Definition and Fundamental Axioms}

Let $E$ and $F$ be two vector spaces over the same field $\mathbb{R}$. A mapping $T : E \to F$ is called a \cle{linear transformation} (or \cle{linear map}) if it satisfies two axioms.

\begin{definition}[Linear Transformation]
A map $T : E \to F$ is linear if and only if, for all vectors $\vec{u}, \vec{v} \in E$ and every scalar $k \in \mathbb{R}$:
\begin{enumerate}
    \item \textbf{Additivity:} $T(\vec{u} + \vec{v}) = T(\vec{u}) + T(\vec{v})$;
    \item \textbf{Homogeneity:} $T(k\vec{u}) = k\,T(\vec{u})$.
\end{enumerate}
\end{definition}

These two conditions can be combined into a single, extremely useful criterion.

\begin{propriete}[Combined Linearity Condition]
$T$ is linear $\iff$ $T(a\vec{u} + b\vec{v}) = a\,T(\vec{u}) + b\,T(\vec{v})$ for all $\vec{u},\vec{v} \in E$ and all scalars $a,b$.
\end{propriete}

\textbf{Proof.}\par
If $T$ is linear, then $T(a\vec{u}+b\vec{v}) = T(a\vec{u}) + T(b\vec{v}) = aT(\vec{u}) + bT(\vec{v})$. Conversely, taking $a=b=1$ gives additivity, and taking $b=0$ gives homogeneity. \hfill$\square$

This combined condition is the most efficient tool for \emph{testing} whether a given mapping is linear. It also shows immediately that a linear transformation preserves \emph{all} linear combinations:
\[
T\Bigl(\sum_{i=1}^n \alpha_i \vec{v}_i\Bigr) = \sum_{i=1}^n \alpha_i T(\vec{v}_i).
\]

\courssub{Immediate Consequences}

\begin{propriete}[Fundamental Properties of Linear Maps]
Let $T : E \to F$ be linear. Then:
\begin{enumerate}
    \item $T(\vec{0}_E) = \vec{0}_F$ (the zero vector maps to the zero vector);
    \item $T(-\vec{u}) = -T(\vec{u})$ for all $\vec{u} \in E$;
    \item $T(\vec{u} - \vec{v}) = T(\vec{u}) - T(\vec{v})$;
    \item $T$ preserves linear combinations: $T\bigl(\sum_i \alpha_i \vec{v}_i\bigr) = \sum_i \alpha_i T(\vec{v}_i)$.
\end{enumerate}
\end{propriete}

\textbf{Proof.}\par
(1) $T(\vec{0}) = T(0\cdot\vec{0}) = 0\cdot T(\vec{0}) = \vec{0}$.
(2) $T(-\vec{u}) = T((-1)\vec{u}) = (-1)T(\vec{u}) = -T(\vec{u})$.
(3) $T(\vec{u}-\vec{v}) = T(\vec{u}+(-\vec{v})) = T(\vec{u})+T(-\vec{v}) = T(\vec{u})-T(\vec{v})$.
(4) follows by induction from additivity and homogeneity. \hfill$\square$

\begin{attention}[The Zero-Vector Test — A Quick Way to Spot Non-Linearity]
If a mapping $T$ does \textbf{not} send the zero vector to the zero vector, it is \textbf{not} linear. This is the very first check you should perform in an examination.
\begin{itemize}
    \item $T(x,y) = (x+1, y)$ gives $T(0,0) = (1,0) \neq (0,0)$, so $T$ is NOT linear.
    \item $T(x,y) = (x^2, y)$ gives $T(0,0) = (0,0)$: it passes the zero test, but fails homogeneity (see below). The zero test can only prove \emph{non}-linearity, never linearity.
\end{itemize}
\end{attention}

\courssub{Standard Examples in the Plane}

The most familiar linear transformations act on the Euclidean plane $\mathbb{R}^2$. We denote the standard basis vectors by $\vec{i} = (1,0)$ and $\vec{j} = (0,1)$.

\begin{exemplebox}[Identity and Zero Transformations]
\begin{itemize}
    \item \textbf{Identity:} $I(\vec{u}) = \vec{u}$. Matrix: $\begin{pmatrix} 1 & 0 \\ 0 & 1 \end{pmatrix}$.
    \item \textbf{Zero map:} $O(\vec{u}) = \vec{0}$. Matrix: $\begin{pmatrix} 0 & 0 \\ 0 & 0 \end{pmatrix}$.
\end{itemize}
Both are trivially linear.
\end{exemplebox}

\begin{exemplebox}[Homothety (Enlargement) of Ratio $k$]
$H_k(x,y) = (kx, ky)$ stretches every vector by the factor $k$ from the origin.
Matrix: $\begin{pmatrix} k & 0 \\ 0 & k \end{pmatrix} = kI$.
\end{exemplebox}

\begin{exemplebox}[Rotation about the Origin through Angle $\theta$]
$R_\theta(x,y) = (x\cos\theta - y\sin\theta,\; x\sin\theta + y\cos\theta)$.
Matrix: $\begin{pmatrix} \cos\theta & -\sin\theta \\ \sin\theta & \cos\theta \end{pmatrix}$.
\end{exemplebox}

\begin{popfigure}
\begin{tikzpicture}[scale=1.3, >=stealth]
    \draw[->] (-2.2,0) -- (2.5,0) node[right] {$x$};
    \draw[->] (0,-1.6) -- (0,2.5) node[above] {$y$};
    \draw[popblue, very thick, ->] (0,0) -- (1.5,1) node[right, popblue] {$\vec{u}=(x,y)$};
    \draw[poporange, very thick, ->] (0,0) -- (-0.16,1.81) node[above left, poporange] {$R_\theta(\vec{u})$};
    \draw[popmuted] (0.6,0) arc (0:33.7:0.6);
    \node[popmuted] at (0.75,0.22) {$\alpha$};
    \draw[popmuted] (1.1,0) arc (0:95:1.1);
    \node[popmuted] at (0.1,1.25) {$\alpha+\theta$};
    \draw[popdark, ->] (0,0) -- (1,0) node[below right] {$\vec{i}$};
    \draw[popdark, ->] (0,0) -- (0,1) node[above left] {$\vec{j}$};
    \draw[popgreen, dashed, ->] (0,0) -- (0.64,-0.77) node[below right] {$R_\theta(\vec{i})$};
    \draw[popgreen, dashed, ->] (0,0) -- (0.77,0.64) node[above right] {$R_\theta(\vec{j})$};
\end{tikzpicture}
\legende{Rotation through $\theta$: the images of $\vec{i}$ and $\vec{j}$ are the columns of the rotation matrix.}
\end{popfigure}

\begin{exemplebox}[Reflection in a Line through the Origin]
Reflection in the line $y = x$: $S(x,y) = (y,x)$, matrix $\begin{pmatrix} 0 & 1 \\ 1 & 0 \end{pmatrix}$.
Reflection in the $x$-axis: $S_x(x,y) = (x,-y)$, matrix $\begin{pmatrix} 1 & 0 \\ 0 & -1 \end{pmatrix}$.
Reflection in the line making angle $\alpha$ with the positive $x$-axis: matrix $\begin{pmatrix} \cos2\alpha & \sin2\alpha \\ \sin2\alpha & -\cos2\alpha \end{pmatrix}$.
\end{exemplebox}

\begin{exemplebox}[Orthogonal Projection onto a Line through the Origin]
Projection onto the $x$-axis: $P_x(x,y) = (x,0)$, matrix $\begin{pmatrix} 1 & 0 \\ 0 & 0 \end{pmatrix}$.
Projection onto the line $y = mx$: matrix $\dfrac{1}{1+m^2}\begin{pmatrix} 1 & m \\ m & m^2 \end{pmatrix}$.
\end{exemplebox}

\begin{exemplebox}[Shear Transformations]
Horizontal shear: $S_h(x,y) = (x+ky, y)$, matrix $\begin{pmatrix} 1 & k \\ 0 & 1 \end{pmatrix}$.
Vertical shear: $S_v(x,y) = (x, y+kx)$, matrix $\begin{pmatrix} 1 & 0 \\ k & 1 \end{pmatrix}$.
\end{exemplebox}

\begin{popfigure}
\begin{tikzpicture}[scale=1.6, >=stealth]
    \draw[->] (-0.5,0) -- (2.6,0) node[right] {$x$};
    \draw[->] (0,-0.5) -- (0,1.6) node[above] {$y$};
    \draw[popblue, thick] (0,0) -- (1,0) -- (1,1) -- (0,1) -- cycle;
    \node[popblue] at (0.4,0.5) {unit square};
    \draw[poporange, thick] (0,0) -- (1,0) -- (2,1) -- (1,1) -- cycle;
    \node[poporange] at (1.6,0.55) {image};
    \draw[popdark, ->, thick] (0,0) -- (1,0) node[midway, below] {$T(\vec{i})=\vec{i}$};
    \draw[poporange, ->, thick] (0,0) -- (1,1) node[midway, above left] {$T(\vec{j})$};
    \foreach \t in {0.25,0.5,0.75} {
        \draw[popblue, thin] (0,\t) -- (1,\t);
        \draw[poporange, thin] (\t,\t) -- (1+\t,\t);
    }
\end{tikzpicture}
\legende{Horizontal shear $T(x,y)=(x+y,y)$: the unit square becomes a parallelogram; horizontal lines stay horizontal and parallel.}
\end{popfigure}

\courssub{Matrix Representation of a Linear Transformation in the Plane}

All the examples above have the form $T(x,y) = (ax+by,\; cx+dy)$. This is no accident.

\begin{propriete}[Matrix of a Linear Transformation of $\mathbb{R}^2$]
Every linear transformation $T : \mathbb{R}^2 \to \mathbb{R}^2$ can be written uniquely as
\[
T(x,y) = (ax+by,\; cx+dy)
\]
for some real numbers $a,b,c,d$. The \cle{matrix of $T$} relative to the standard basis $\{\vec{i},\vec{j}\}$ is
\[
\mathbf{M}_T = \begin{pmatrix} a & b \\ c & d \end{pmatrix},
\]
whose \textbf{first column is $T(\vec{i})$} and whose \textbf{second column is $T(\vec{j})$}.
\end{propriete}

\textbf{Proof.}\par
Write $T(\vec{i}) = (a,c)$ and $T(\vec{j}) = (b,d)$. For any $(x,y) = x\vec{i}+y\vec{j}$, linearity gives
\[
T(x,y) = xT(\vec{i})+yT(\vec{j}) = x(a,c)+y(b,d) = (ax+by,\; cx+dy).
\]
Conversely, any map of this form satisfies the combined linearity condition, since each component is a linear form in $(x,y)$. \hfill$\square$

\begin{methode}[Building the Matrix from a Geometric Description]
To find the matrix of a linear transformation $T$ described geometrically:
\begin{enumerate}
    \item Determine the images of the basis vectors $\vec{i}=(1,0)$ and $\vec{j}=(0,1)$.
    \item Write these images as column vectors.
    \item The matrix is $\mathbf{M}_T = \bigl[\,T(\vec{i})\;\; T(\vec{j})\,\bigr]$.
\end{enumerate}
\end{methode}

\begin{exoresolu}[Finding the Matrix of a Rotation through $60^\circ$]
\noindent\textbf{Statement:} Find the matrix of the rotation $R_{60^\circ}$ about the origin.
\tcblower
\noindent\textbf{Solution:}
$\cos60^\circ = \tfrac12$ and $\sin60^\circ = \tfrac{\sqrt3}{2}$.
$R_{60^\circ}(\vec{i}) = \bigl(\tfrac12,\; \tfrac{\sqrt3}{2}\bigr)$ and $R_{60^\circ}(\vec{j}) = \bigl(-\tfrac{\sqrt3}{2},\; \tfrac12\bigr)$.
Hence
\[
\mathbf{M} = \begin{pmatrix} \frac12 & -\frac{\sqrt3}{2} \\[2pt] \frac{\sqrt3}{2} & \frac12 \end{pmatrix}.
\]
\end{exoresolu}

\courssub{Extension to Three-Dimensional Space}

Everything generalises naturally to $\mathbb{R}^3$ with $\vec{i}=(1,0,0)$, $\vec{j}=(0,1,0)$, $\vec{k}=(0,0,1)$.

\begin{propriete}[Matrix of a Linear Transformation of $\mathbb{R}^3$]
Every linear transformation $T : \mathbb{R}^3 \to \mathbb{R}^3$ has the form
\[
T(x,y,z) = (a_{11}x + a_{12}y + a_{13}z,\; a_{21}x + a_{22}y + a_{23}z,\; a_{31}x + a_{32}y + a_{33}z)
\]
and is represented by the $3\times3$ matrix
\[
\mathbf{M}_T = \begin{pmatrix}
a_{11} & a_{12} & a_{13} \\
a_{21} & a_{22} & a_{23} \\
a_{31} & a_{32} & a_{33}
\end{pmatrix}
= \bigl[\,T(\vec{i})\;\; T(\vec{j})\;\; T(\vec{k})\,\bigr].
\]
\end{propriete}

Common examples in $\mathbb{R}^3$:
\begin{itemize}
    \item Rotation about the $z$-axis through $\theta$: $\begin{pmatrix} \cos\theta & -\sin\theta & 0 \\ \sin\theta & \cos\theta & 0 \\ 0 & 0 & 1 \end{pmatrix}$.
    \item Reflection in the plane $z=0$: $\begin{pmatrix} 1 & 0 & 0 \\ 0 & 1 & 0 \\ 0 & 0 & -1 \end{pmatrix}$.
    \item Homothety of ratio $k$: $k\mathbf{I}_3$.
\end{itemize}

\courssub{Operations on Linear Transformations}

Linear transformations can be added, multiplied by scalars, and composed; these operations correspond exactly to the usual matrix operations.

\begin{propriete}[Algebra of Linear Transformations]
Let $T, S : E \to F$ be linear and $\lambda \in \mathbb{R}$.
\begin{enumerate}
    \item \textbf{Sum:} $(T+S)(\vec{u}) = T(\vec{u}) + S(\vec{u})$ is linear, and $\mathbf{M}_{T+S} = \mathbf{M}_T + \mathbf{M}_S$.
    \item \textbf{Scalar multiple:} $(\lambda T)(\vec{u}) = \lambda\,T(\vec{u})$ is linear, and $\mathbf{M}_{\lambda T} = \lambda\mathbf{M}_T$.
    \item \textbf{Composition:} if $T : E \to F$ and $S : F \to G$ are linear, then $S \circ T : E \to G$, defined by $(S\circ T)(\vec{u}) = S(T(\vec{u}))$, is linear, and $\mathbf{M}_{S\circ T} = \mathbf{M}_S \mathbf{M}_T$.
\end{enumerate}
\end{propriete}

\begin{attention}[Order of Composition]
$S \circ T$ means \textbf{apply $T$ first, then $S$}. The matrix is $\mathbf{M}_S \mathbf{M}_T$. In general $\mathbf{M}_S\mathbf{M}_T \neq \mathbf{M}_T\mathbf{M}_S$: composition of transformations, like matrix multiplication, is not commutative.
\end{attention}

\begin{exemplebox}[Composition of Two Shears]
Let $T_1(x,y) = (x+y, y)$ with matrix $\begin{pmatrix}1&1\\0&1\end{pmatrix}$ and $T_2(x,y) = (x, y+x)$ with matrix $\begin{pmatrix}1&0\\1&1\end{pmatrix}$.
\[
\mathbf{M}_{T_2\circ T_1} = \begin{pmatrix}1&0\\1&1\end{pmatrix}\begin{pmatrix}1&1\\0&1\end{pmatrix} = \begin{pmatrix}1&1\\1&2\end{pmatrix},
\qquad
\mathbf{M}_{T_1\circ T_2} = \begin{pmatrix}1&1\\0&1\end{pmatrix}\begin{pmatrix}1&0\\1&1\end{pmatrix} = \begin{pmatrix}2&1\\1&1\end{pmatrix}.
\]
The two composites are different.
\end{exemplebox}

\courssub{Recognising Non-Linear Maps — GCE Examination Traps}

A mapping $T : \mathbb{R}^2 \to \mathbb{R}^2$ given by algebraic formulas is \textbf{linear} if and only if each component is a \textbf{homogeneous linear polynomial} in the input variables: a sum of terms of the form $a x_i$, with no constant term, no products and no powers of variables.

\begin{methode}[How to Test Linearity Quickly]
\begin{enumerate}
    \item \textbf{Zero test:} compute $T(\vec{0})$. If $T(\vec{0}) \neq \vec{0}$, stop: $T$ is \textbf{not} linear.
    \item \textbf{Form inspection:} constant terms, products ($xy$, $yz$), powers ($x^2$), or non-linear functions ($\sin x$, $e^x$, $\sqrt{x}$, $|x|$) $\Rightarrow$ not linear.
    \item Otherwise, verify the combined condition $T(a\vec{u}+b\vec{v}) = aT(\vec{u})+bT(\vec{v})$ algebraically, or exhibit the matrix.
\end{enumerate}
\end{methode}

\begin{exoresolu}[Identifying Linear and Non-Linear Maps]
\noindent\textbf{Statement:} Which of the following transformations of $\mathbb{R}^2$ are linear?
\begin{enumerate}
    \item $T_1(x,y) = (2x - y,\; 3x + 4y)$
    \item $T_2(x,y) = (x+1,\; y-2)$
    \item $T_3(x,y) = (x^2,\; y)$
    \item $T_4(x,y) = (x+y,\; xy)$
    \item $T_5(x,y) = (0,\; 0)$
\end{enumerate}
\tcblower
\noindent\textbf{Solution:}
\begin{enumerate}
    \item \textbf{Linear.} Each component is a homogeneous linear polynomial; matrix $\begin{pmatrix}2&-1\\3&4\end{pmatrix}$.
    \item \textbf{Not linear.} $T_2(0,0) = (1,-2) \neq (0,0)$ (constant terms).
    \item \textbf{Not linear.} $T_3(0,0)=(0,0)$, but $T_3(2,0)=(4,0)$ while $2T_3(1,0)=(2,0)$: homogeneity fails (term $x^2$).
    \item \textbf{Not linear.} $T_4(1,1)=(2,1)$ while $T_4(1,0)+T_4(0,1)=(1,0)+(1,0)=(2,0)$: additivity fails (term $xy$).
    \item \textbf{Linear.} This is the zero transformation.
\end{enumerate}
\end{exoresolu}

\begin{popfigure}
\begin{tikzpicture}[scale=0.75, >=stealth]
    % Linear map
    \foreach \x in {-2,-1,0,1,2} \draw[popblue, thin] (\x,-2) -- (\x,2);
    \foreach \y in {-2,-1,0,1,2} \draw[popblue, thin] (-2,\y) -- (2,\y);
    \draw[popdark, ->, thick] (0,0) -- (1,0);
    \draw[popdark, ->, thick] (0,0) -- (0,1);
    \node at (0,-2.7) {original grid};
    \begin{scope}[xshift=7cm]
        \foreach \x in {-4,-2,0,2,4} \draw[poporange, thin] (\x,-2) -- (\x,2);
        \foreach \y in {-2,-1,0,1,2} \draw[poporange, thin] (-4,\y) -- (4,\y);
        \draw[popdark, ->, thick] (0,0) -- (2,0);
        \draw[popdark, ->, thick] (0,0) -- (0,1);
        \node at (0,-2.7) {image under $T(x,y)=(2x,y)$};
    \end{scope}
    \draw[->, popmuted, thick] (2.6,0) -- (4.4,0) node[midway, above] {$T$};
    % Non-linear map
    \begin{scope}[yshift=-7cm]
        \foreach \x in {-2,-1,0,1,2} \draw[popblue, thin] (\x,-2) -- (\x,2);
        \foreach \y in {-2,-1,0,1,2} \draw[popblue, thin] (-2,\y) -- (2,\y);
        \draw[popdark, ->, thick] (0,0) -- (1,0);
        \draw[popdark, ->, thick] (0,0) -- (0,1);
        \node at (0,-2.7) {original grid};
    \end{scope}
    \begin{scope}[xshift=7cm, yshift=-7cm]
        \foreach \x in {0,0.25,1,2.25,4} \draw[poporange, thin] (\x,-2) -- (\x,2);
        \foreach \y in {-2,-1,0,1,2} \draw[poporange, thin] (0,\y) -- (4,\y);
        \draw[popdark, ->, thick] (0,0) -- (1,0);
        \draw[popdark, ->, thick] (0,0) -- (0,1);
        \node at (2,-2.7) {image under $N(x,y)=(x^2,y)$};
    \end{scope}
    \draw[->, popmuted, thick] (2.6,-7) -- (4.4,-7) node[midway, above] {$N$};
\end{tikzpicture}
\legende{Top: a linear map sends a rectangular grid to a parallelogram grid. Bottom: the non-linear map $N(x,y)=(x^2,y)$ collapses the lines $x=c$ and $x=-c$ onto the same line and distorts the spacing.}
\end{popfigure}

\begin{aretenir}[Key Points to Remember]
\begin{itemize}
    \item $T$ is linear $\iff$ $T(a\vec{u}+b\vec{v}) = aT(\vec{u})+bT(\vec{v})$ for all vectors and scalars.
    \item Every linear map sends $\vec{0}$ to $\vec{0}$; if $T(\vec{0})\neq\vec{0}$, $T$ is not linear.
    \item In $\mathbb{R}^2$, $T(x,y) = (ax+by,\; cx+dy)$ with matrix $\begin{pmatrix}a&b\\c&d\end{pmatrix}$ whose columns are $T(\vec{i})$ and $T(\vec{j})$.
    \item In $\mathbb{R}^3$, the matrix is $3\times3$ with columns $T(\vec{i}), T(\vec{j}), T(\vec{k})$.
    \item Composition corresponds to matrix multiplication: $\mathbf{M}_{S\circ T} = \mathbf{M}_S\mathbf{M}_T$.
    \item \textbf{Examination trap:} constant terms, squares, products ($xy$) or transcendental functions in the formula $\Rightarrow$ NOT linear.
\end{itemize}
\end{aretenir}

\begin{exercice}[Practice Exercise 1]
\begin{enumerate}
    \item Determine whether the following maps are linear. Justify your answer.
        \begin{enumerate}
            \item $T(x,y) = (3x-2y,\; x+5y)$
            \item $T(x,y) = (x+2,\; y-3)$
            \item $T(x,y) = (x^2+y^2,\; 0)$
            \item $T(x,y) = (|x|,\; y)$
        \end{enumerate}
    \item Find the matrix of the reflection of $\mathbb{R}^2$ in the line $y = -x$.
    \item A linear transformation $T : \mathbb{R}^2 \to \mathbb{R}^2$ satisfies $T(1,2) = (3,4)$ and $T(2,1) = (5,6)$. Find $T(x,y)$ and its matrix.
    \item Let $R_\theta$ be the rotation through $\theta$ and $S$ the reflection in the $x$-axis. Write the matrices of $R_\theta \circ S$ and $S \circ R_\theta$. Are they equal?
\end{enumerate}
\end{exercice}

\begin{corrige}[Exercise 1]
\begin{enumerate}
    \item \textbf{(a) Linear:} homogeneous linear components; matrix $\begin{pmatrix}3&-2\\1&5\end{pmatrix}$.
          \textbf{(b) Not linear:} $T(0,0)=(2,-3)\neq(0,0)$.
          \textbf{(c) Not linear:} $T(2,0)=(4,0)$ but $2T(1,0)=(2,0)$; homogeneity fails.
          \textbf{(d) Not linear:} $T(-1,0)=(1,0)$ but $-T(1,0)=(-1,0)$; the absolute value breaks homogeneity.
    \item $(1,0)\mapsto(0,-1)$ and $(0,1)\mapsto(-1,0)$, so the matrix is $\begin{pmatrix}0&-1\\-1&0\end{pmatrix}$.
    \item Write $(x,y) = \alpha(1,2)+\beta(2,1)$. Solving $\alpha+2\beta=x$, $2\alpha+\beta=y$ gives $\alpha=\dfrac{2y-x}{3}$, $\beta=\dfrac{2x-y}{3}$.
        Then
        \begin{align*}
        T(x,y) &= \alpha(3,4)+\beta(5,6)\\
        &= \tfrac{1}{3}\bigl(3(2y-x)+5(2x-y),\; 4(2y-x)+6(2x-y)\bigr)\\
        &= \tfrac{1}{3}(7x+y,\; 8x+2y).
        \end{align*}
        Matrix: $\dfrac{1}{3}\begin{pmatrix}7&1\\8&2\end{pmatrix}$. Check: $T(1,2)=\tfrac13(9,12)=(3,4)$ and $T(2,1)=\tfrac13(15,18)=(5,6)$. \hfill$\square$
    \item $\mathbf{M}_{R_\theta\circ S} = \begin{pmatrix}\cos\theta&-\sin\theta\\\sin\theta&\cos\theta\end{pmatrix}\begin{pmatrix}1&0\\0&-1\end{pmatrix} = \begin{pmatrix}\cos\theta&\sin\theta\\\sin\theta&-\cos\theta\end{pmatrix}$;
        $\mathbf{M}_{S\circ R_\theta} = \begin{pmatrix}1&0\\0&-1\end{pmatrix}\begin{pmatrix}\cos\theta&-\sin\theta\\\sin\theta&\cos\theta\end{pmatrix} = \begin{pmatrix}\cos\theta&-\sin\theta\\-\sin\theta&-\cos\theta\end{pmatrix}$.
        They are \textbf{not equal} unless $\sin\theta=0$.
\end{enumerate}
\end{corrige}

\coursec{Kernel and Image of a Linear Transformation}

A projection of $\mathbb{R}^3$ onto a plane squashes a whole direction onto the origin: infinitely many vectors share the same image. To understand a linear transformation we must answer two questions: \emph{which vectors are sent to zero?} and \emph{which vectors can be reached as images?} The answers are the \cle{kernel} and the \cle{image}.

\courssub{Definitions}

\begin{definition}[Kernel and Image]
Let $T : E \to F$ be a linear transformation.
\begin{enumerate}
    \item The \cle{kernel} of $T$ is the set of all vectors of $E$ mapped to $\vec{0}_F$:
    \[
    \ker T = \{\vec{u} \in E : T(\vec{u}) = \vec{0}_F\}.
    \]
    \item The \cle{image} of $T$ is the set of all vectors of $F$ that are images of vectors of $E$:
    \[
    \operatorname{Im} T = \{T(\vec{u}) : \vec{u} \in E\}.
    \]
\end{enumerate}
\end{definition}

\begin{propriete}[Kernel and Image are Subspaces]
$\ker T$ is a vector subspace of $E$ and $\operatorname{Im} T$ is a vector subspace of $F$. In particular both contain the zero vector.
\end{propriete}

\textbf{Proof.}\par
$T(\vec{0}_E)=\vec{0}_F$, so $\vec{0}_E\in\ker T$. If $\vec{u},\vec{v}\in\ker T$ and $a,b\in\mathbb{R}$, then $T(a\vec{u}+b\vec{v}) = aT(\vec{u})+bT(\vec{v}) = \vec{0}$, so $a\vec{u}+b\vec{v}\in\ker T$. Similarly, if $\vec{w}_1=T(\vec{u}_1)$ and $\vec{w}_2=T(\vec{u}_2)$, then $a\vec{w}_1+b\vec{w}_2 = T(a\vec{u}_1+b\vec{u}_2)\in\operatorname{Im}T$. \hfill$\square$

\begin{propriete}[The Image is Spanned by the Columns]
If $T : \mathbb{R}^2 \to \mathbb{R}^2$ has matrix $\mathbf{M}_T = \bigl[\,T(\vec{i})\;\;T(\vec{j})\,\bigr]$, then
\[
\operatorname{Im} T = \operatorname{span}\{T(\vec{i}),\, T(\vec{j})\},
\]
the set of all linear combinations of the columns of $\mathbf{M}_T$. The same holds in $\mathbb{R}^3$ with the three columns.
\end{propriete}

\textbf{Proof.}\par
$T(x,y) = xT(\vec{i}) + yT(\vec{j})$, so every image is a linear combination of the two columns, and conversely every such combination is an image. \hfill$\square$

\begin{methode}[Finding Kernel and Image from the Matrix]
Given $T(x,y) = (ax+by,\; cx+dy)$:
\begin{enumerate}
    \item \textbf{Kernel:} solve the homogeneous system $ax+by=0$, $cx+dy=0$. The solution set is $\ker T$.
    \item \textbf{Image:} take the span of the columns $T(\vec{i})=(a,c)$ and $T(\vec{j})=(b,d)$; reduce to a simplest description (a point, a line through $O$, or the whole plane).
\end{enumerate}
\end{methode}

\begin{exoresolu}[Kernel and Image of a Projection]
\noindent\textbf{Statement:} Let $T(x,y) = (x+y,\; 2x+2y)$. Determine $\ker T$ and $\operatorname{Im} T$.
\tcblower
\noindent\textbf{Solution:}
\textbf{Kernel:} solve $x+y=0$ and $2x+2y=0$; both reduce to $y=-x$. Hence
\[
\ker T = \{(t,-t) : t\in\mathbb{R}\} = \operatorname{span}\{(1,-1)\},
\]
the line $y=-x$.
\textbf{Image:} the columns are $T(\vec{i})=(1,2)$ and $T(\vec{j})=(1,2)$, so
\[
\operatorname{Im} T = \operatorname{span}\{(1,2)\} = \{(s,2s) : s\in\mathbb{R}\},
\]
the line $y=2x$.
Note that $\dim\ker T + \dim\operatorname{Im} T = 1+1 = 2 = \dim\mathbb{R}^2$.
\end{exoresolu}

\courssub{The Rank--Nullity Theorem}

\begin{definition}[Rank and Nullity]
The \cle{rank} of $T$ is $\operatorname{rank}(T) = \dim(\operatorname{Im} T)$; the \cle{nullity} of $T$ is $\dim(\ker T)$.
\end{definition}

\begin{propriete}[Rank--Nullity Theorem]
For any linear transformation $T : E \to F$ with $E$ finite-dimensional,
\[
\dim(\ker T) + \dim(\operatorname{Im} T) = \dim E.
\]
In the plane: $\dim\ker T + \operatorname{rank}(T) = 2$; in space: $\dim\ker T + \operatorname{rank}(T) = 3$.
\end{propriete}

\begin{exemplebox}[The Three Possibilities in the Plane]
For $T : \mathbb{R}^2 \to \mathbb{R}^2$:
\begin{itemize}
    \item $\operatorname{rank}(T)=2$: $\ker T = \{\vec{0}\}$, $\operatorname{Im} T = \mathbb{R}^2$ ($T$ is invertible);
    \item $\operatorname{rank}(T)=1$: $\ker T$ and $\operatorname{Im} T$ are both lines through $O$;
    \item $\operatorname{rank}(T)=0$: $T$ is the zero map, $\ker T = \mathbb{R}^2$.
\end{itemize}
\end{exemplebox}

\begin{exoresolu}[Kernel and Image in $\mathbb{R}^3$]
\noindent\textbf{Statement:} Let $T(x,y,z) = (x+y+z,\; x-y,\; 2x+z)$. Find $\ker T$ and $\operatorname{Im} T$, and verify the rank--nullity theorem.
\tcblower
\noindent\textbf{Solution:}
\textbf{Kernel:} solve
\[
\begin{cases} x+y+z = 0 \\ x - y = 0 \\ 2x + z = 0 \end{cases}
\Rightarrow y = x,\quad z = -2x,\quad x+x-2x = 0.
\]
All three equations are satisfied by $(t, t, -2t)$, so $\ker T = \operatorname{span}\{(1,1,-2)\}$, a line through $O$; nullity $=1$.
\textbf{Image:} the columns are $(1,1,2)$, $(1,-1,0)$, $(1,0,1)$. Since $(1,0,1) = \tfrac12(1,1,2) + \tfrac12(1,-1,0)$, the image is spanned by the first two columns, which are not proportional:
\[
\operatorname{Im} T = \operatorname{span}\{(1,1,2),\, (1,-1,0)\},
\]
a plane through $O$; rank $=2$. Check: $1 + 2 = 3 = \dim\mathbb{R}^3$. \hfill$\square$
\end{exoresolu}

\begin{aretenir}[Key Points]
\begin{itemize}
    \item $\ker T = \{\vec{u} : T(\vec{u})=\vec{0}\}$: solve the homogeneous system given by the matrix.
    \item $\operatorname{Im} T$ is the span of the columns of $\mathbf{M}_T$.
    \item Rank--nullity: $\dim\ker T + \dim\operatorname{Im} T = \dim E$.
    \item $T$ is one-to-one (injective) $\iff \ker T = \{\vec{0}\}$.
\end{itemize}
\end{aretenir}

\coursec{Inverse of a Linear Transformation and Invariant Points}

\courssub{Invertible Linear Transformations}

\begin{definition}[Inverse Transformation]
A linear transformation $T : E \to E$ is \cle{invertible} if there exists a map $T^{-1} : E \to E$ such that
\[
T^{-1}\circ T = T \circ T^{-1} = I.
\]
$T^{-1}$ is then automatically linear, and $\mathbf{M}_{T^{-1}} = \mathbf{M}_T^{-1}$.
\end{definition}

\begin{propriete}[Invertibility Criterion]
For $T : \mathbb{R}^2 \to \mathbb{R}^2$ with matrix $\mathbf{M}$, the following are equivalent:
\begin{enumerate}
    \item $T$ is invertible;
    \item $\det \mathbf{M} \neq 0$;
    \item $\ker T = \{\vec{0}\}$;
    \item $\operatorname{Im} T = \mathbb{R}^2$.
\end{enumerate}
When $\mathbf{M} = \begin{pmatrix} a & b \\ c & d \end{pmatrix}$ with $ad-bc\neq 0$,
\[
\mathbf{M}^{-1} = \frac{1}{ad-bc}\begin{pmatrix} d & -b \\ -c & a \end{pmatrix}.
\]
\end{propriete}

\begin{exemplebox}[Inverses of the Standard Transformations]
\begin{itemize}
    \item Rotation: $R_\theta^{-1} = R_{-\theta}$ (rotate back).
    \item Homothety: $H_k^{-1} = H_{1/k}$ for $k\neq 0$.
    \item Reflection: $S^{-1} = S$ (a reflection is its own inverse: $S\circ S = I$).
    \item A projection $P$ (other than $I$) is \textbf{never} invertible: $\det \mathbf{M}_P = 0$ and $\ker P \neq \{\vec{0}\}$.
\end{itemize}
\end{exemplebox}

\begin{exoresolu}[Computing an Inverse Transformation]
\noindent\textbf{Statement:} Let $T(x,y) = (2x+y,\; 3x+2y)$. Show that $T$ is invertible and find $T^{-1}(x,y)$.
\tcblower
\noindent\textbf{Solution:}
$\mathbf{M} = \begin{pmatrix}2&1\\3&2\end{pmatrix}$, $\det\mathbf{M} = 4-3 = 1 \neq 0$, so $T$ is invertible.
\[
\mathbf{M}^{-1} = \frac{1}{1}\begin{pmatrix}2&-1\\-3&2\end{pmatrix},
\qquad\text{hence}\qquad
T^{-1}(x,y) = (2x - y,\; -3x + 2y).
\]
Check: $T^{-1}(T(1,1)) = T^{-1}(3,5) = (6-5,\, -9+10) = (1,1)$. \hfill$\square$
\end{exoresolu}

\courssub{Invariant (Fixed) Points}

\begin{definition}[Invariant Point]
A point $P$ with position vector $\vec{u}$ is an \cle{invariant point} (or \emph{fixed point}) of $T$ if $T(\vec{u}) = \vec{u}$: the point is left unmoved by the transformation.
\end{definition}

Since $T(\vec{0}) = \vec{0}$, the origin is always an invariant point of a linear transformation. The interesting question is whether there are \emph{others}. The invariant points are exactly the solutions of
\[
(\mathbf{M} - \mathbf{I})\vec{u} = \vec{0},
\qquad\text{i.e.}\qquad
\ker(T - I).
\]
The set of invariant points is therefore always a subspace: $\{\vec{0}\}$, a line through $O$, a plane through $O$, or the whole space.

\begin{methode}[Finding Invariant Points]
\begin{enumerate}
    \item Write the condition $T(x,y) = (x,y)$.
    \item Bring everything to one side: this gives a homogeneous system $(\mathbf{M}-\mathbf{I})\begin{pmatrix}x\\y\end{pmatrix} = \begin{pmatrix}0\\0\end{pmatrix}$.
    \item Solve; describe the solution set geometrically.
\end{enumerate}
\end{methode}

\begin{exoresolu}[Invariant Points of a Reflection and of a Shear]
\noindent\textbf{Statement:} Find the invariant points of (a) the reflection $S$ in the line $y=x$; (b) the shear $T(x,y) = (x+2y,\; y)$.
\tcblower
\noindent\textbf{Solution:}
\textbf{(a)} $S(x,y) = (y,x) = (x,y) \iff y = x$. The invariant points form the line $y=x$ — exactly the mirror line, as expected.
\textbf{(b)} $T(x,y) = (x,y) \iff x+2y = x \iff y = 0$. The invariant points form the $x$-axis: every point of the $x$-axis is fixed, every other point slides horizontally. \hfill$\square$
\end{exoresolu}

\begin{attention}[Invariant Points vs Invariant Lines]
Do not confuse the two notions:
\begin{itemize}
    \item An \textbf{invariant point} satisfies $T(\vec{u}) = \vec{u}$: the point does not move.
    \item A line $\ell$ is an \textbf{invariant line} if $T(\ell) \subseteq \ell$: the line is mapped \emph{into itself}, but its individual points may move along it. For the homothety $H_2$, every line through $O$ is invariant, yet the only invariant \emph{point} is $O$.
\end{itemize}
\end{attention}

\begin{exoresolu}[A Line of Invariant Points]
\noindent\textbf{Statement:} Find the invariant points of $T(x,y) = (3x-2y,\; 2x-y)$.
\tcblower
\noindent\textbf{Solution:}
$T(x,y) = (x,y)$ gives
\[
\begin{cases} 3x - 2y = x \\ 2x - y = y \end{cases}
\iff
\begin{cases} 2x - 2y = 0 \\ 2x - 2y = 0 \end{cases}
\iff x = y.
\]
The invariant points form the line $y=x$. (Indeed $\mathbf{M}-\mathbf{I} = \begin{pmatrix}2&-2\\2&-2\end{pmatrix}$ has rank $1$, so its kernel is a line.) \hfill$\square$
\end{exoresolu}

\coursec{Images of Lines and Planes}

A photograph enlarged by a printer keeps straight lines straight: this is a general fact about linear maps. We now make it precise and learn how to compute the image of a given line or plane.

\courssub{Image of a Line}

\begin{propriete}[Linear Maps Preserve Lines]
Let $T : \mathbb{R}^2 \to \mathbb{R}^2$ be linear and let $\ell$ be the line through the point $A$ with direction vector $\vec{d}$, i.e. $\ell = \{A + t\vec{d} : t\in\mathbb{R}\}$. Then
\[
T(\ell) = \{T(A) + t\,T(\vec{d}) : t\in\mathbb{R}\}.
\]
\begin{itemize}
    \item If $T(\vec{d}) \neq \vec{0}$, then $T(\ell)$ is the \textbf{line} through $T(A)$ with direction $T(\vec{d})$.
    \item If $T(\vec{d}) = \vec{0}$, then $T(\ell)$ collapses to the single \textbf{point} $T(A)$.
\end{itemize}
In particular, the image of a line through the origin is a line through the origin (or just $\{O\}$).
\end{propriete}

\textbf{Proof.}\par
$T(A + t\vec{d}) = T(A) + tT(\vec{d})$ by linearity; as $t$ ranges over $\mathbb{R}$ this describes exactly the stated set. \hfill$\square$

\begin{methode}[Finding the Image of a Line]
Given a line $\ell$ and a linear map $T$:
\begin{enumerate}
    \item Write $\ell$ in parametric form: pick a point $A$ on $\ell$ and a direction vector $\vec{d}$.
    \item Compute $T(A)$ and $T(\vec{d})$.
    \item If $T(\vec{d})\neq\vec{0}$, write the Cartesian equation of the line through $T(A)$ in direction $T(\vec{d})$.
\end{enumerate}
Alternative for a line $ax+by=c$: express $x,y$ from the parametric form, substitute into $T$, then eliminate the parameter.
\end{methode}

\begin{exoresolu}[Image of a Line under a Linear Map]
\noindent\textbf{Statement:} Let $T(x,y) = (x+2y,\; 3x-y)$. Find the image of the line $\ell : x + y = 1$.
\tcblower
\noindent\textbf{Solution:}
Parametrise $\ell$: $A = (1,0)$ lies on $\ell$ and $\vec{d} = (1,-1)$ is a direction vector, so points of $\ell$ are $(1+t,\; -t)$.
\[
T(1+t,\, -t) = \bigl((1+t) + 2(-t),\; 3(1+t) - (-t)\bigr) = (1 - t,\; 3 + 4t).
\]
So the image is the line through $T(A) = (1,3)$ with direction $T(\vec{d}) = (-1, 4)$. Eliminating $t$: from $x = 1-t$ we get $t = 1-x$, and substituting in $y = 3+4t$ gives
\[
y = 3 + 4(1-x) \iff 4x + y = 7.
\]
Hence $T(\ell)$ is the line $4x + y = 7$. \hfill$\square$
\end{exoresolu}

\begin{exoresolu}[A Line that Collapses to a Point]
\noindent\textbf{Statement:} Let $T(x,y) = (x+y,\; 2x+2y)$ (as in a previous example, $\ker T$ is the line $y=-x$). Find the image of the line $\ell : y = -x + 3$.
\tcblower
\noindent\textbf{Solution:}
Points of $\ell$: $(t,\; 3-t)$. Then
\[
T(t,\, 3-t) = \bigl(t + 3 - t,\; 2t + 2(3-t)\bigr) = (3, 6).
\]
Every point of $\ell$ has the same image: $T(\ell) = \{(3,6)\}$, a single point. This happens because the direction $(1,-1)$ of $\ell$ lies in $\ker T$. \hfill$\square$
\end{exoresolu}

\courssub{Image of a Plane in $\mathbb{R}^3$}

\begin{propriete}[Image of a Plane]
Let $T : \mathbb{R}^3 \to \mathbb{R}^3$ be linear and let $\Pi$ be the plane through $A$ spanned by the independent directions $\vec{d}_1, \vec{d}_2$:
\[
\Pi = \{A + s\vec{d}_1 + t\vec{d}_2 : s,t\in\mathbb{R}\}.
\]
Then $T(\Pi) = \{T(A) + s\,T(\vec{d}_1) + t\,T(\vec{d}_2) : s,t\in\mathbb{R}\}$, which is:
\begin{itemize}
    \item a \textbf{plane} if $T(\vec{d}_1)$ and $T(\vec{d}_2)$ are not proportional;
    \item a \textbf{line} if they are proportional (and not both zero);
    \item a \textbf{point} if both are zero.
\end{itemize}
\end{propriete}

\begin{exoresolu}[Image of a Plane under a Projection]
\noindent\textbf{Statement:} Let $T(x,y,z) = (x,\; y,\; 0)$ be the orthogonal projection onto the plane $z=0$. Find the image of the plane $\Pi : x + y + z = 3$.
\tcblower
\noindent\textbf{Solution:}
Points of $\Pi$: $(x, y, 3-x-y)$. Then
\[
T(x,\, y,\, 3-x-y) = (x,\, y,\; 0).
\]
As $x$ and $y$ range freely over $\mathbb{R}$, the image is the whole plane $z=0$: $T(\Pi) = \{(x,y,0) : x,y\in\mathbb{R}\}$. The projection flattens the tilted plane $\Pi$ onto the coordinate plane $z=0$. \hfill$\square$
\end{exoresolu}

\begin{attention}[Common Mistakes]
\begin{itemize}
    \item The image of a line \emph{not} through the origin need not pass through the origin: compute $T(A)$ carefully.
    \item Do not forget the degenerate case: if the direction vector lies in $\ker T$, the image is a point, not a line.
    \item Parallel lines remain parallel (or collapse): $T$ sends direction $\vec{d}$ to $T(\vec{d})$ for \emph{every} line with that direction.
\end{itemize}
\end{attention}

\coursec{Eigenvalues and Eigenvectors}

Under a reflection in the line $y=x$, the vector $(1,1)$ is fixed and the vector $(1,-1)$ is flipped to its opposite: both keep their \emph{line of action}. Under a general linear map, most directions are rotated, but some special directions are preserved — the map merely stretches, compresses or reverses vectors along them. These are the \cle{eigendirections}, and they reveal the deep structure of the transformation.

\courssub{Definitions}

\begin{definition}[Eigenvalue and Eigenvector]
Let $T : E \to E$ be a linear transformation. A scalar $\lambda$ is an \cle{eigenvalue} of $T$ if there exists a \textbf{non-zero} vector $\vec{u}$ such that
\[
T(\vec{u}) = \lambda \vec{u}.
\]
Such a vector $\vec{u}$ is called an \cle{eigenvector} of $T$ associated with $\lambda$. The set
\[
E_\lambda = \{\vec{u} : T(\vec{u}) = \lambda\vec{u}\} = \ker(T - \lambda I)
\]
is the \cle{eigenspace} associated with $\lambda$; it is a subspace of $E$.
\end{definition}

\begin{attention}[The Zero Vector is Never an Eigenvector]
$T(\vec{0}) = \lambda\vec{0}$ holds for \emph{every} $\lambda$, so $\vec{0}$ is excluded from the definition of an eigenvector. However, $\lambda = 0$ \emph{can} be an eigenvalue: it is one exactly when $\ker T \neq \{\vec{0}\}$, i.e. when $T$ is not invertible.
\end{attention}

\courssub{The Characteristic Equation}

The condition $T(\vec{u}) = \lambda\vec{u}$ with $\vec{u}\neq\vec{0}$ means $(\mathbf{M} - \lambda\mathbf{I})\vec{u} = \vec{0}$ has a non-trivial solution, which happens exactly when the matrix $\mathbf{M} - \lambda\mathbf{I}$ is singular.

\begin{propriete}[Characteristic Equation]
The eigenvalues of $T$ with matrix $\mathbf{M}$ are the roots of the \cle{characteristic equation}
\[
\det(\mathbf{M} - \lambda\mathbf{I}) = 0.
\]
For $\mathbf{M} = \begin{pmatrix} a & b \\ c & d \end{pmatrix}$ this is the quadratic
\[
\lambda^2 - (a+d)\lambda + (ad - bc) = 0,
\qquad\text{i.e.}\qquad
\lambda^2 - (\operatorname{tr}\mathbf{M})\lambda + \det\mathbf{M} = 0,
\]
where $\operatorname{tr}\mathbf{M} = a+d$ is the \cle{trace} of $\mathbf{M}$.
\end{propriete}

\textbf{Proof.}\par
$\det(\mathbf{M}-\lambda\mathbf{I}) = \begin{vmatrix} a-\lambda & b \\ c & d-\lambda \end{vmatrix} = (a-\lambda)(d-\lambda) - bc = \lambda^2 - (a+d)\lambda + (ad-bc)$. \hfill$\square$

\begin{methode}[Finding Eigenvalues and Eigenvectors of a $2\times2$ Matrix]
\begin{enumerate}
    \item Compute $\operatorname{tr}\mathbf{M}$ and $\det\mathbf{M}$; write the characteristic equation $\lambda^2 - (\operatorname{tr}\mathbf{M})\lambda + \det\mathbf{M} = 0$.
    \item Solve the quadratic: the roots $\lambda_1, \lambda_2$ are the eigenvalues.
    \item For each eigenvalue $\lambda$, solve $(\mathbf{M} - \lambda\mathbf{I})\vec{u} = \vec{0}$ to find the eigenspace $E_\lambda$ (a line through $O$ when $\lambda$ is a simple root).
    \item \textbf{Check:} $\lambda_1 + \lambda_2 = \operatorname{tr}\mathbf{M}$ and $\lambda_1\lambda_2 = \det\mathbf{M}$.
\end{enumerate}
\end{methode}

\begin{exoresolu}[Eigenvalues and Eigenvectors — Two Distinct Real Eigenvalues]
\noindent\textbf{Statement:} Find the eigenvalues and eigenvectors of $T(x,y) = (4x + 2y,\; x + 3y)$.
\tcblower
\noindent\textbf{Solution:}
$\mathbf{M} = \begin{pmatrix}4&2\\1&3\end{pmatrix}$; $\operatorname{tr}\mathbf{M} = 7$, $\det\mathbf{M} = 12 - 2 = 10$.
Characteristic equation: $\lambda^2 - 7\lambda + 10 = 0 \iff (\lambda-2)(\lambda-5) = 0$, so $\lambda_1 = 2$, $\lambda_2 = 5$.
\textbf{For $\lambda_1 = 2$:} $(\mathbf{M} - 2\mathbf{I})\vec{u} = \vec{0}$:
\[
\begin{cases} 2x + 2y = 0 \\ x + y = 0 \end{cases} \iff y = -x,
\qquad E_2 = \operatorname{span}\{(1,-1)\}.
\]
\textbf{For $\lambda_2 = 5$:} $(\mathbf{M} - 5\mathbf{I})\vec{u} = \vec{0}$:
\[
\begin{cases} -x + 2y = 0 \\ x - 2y = 0 \end{cases} \iff x = 2y,
\qquad E_5 = \operatorname{span}\{(2,1)\}.
\]
Check: $T(1,-1) = (2,-2) = 2(1,-1)$ and $T(2,1) = (10,5) = 5(2,1)$. \hfill$\square$
\end{exoresolu}

\begin{exoresolu}[A Repeated Eigenvalue and No Real Eigenvalue]
\noindent\textbf{Statement:} Find the eigenvalues of (a) $T(x,y) = (x+2y,\; y)$ (shear); (b) $R_{90^\circ}(x,y) = (-y,\; x)$.
\tcblower
\noindent\textbf{Solution:}
\textbf{(a)} $\mathbf{M} = \begin{pmatrix}1&2\\0&1\end{pmatrix}$; characteristic equation $(1-\lambda)^2 = 0$, so $\lambda = 1$ is a \textbf{double} eigenvalue. Eigenspace: $2y = 0$, so $E_1 = \operatorname{span}\{(1,0)\}$, the $x$-axis only. A shear has a single eigendirection.
\textbf{(b)} $\mathbf{M} = \begin{pmatrix}0&-1\\1&0\end{pmatrix}$; characteristic equation $\lambda^2 + 1 = 0$, which has \textbf{no real root}. A rotation through $90^\circ$ has no real eigenvalue — geometrically obvious, since no direction of the plane is preserved. \hfill$\square$
\end{exoresolu}

\begin{popfigure}
\begin{tikzpicture}[scale=1.1, >=stealth]
    \draw[->] (-3,0) -- (3,0) node[right] {$x$};
    \draw[->] (0,-2.5) -- (0,2.5) node[above] {$y$};
    \draw[popline, thin] (-2.4,2.4) -- (2.4,-2.4);
    \draw[popline, thin] (-1.2,-0.6) -- (2.4,1.2);
    \draw[popblue, very thick, ->] (0,0) -- (1,-1) node[below right, popblue] {$\vec{u}_1$};
    \draw[popblue, very thick, ->] (0,0) -- (2,-2) node[below left, popblue] {$T(\vec{u}_1)=2\vec{u}_1$};
    \draw[poporange, very thick, ->] (0,0) -- (0.8,0.4) node[above, poporange] {$\vec{u}_2$};
    \draw[poporange, very thick, ->] (0,0) -- (2,1) node[above right, poporange] {$T(\vec{u}_2)=5\vec{u}_2$};
    \node[popmuted] at (-2.2,-1.9) {$E_2$};
    \node[popmuted] at (-1.1,-0.75) {$E_5$};
\end{tikzpicture}
\legende{The transformation $T(x,y)=(4x+2y,\;x+3y)$ doubles vectors along the line $E_2$ ($y=-x$) and multiplies by $5$ vectors along the line $E_5$ ($x=2y$).}
\end{popfigure}

\courssub{Eigenvalues of the Standard Transformations}

\begin{center}
\begin{tabularx}{\linewidth}{|l|X|X|}
\hline
\rowcolor{popblueL} \textbf{Transformation} & \textbf{Eigenvalues} & \textbf{Eigenspaces} \\
\hline
Homothety $H_k$ & $k$ (double) & the whole plane $\mathbb{R}^2$ \\
\hline
Rotation $R_\theta$ ($\theta \neq 0, \pi$) & no real eigenvalue & --- \\
\hline
Rotation $R_\pi$ (half-turn) & $-1$ (double) & the whole plane $\mathbb{R}^2$ \\
\hline
Reflection in a line $\ell$ & $1$ and $-1$ & $E_1 = \ell$; $E_{-1} = \ell^\perp$ \\
\hline
Projection onto a line $\ell$ & $1$ and $0$ & $E_1 = \ell$; $E_0 = \ker P = \ell^\perp$ \\
\hline
Shear $T(x,y)=(x+ky,\,y)$ & $1$ (double) & $E_1 = $ the $x$-axis \\
\hline
\end{tabularx}
\end{center}

\begin{propriete}[Eigenvalues and Invertibility]
Let $T$ have matrix $\mathbf{M}$.
\begin{enumerate}
    \item $0$ is an eigenvalue of $T$ $\iff$ $\det\mathbf{M} = 0$ $\iff$ $T$ is not invertible.
    \item If $T$ is invertible and $\lambda$ is an eigenvalue of $T$, then $\lambda \neq 0$ and $\dfrac{1}{\lambda}$ is an eigenvalue of $T^{-1}$, with the same eigenvectors.
    \item The eigenvalue $\lambda = 1$ occurs exactly when $T$ has invariant points other than $O$; the eigenspace $E_1$ is precisely the set of invariant points.
\end{enumerate}
\end{propriete}

\textbf{Proof of (2).}\par
If $T(\vec{u}) = \lambda\vec{u}$ with $\vec{u}\neq\vec{0}$ and $T$ invertible, then $\lambda\neq 0$ (otherwise $\vec{u}\in\ker T$) and applying $T^{-1}$: $\vec{u} = \lambda\,T^{-1}(\vec{u})$, so $T^{-1}(\vec{u}) = \frac{1}{\lambda}\vec{u}$. \hfill$\square$

\begin{savaistu}[Why "Eigen"?]
The prefix \emph{eigen-} is German for "own" or "characteristic": eigenvectors are the transformation's \emph{own} directions. Eigenvalues are everywhere in modern science and technology: the stability of bridges and electricity pylons, the vibration modes of musical instruments, quantum mechanics (energy levels are eigenvalues), and the PageRank algorithm that orders web search results all rest on eigenvalue computations.
\end{savaistu}

\begin{aretenir}[Key Points]
\begin{itemize}
    \item $\lambda$ is an eigenvalue $\iff$ $T(\vec{u}) = \lambda\vec{u}$ for some $\vec{u}\neq\vec{0}$ $\iff$ $\det(\mathbf{M}-\lambda\mathbf{I}) = 0$.
    \item For a $2\times2$ matrix: $\lambda^2 - (\operatorname{tr}\mathbf{M})\lambda + \det\mathbf{M} = 0$; check with $\lambda_1+\lambda_2 = \operatorname{tr}\mathbf{M}$, $\lambda_1\lambda_2 = \det\mathbf{M}$.
    \item The eigenspace $E_\lambda = \ker(T-\lambda I)$ always contains $\vec{0}$; eigenvectors are its \emph{non-zero} elements.
    \item $E_1$ is the set of invariant points; $E_0 = \ker T$.
    \item A rotation (other than $0$ or $\pi$) has no real eigenvalue; a reflection always has $1$ and $-1$.
\end{itemize}
\end{aretenir}

\begin{exercice}[Practice Exercise 2]
\begin{enumerate}
    \item Let $T(x,y) = (2x + y,\; x + 2y)$. Find $\ker T$, $\operatorname{Im} T$, and state whether $T$ is invertible.
    \item Find the invariant points of $T(x,y) = (5x - 4y,\; 2x - y)$.
    \item Let $T(x,y) = (3x - y,\; -x + 3y)$. Find the image of the line $\ell : 2x - y = 1$.
    \item Find the eigenvalues and eigenspaces of $T(x,y) = (2x + 3y,\; 3x + 2y)$. Interpret geometrically.
    \item Show that if $\lambda$ is an eigenvalue of $T$, then $\lambda^2$ is an eigenvalue of $T\circ T$.
\end{enumerate}
\end{exercice}

\begin{corrige}[Exercise 2]
\begin{enumerate}
    \item $\det\mathbf{M} = 4 - 1 = 3 \neq 0$, so $\ker T = \{(0,0)\}$, $\operatorname{Im} T = \mathbb{R}^2$, and $T$ is invertible with $T^{-1}(x,y) = \tfrac13(2x - y,\; -x + 2y)$.
    \item $T(x,y) = (x,y) \iff 4x - 4y = 0$ and $2x - 2y = 0 \iff x = y$: the invariant points form the line $y = x$.
    \item Parametrise $\ell$: $(t,\; 2t-1)$. Then $T(t, 2t-1) = (3t - 2t + 1,\; -t + 6t - 3) = (t+1,\; 5t-3)$. With $x = t+1$, $t = x-1$, so $y = 5(x-1) - 3 = 5x - 8$. The image is the line $5x - y = 8$.
    \item $\operatorname{tr}\mathbf{M} = 4$, $\det\mathbf{M} = 4 - 9 = -5$; characteristic equation $\lambda^2 - 4\lambda - 5 = 0$, giving $\lambda_1 = 5$, $\lambda_2 = -1$.
        $E_5$: $-3x + 3y = 0 \iff y = x$, so $E_5 = \operatorname{span}\{(1,1)\}$.
        $E_{-1}$: $3x + 3y = 0 \iff y = -x$, so $E_{-1} = \operatorname{span}\{(1,-1)\}$.
        Geometrically, $T$ stretches by $5$ along $y=x$ and reflects along $y=-x$: it is the reflection in $y=x$ composed with a stretch of factor $5$.
    \item If $T(\vec{u}) = \lambda\vec{u}$ with $\vec{u}\neq\vec{0}$, then $(T\circ T)(\vec{u}) = T(\lambda\vec{u}) = \lambda T(\vec{u}) = \lambda^2\vec{u}$, so $\lambda^2$ is an eigenvalue of $T\circ T$ with the same eigenvector. \hfill$\square$
\end{enumerate}
\end{corrige}

\coursec{Activities of Integration, Assessment and Remediation}

\begin{experience}[Guided Study — The Printer's Transformation]
A printing workshop in Yaound\'e uses a machine that maps every point $(x,y)$ of a design to $(x', y')$ with
\[
T(x,y) = (x + 0.5y,\; 0.5y).
\]
\begin{enumerate}
    \item Show that $T$ is linear and write its matrix.
    \item Compute $\det\mathbf{M}_T$. Is the machine's transformation invertible? What does this mean practically?
    \item Find the invariant points of $T$.
    \item The design contains the line $y = 2$. Find its image. Are horizontal lines preserved?
    \item Find the eigenvalues and eigenspaces of $T$. Interpret the factor by which lengths along each eigendirection are multiplied.
\end{enumerate}
\textbf{Guided answers:} (1) homogeneous linear components; $\mathbf{M} = \begin{pmatrix}1&0.5\\0&0.5\end{pmatrix}$. (2) $\det\mathbf{M} = 0.5 \neq 0$: invertible, the original design can be recovered. (3) $T(x,y)=(x,y) \iff y = 0$: the $x$-axis is fixed. (4) $T(t,2) = (t+1, 1)$: the image is the line $y = 1$; horizontal lines stay horizontal but are shifted and compressed vertically. (5) Characteristic equation $(1-\lambda)(0.5-\lambda) = 0$: $\lambda_1 = 1$ with $E_1 = \operatorname{span}\{(1,0)\}$ (horizontal lengths unchanged), $\lambda_2 = 0.5$ with $E_2$: $0.5x + 0.5y = 0$, i.e. $E_2 = \operatorname{span}\{(1,-1)\}$ (lengths along this direction halved).
\end{experience}

\begin{exercice}[Assessment — GCE Style]
\begin{enumerate}
    \item The transformation $T$ of the plane has matrix $\mathbf{M} = \begin{pmatrix} 3 & 1 \\ 2 & 2 \end{pmatrix}$.
        \begin{enumerate}
            \item Find the image of the point $(2, -1)$.
            \item Find the eigenvalues of $T$ and an eigenvector for each.
            \item Deduce the equations of two lines through the origin that are invariant under $T$.
            \item Find $T^{-1}$ and the image of the line $x + y = 4$ under $T$.
        \end{enumerate}
    \item Let $S$ be the reflection in the line $y = \sqrt{3}\,x$.
        \begin{enumerate}
            \item Using $\alpha = 60^\circ$ in the reflection matrix, write the matrix of $S$.
            \item Verify that $S\circ S = I$.
            \item State the eigenvalues of $S$ and describe the eigenspaces geometrically.
        \end{enumerate}
    \item $T : \mathbb{R}^3 \to \mathbb{R}^3$ is defined by $T(x,y,z) = (x - z,\; y + z,\; x + y)$.
        \begin{enumerate}
            \item Find $\ker T$ and $\operatorname{Im} T$; verify the rank--nullity theorem.
            \item Is $T$ invertible? Justify.
        \end{enumerate}
\end{enumerate}
\end{exercice}

\begin{corrige}[Assessment]
\begin{enumerate}
    \item \textbf{(a)} $T(2,-1) = (6-1,\; 4-2) = (5,2)$.
        \textbf{(b)} $\operatorname{tr}\mathbf{M} = 5$, $\det\mathbf{M} = 4$; $\lambda^2 - 5\lambda + 4 = 0$ gives $\lambda_1 = 1$, $\lambda_2 = 4$.
        $E_1$: $2x + y = 0$, eigenvector $(1,-2)$. $E_4$: $-x + y = 0$, eigenvector $(1,1)$.
        \textbf{(c)} The invariant lines through $O$ are the eigenspaces: $y = -2x$ (points fixed) and $y = x$ (points stretched by $4$).
        \textbf{(d)} $\mathbf{M}^{-1} = \tfrac14\begin{pmatrix}2&-1\\-2&3\end{pmatrix}$, so $T^{-1}(x,y) = \tfrac14(2x-y,\; -2x+3y)$.
        Line $x+y=4$: parametrise $(t, 4-t)$; $T(t,4-t) = (3t + 4 - t,\; 2t + 8 - 2t) = (2t+4,\; 8)$. The image is the horizontal line $y = 8$.
    \item \textbf{(a)} $\cos120^\circ = -\tfrac12$, $\sin120^\circ = \tfrac{\sqrt3}{2}$: $\mathbf{M}_S = \begin{pmatrix} -\frac12 & \frac{\sqrt3}{2} \\[2pt] \frac{\sqrt3}{2} & \frac12 \end{pmatrix}$.
        \textbf{(b)} $\mathbf{M}_S^2 = \begin{pmatrix} \frac14+\frac34 & -\frac{\sqrt3}{4}+\frac{\sqrt3}{4} \\[2pt] -\frac{\sqrt3}{4}+\frac{\sqrt3}{4} & \frac34+\frac14 \end{pmatrix} = \begin{pmatrix}1&0\\0&1\end{pmatrix} = \mathbf{I}$.
        \textbf{(c)} Eigenvalues $1$ and $-1$; $E_1$ is the mirror line $y = \sqrt3\,x$ (fixed points), $E_{-1}$ is the perpendicular line through $O$ (vectors reversed).
    \item \textbf{(a)} Kernel: $x = z$, $y = -z$, $x + y = 0$; the third equation gives $z - z = 0$, automatically satisfied, so $\ker T = \operatorname{span}\{(1,-1,1)\}$, nullity $1$.
        Columns: $(1,0,1)$, $(0,1,1)$, $(-1,1,0)$; note $(-1,1,0) = (0,1,1) - (1,0,1)$, so $\operatorname{Im} T = \operatorname{span}\{(1,0,1), (0,1,1)\}$, a plane, rank $2$. Check: $1 + 2 = 3$. \hfill$\square$
        \textbf{(b)} No: $\ker T \neq \{\vec{0}\}$ (equivalently $\det\mathbf{M} = 0$), so $T$ is not invertible.
\end{enumerate}
\end{corrige}

\begin{attention}[Remediation — Final Checklist Before the Examination]
\begin{itemize}
    \item Always test $T(\vec{0})$ first when linearity is in doubt.
    \item Columns of the matrix = images of the basis vectors, \emph{in order}.
    \item $\mathbf{M}_{S\circ T} = \mathbf{M}_S\mathbf{M}_T$: rightmost transformation is applied first.
    \item Kernel: solve $\mathbf{M}\vec{u} = \vec{0}$. Image: span of the columns. Never confuse the two.
    \item Invariant points: solve $T(\vec{u}) = \vec{u}$, i.e. $(\mathbf{M}-\mathbf{I})\vec{u} = \vec{0}$ — this is the eigenspace $E_1$.
    \item Eigenvectors must be \textbf{non-zero}; always verify your answer by computing $T(\vec{u})$ and comparing with $\lambda\vec{u}$.
    \item Check eigenvalues with the trace and the determinant: $\lambda_1+\lambda_2 = \operatorname{tr}\mathbf{M}$, $\lambda_1\lambda_2 = \det\mathbf{M}$.
\end{itemize}
\end{attention}

\coursec{Kernel and Image of a Linear Transformation}

In the previous section we discovered that a linear transformation $T$ is completely determined by its matrix $A$, and that the matrix acts on vectors by multiplication: $T(\vec{u}) = A\vec{u}$. A natural question now arises: \emph{what does a linear transformation do to the space as a whole?} Does it squash the plane onto a line? Does it rotate everything faithfully? Does it send some non-zero vectors to zero?

Two sets answer these questions completely. The \cle{kernel} tells us \emph{what is destroyed} (sent to the zero vector), and the \cle{image} tells us \emph{what is produced} (the set of all outputs). Together they give a precise ``identity card'' of any linear transformation, and they are the key to deciding whether $T$ is injective, surjective, or invertible --- which we shall need in the next section.

\courssub{The Kernel of a Linear Transformation}

\textbf{An observation to start with.}\par
Consider the orthogonal projection $P$ of the plane onto the $x$-axis. Every point $(x,y)$ is sent to $(x,0)$. Now look at the points on the $y$-axis: $(0,1)$, $(0,-3)$, $(0,7)$ are \emph{all} sent to $(0,0)$. The whole $y$-axis is ``crushed'' onto the origin. This collection of vectors that $P$ annihilates is exactly what we call the kernel of $P$.

\begin{definition}[Kernel of a linear transformation]
Let $T : E \to F$ be a linear transformation. The \cle{kernel} of $T$, written $\ker T$, is the set of all vectors of $E$ that are sent to the zero vector of $F$:
\[ \ker T = \{\, \vec{u} \in E \;:\; T(\vec{u}) = \vec{0} \,\}. \]
If $T$ has matrix $A$, then $\ker T$ is the set of solutions of the homogeneous system $AX = 0$.
\end{definition}

Note carefully: $\ker T$ lives in the \emph{domain} $E$. It is never empty, because linearity forces $T(\vec{0}) = \vec{0}$, so $\vec{0} \in \ker T$ always.

\begin{propriete}[The kernel is a subspace]
$\ker T$ is a vector subspace of $E$: it contains $\vec{0}$, and it is closed under addition and scalar multiplication.
\end{propriete}

\textbf{Proof (examinable).} We verify the three subspace conditions using linearity of $T$.
\begin{enumerate}
\item $T(\vec{0}) = \vec{0}$, so $\vec{0} \in \ker T$.
\item If $\vec{u}, \vec{v} \in \ker T$, then $T(\vec{u} + \vec{v}) = T(\vec{u}) + T(\vec{v}) = \vec{0} + \vec{0} = \vec{0}$, so $\vec{u} + \vec{v} \in \ker T$.
\item If $\vec{u} \in \ker T$ and $\lambda \in \mathbb{R}$, then $T(\lambda \vec{u}) = \lambda T(\vec{u}) = \lambda \vec{0} = \vec{0}$, so $\lambda \vec{u} \in \ker T$. \hfill$\square$
\end{enumerate}

This is excellent news: instead of listing infinitely many vectors, we can describe $\ker T$ by a \emph{basis}. In the plane, the kernel is either the single point $\{\vec{0}\}$, a line through the origin, or the whole plane.

\begin{methode}[Finding the kernel in practice]
To find $\ker T$ when $T$ has matrix $A$:
\begin{enumerate}
\item Write the homogeneous system $AX = 0$, where $X = \begin{pmatrix} x \\ y \end{pmatrix}$.
\item Solve it (by elimination or by inspection).
\item Describe the solution set geometrically: the origin alone, a line through the origin (give a direction vector), or a plane through the origin in 3D.
\end{enumerate}
\end{methode}

\begin{exoresolu}[Kernel of a plane transformation]
Let $T$ be the linear transformation of the plane with matrix $A = \begin{pmatrix} 1 & 2 \\ 2 & 4 \end{pmatrix}$. Determine $\ker T$ and describe it geometrically.
\tcblower
\textbf{Solution.} We solve $AX = 0$:
\[ \begin{cases} x + 2y = 0 \\ 2x + 4y = 0 \end{cases} \]
The second equation is just twice the first, so the only condition is $x = -2y$. Setting $y = t$ (a free parameter), the solutions are
\[ \begin{pmatrix} x \\ y \end{pmatrix} = t \begin{pmatrix} -2 \\ 1 \end{pmatrix}, \qquad t \in \mathbb{R}. \]
Hence $\ker T$ is the \textbf{line through the origin} spanned by $\begin{pmatrix} -2 \\ 1 \end{pmatrix}$, i.e. the line $x + 2y = 0$. Every vector on this line is crushed to $\vec{0}$ by $T$.
\end{exoresolu}

\courssub{The Image of a Linear Transformation}

\textbf{An observation to start with.}\par
Return to the projection $P$ onto the $x$-axis. Whatever point $(x,y)$ of the plane we choose, its image $(x,0)$ lies \emph{on} the $x$-axis; and every point of the $x$-axis is reached (it is its own image). So the totality of outputs of $P$ is exactly the $x$-axis. This set of all outputs is the image of $P$.

\begin{definition}[Image of a linear transformation]
Let $T : E \to F$ be a linear transformation. The \cle{image} of $T$, written $\operatorname{Im} T$, is the set of all outputs of $T$:
\[ \operatorname{Im} T = \{\, T(\vec{u}) \;:\; \vec{u} \in E \,\} = \{\, \vec{w} \in F \;:\; \vec{w} = T(\vec{u}) \text{ for some } \vec{u} \in E \,\}. \]
\end{definition}

Note carefully: $\operatorname{Im} T$ lives in the \emph{codomain} $F$.

\begin{propriete}[The image is a subspace]
$\operatorname{Im} T$ is a vector subspace of $F$.
\end{propriete}

\textbf{Proof (examinable).} (i) $\vec{0} = T(\vec{0}) \in \operatorname{Im} T$. (ii) If $\vec{w}_1 = T(\vec{u}_1)$ and $\vec{w}_2 = T(\vec{u}_2)$, then $\vec{w}_1 + \vec{w}_2 = T(\vec{u}_1 + \vec{u}_2) \in \operatorname{Im} T$. (iii) If $\vec{w} = T(\vec{u})$ and $\lambda \in \mathbb{R}$, then $\lambda \vec{w} = T(\lambda \vec{u}) \in \operatorname{Im} T$. \hfill$\square$

How do we compute it concretely? Write $\vec{u} = x\vec{e}_1 + y\vec{e}_2$ in the standard basis. By linearity,
\[ T(\vec{u}) = x\,T(\vec{e}_1) + y\,T(\vec{e}_2). \]
But $T(\vec{e}_1)$ and $T(\vec{e}_2)$ are precisely the \textbf{columns} of the matrix $A$! So every output is a linear combination of the columns.

\begin{methode}[Finding the image in practice]
The image of $T$ is the \emph{span of the columns} of its matrix $A$:
\[ \operatorname{Im} T = \operatorname{span}\{\, T(\vec{e}_1),\ T(\vec{e}_2) \,\} = \operatorname{span}\{\text{columns of } A\}. \]
\begin{enumerate}
\item Read off the columns of $A$.
\item If they are not parallel, they span the whole plane: $\operatorname{Im} T = \mathbb{R}^2$.
\item If they are parallel (one is a multiple of the other), $\operatorname{Im} T$ is the line through the origin in that common direction.
\end{enumerate}
\end{methode}

\begin{exoresolu}[Image of a plane transformation]
For the same matrix $A = \begin{pmatrix} 1 & 2 \\ 2 & 4 \end{pmatrix}$, find $\operatorname{Im} T$.
\tcblower
\textbf{Solution.} The columns are $\begin{pmatrix} 1 \\ 2 \end{pmatrix}$ and $\begin{pmatrix} 2 \\ 4 \end{pmatrix} = 2\begin{pmatrix} 1 \\ 2 \end{pmatrix}$. They are parallel, so
\[ \operatorname{Im} T = \operatorname{span}\left\{ \begin{pmatrix} 1 \\ 2 \end{pmatrix} \right\}, \]
the line through the origin of slope $2$ (equation $y = 2x$). Notice the beautiful picture: $T$ crushes the kernel line $x + 2y = 0$ to a point, and the whole plane collapses onto the image line $y = 2x$.
\end{exoresolu}

\begin{attention}[Do not confuse kernel and image]
The \textbf{kernel} is a set of \emph{inputs}: it lives in the \textbf{domain} $E$. The \textbf{image} is a set of \emph{outputs}: it lives in the \textbf{codomain} $F$. A frequent GCE error is to solve $AX = 0$ and then declare the answer to be ``the image''. Remember: \emph{kernel = what is killed; image = what is reached}.
\end{attention}

\begin{popfigure}
\begin{tikzpicture}[scale=1.1]
% Left oval: E
\draw[thick, popblue] (0,0) ellipse (1.6 and 2.0);
\node[popblue, font=\large\bfseries] at (0,2.35) {$E$};
% kernel inside E
\draw[thick, poporange, fill=poporange!20] (0,-0.4) ellipse (0.55 and 0.9);
\node[popdark, font=\small] at (0,-0.4) {$\ker T$};
\fill[popdark] (0,-0.4) circle (1.5pt);
% Right oval: F
\draw[thick, popgreen] (6,0) ellipse (1.6 and 2.0);
\node[popgreen!60!black, font=\large\bfseries] at (6,2.35) {$F$};
% zero in F
\fill[popdark] (6,-0.4) circle (1.5pt);
\node[popdark, font=\small, right] at (6.1,-0.55) {$\vec{0}$};
% Im T inside F
\draw[thick, poppurple, fill=poppurple!15] (6,0.55) ellipse (0.7 and 1.0);
\node[popdark, font=\small] at (6,0.75) {$\operatorname{Im} T$};
% arrows
\draw[-latex, thick, poporange] (0.6,-0.4) to[bend left=15] node[above, font=\small]{$T$} (5.2,-0.4);
\draw[-latex, thick, poppurple] (0.9,0.9) to[bend left=25] (4.9,1.1);
\node[font=\small, popmuted] at (3,1.9) {$T(E) = \operatorname{Im} T$};
\end{tikzpicture}
\legende{The kernel (part of the domain $E$) maps to $\vec{0}$; the image is the part of the codomain $F$ actually reached by $T$.}
\end{popfigure}

\courssub{Geometric Interpretation: the Case of a Projection}

Projection onto the $x$-axis is the perfect mental model. Its matrix is
\[ A = \begin{pmatrix} 1 & 0 \\ 0 & 0 \end{pmatrix}. \]
Solving $AX = 0$ gives $x = 0$: the kernel is the $y$-axis, the direction that is ``projected away''. The columns of $A$ are $\begin{pmatrix} 1 \\ 0 \end{pmatrix}$ and $\begin{pmatrix} 0 \\ 0 \end{pmatrix}$, so the image is the $x$-axis, the line \emph{onto which} we project.

\begin{popfigure}
\begin{tikzpicture}[scale=1.0]
% grid points and their projections
\draw[thin, gray!40] (-0.5,-2.6) grid (4.5,2.6);
% kernel: y-axis highlighted
\draw[line width=2.2pt, poporange] (0,-2.6) -- (0,2.6);
\node[poporange, font=\small\bfseries, align=center] at (-0.05,2.85) {kernel: $y$-axis};
% image: x-axis highlighted
\draw[line width=2.2pt, popblue] (-0.5,0) -- (4.5,0);
\node[popblue, font=\small\bfseries, align=center] at (3.6,-0.35) {image: $x$-axis};
% example points
\fill[popdark] (3,2) circle (2pt);
\node[popdark, font=\small, above right] at (3,2) {$(3,2)$};
\draw[-latex, thick, poppurple, dashed] (3,2) -- (3,0.06);
\fill[popblue] (3,0) circle (2pt);
\node[popblue, font=\small, below right] at (3.05,0.05) {$(3,0)$};
\fill[poporange] (0,1.5) circle (2pt);
\draw[-latex, thick, poporange] (0,1.5) -- (0,0.08);
\node[popdark, font=\small, left] at (-0.05,1.5) {$(0,1.5) \to \vec{0}$};
\end{tikzpicture}
\legende{Orthogonal projection onto the $x$-axis: the kernel ($y$-axis, orange) is crushed to $\vec{0}$; the image ($x$-axis, blue) receives every output.}
\end{popfigure}

The same phenomenon occurs for any singular $2 \times 2$ matrix: the whole plane is collapsed onto a single line (the image), and all vectors pointing in one direction (the kernel) are annihilated.

\begin{popfigure}
\begin{tikzpicture}[scale=1.0]
% left: original plane with kernel direction
\draw[thin, gray!40] (-2.2,-2.2) grid (2.2,2.2);
\draw[->, gray] (-2.3,0) -- (2.4,0);
\draw[->, gray] (0,-2.3) -- (0,2.4);
% kernel line x+2y=0 i.e. direction (-2,1)
\draw[line width=2pt, poporange] (2,-1) -- (-2,1);
\node[poporange, font=\small\bfseries] at (-1.7,1.45) {$\ker T$};
\fill[popdark] (1,1) circle (1.8pt);
\fill[popdark] (-1,2) circle (1.8pt);
\node[font=\small, popmuted] at (0,-2.6) {plane $E$};
% arrow
\draw[-latex, very thick, popdark] (2.7,0) -- (4.1,0) node[midway, above, font=\small]{$T$};
% right: collapsed line y=2x
\begin{scope}[xshift=6.6cm]
\draw[thin, gray!40] (-2.2,-2.2) grid (2.2,2.2);
\draw[->, gray] (-2.3,0) -- (2.4,0);
\draw[->, gray] (0,-2.3) -- (0,2.4);
\draw[line width=2pt, popblue] (-1,-2) -- (1,2);
\node[popblue, font=\small\bfseries] at (1.35,1.7) {$\operatorname{Im} T$};
\fill[popdark] (0,0) circle (2pt);
\node[font=\small, popmuted] at (0,-2.6) {plane $F$};
\end{scope}
\end{tikzpicture}
\legende{The singular transformation $A = \begin{pmatrix} 1 & 2 \\ 2 & 4 \end{pmatrix}$ collapses the whole plane onto the line $\operatorname{Im} T : y = 2x$; the kernel line $x + 2y = 0$ is sent to the origin.}
\end{popfigure}

\courssub{Kernel, Injectivity and Determinants}

When is a transformation \emph{injective} (one-to-one), i.e. when do two different inputs never give the same output? Suppose $T(\vec{u}) = T(\vec{v})$. By linearity, $T(\vec{u} - \vec{v}) = \vec{0}$, so $\vec{u} - \vec{v} \in \ker T$. If the kernel contains only $\vec{0}$, we are forced to conclude $\vec{u} = \vec{v}$. This proves the fundamental criterion:

\begin{propriete}[Injectivity criterion]
A linear transformation $T$ is \cle{injective} if and only if $\ker T = \{\vec{0}\}$.
\end{propriete}

This is a remarkable simplification: to test injectivity, we do not need to compare all pairs of vectors --- we only solve one homogeneous system.

For a \emph{square} matrix, solving $AX = 0$ connects directly with the determinant. The system $AX = 0$ has a unique solution (namely $X = 0$) exactly when $A$ is non-singular:

\begin{propriete}[Kernel and determinant]
Let $T$ be a linear transformation of the plane (or of space) with square matrix $A$. Then
\[ \ker T = \{\vec{0}\} \iff \det A \neq 0. \]
In that case $T$ is injective, and for plane transformations it is also \cle{surjective}: $\operatorname{Im} T$ is the whole plane, because the two columns of $A$ are then non-parallel and span $\mathbb{R}^2$.
\end{propriete}

Conversely, $\det A = 0$ means the columns are parallel: the image collapses to a line (or a point), and a whole line of vectors is sent to $\vec{0}$.

\begin{exemplebox}[Quick test]
$A = \begin{pmatrix} 1 & 2 \\ 2 & 4 \end{pmatrix}$: $\det A = 4 - 4 = 0$, so $\ker T$ is a line and $T$ is neither injective nor surjective --- confirming our computations above.
$B = \begin{pmatrix} 2 & 1 \\ 1 & 3 \end{pmatrix}$: $\det B = 6 - 1 = 5 \neq 0$, so $\ker T = \{\vec{0}\}$ and $\operatorname{Im} T = \mathbb{R}^2$: $T$ is a bijection of the plane.
\end{exemplebox}

\begin{savaistu}[GCE extra: nothing is lost, everything is transformed]
Here is a powerful intuition for checking your answers (stated informally, as the GCE requires): \emph{what is lost in the kernel reappears as the size of the image}. In the plane: if the kernel is a line (dimension $1$), the image is a line (dimension $1$), and $1 + 1 = 2$, the dimension of the plane. If the kernel is just $\{\vec{0}\}$ (dimension $0$), the image is the whole plane (dimension $2$). In 3D, a kernel that is a plane (dimension $2$) forces the image to be a line (dimension $1$): $2 + 1 = 3$. This ``conservation law'' (the rank--nullity theorem, studied at university) lets you spot impossible answers instantly.
\end{savaistu}

\begin{exoresolu}[Kernel and image in space]
Let $T : \mathbb{R}^3 \to \mathbb{R}^3$ have matrix $A = \begin{pmatrix} 1 & 0 & 1 \\ 0 & 1 & 1 \\ 1 & 1 & 2 \end{pmatrix}$. Find $\ker T$ and $\operatorname{Im} T$, and check the ``conservation'' of dimensions.
\tcblower
\textbf{Solution.} Solve $AX = 0$:
\[ \begin{cases} x + z = 0 \\ y + z = 0 \\ x + y + 2z = 0 \end{cases} \]
The third equation is the sum of the first two, so it adds nothing. With $z = t$: $x = -t$, $y = -t$. Hence
\[ \ker T = \operatorname{span}\left\{ \begin{pmatrix} -1 \\ -1 \\ 1 \end{pmatrix} \right\}, \quad \text{a line through the origin (dimension } 1\text{)}. \]
The columns of $A$ are $\vec{c}_1 = \begin{pmatrix} 1 \\ 0 \\ 1 \end{pmatrix}$, $\vec{c}_2 = \begin{pmatrix} 0 \\ 1 \\ 1 \end{pmatrix}$, $\vec{c}_3 = \vec{c}_1 + \vec{c}_2$. Since $\vec{c}_1$ and $\vec{c}_2$ are not parallel,
\[ \operatorname{Im} T = \operatorname{span}\{ \vec{c}_1, \vec{c}_2 \}, \quad \text{a plane through the origin (dimension } 2\text{)}. \]
Check: $1 + 2 = 3 = \dim \mathbb{R}^3$. \checkmark
\end{exoresolu}

\begin{exercice}[Kernel and image]
For each matrix, find $\ker T$ and $\operatorname{Im} T$, state whether $T$ is injective and/or surjective, and verify the dimension check.
\begin{enumerate}
\item $A = \begin{pmatrix} 3 & -6 \\ -1 & 2 \end{pmatrix}$
\item $B = \begin{pmatrix} 1 & 1 \\ 0 & 1 \end{pmatrix}$
\item $C = \begin{pmatrix} 0 & 0 \\ 0 & 0 \end{pmatrix}$ (the zero transformation)
\end{enumerate}
\end{exercice}

\begin{corrige}[Exercises on Kernel and Image]
\textbf{1.} $\det A = 6 - 6 = 0$. Kernel: $3x - 6y = 0 \Rightarrow x = 2y$, so $\ker T = \operatorname{span}\{(2,1)\}$, a line. Columns are parallel: $\operatorname{Im} T = \operatorname{span}\{(3,-1)\}$, a line. Neither injective nor surjective. Check: $1 + 1 = 2$. \checkmark

\textbf{2.} $\det B = 1 \neq 0$. Kernel: $x + y = 0$ and $y = 0$ force $x = y = 0$, so $\ker T = \{\vec{0}\}$. Columns $(1,0)$ and $(1,1)$ are non-parallel: $\operatorname{Im} T = \mathbb{R}^2$. $T$ is a bijection. Check: $0 + 2 = 2$. \checkmark

\textbf{3.} Every vector satisfies $CX = 0$: $\ker T = \mathbb{R}^2$ (the whole plane). The only output is $\vec{0}$: $\operatorname{Im} T = \{\vec{0}\}$. Check: $2 + 0 = 2$. \checkmark
\end{corrige}

\begin{aretenir}[Key points]
\begin{itemize}
\item $\ker T = \{\vec{u} : T(\vec{u}) = \vec{0}\}$: a subspace of the \textbf{domain}, found by solving $AX = 0$.
\item $\operatorname{Im} T = \{T(\vec{u})\}$: a subspace of the \textbf{codomain}, spanned by the \textbf{columns} of $A$.
\item $T$ injective $\iff \ker T = \{\vec{0}\} \iff \det A \neq 0$ (square matrices).
\item $T$ surjective (plane) $\iff \operatorname{Im} T = \mathbb{R}^2 \iff \det A \neq 0$.
\item For a projection: the kernel is the direction projected away; the image is the line (or plane) onto which we project.
\item Dimension check: $\dim \ker T + \dim \operatorname{Im} T = \dim E$.
\end{itemize}
\end{aretenir}

We now know exactly when a linear transformation is a bijection: when its kernel is trivial. In the next section we shall exploit this to \emph{undo} the transformation --- defining its inverse --- and to hunt for the points that the transformation leaves fixed, the invariant points.

\coursec{Inverse of a Linear Transformation and Invariant Points}

In the previous section we saw that the kernel and the image tell us whether a linear transformation $T$ ``collapses'' information: if $\ker T = \{\vec{0}\}$ the transformation is injective, and if $\operatorname{Im} T$ is the whole target space it is surjective. When both hold, every output vector comes from exactly one input vector — and this means we can \emph{undo} the transformation. That is the subject of this section: the \cle{inverse} of a linear transformation, and the closely related question of which points are left \emph{unchanged} by $T$, the \cle{invariant points}.

\courssub{Bijective Linear Transformations and the Inverse}

\begin{experience}[Can we go backwards?]
Consider the plane transformation $T$ with matrix $A = \begin{pmatrix} 2 & 1 \\ 1 & 1 \end{pmatrix}$. A point $P(1,2)$ is sent to $T(P) = (4,3)$. Now suppose someone tells you only the image $Q(4,3)$ and asks: \emph{which point was sent here?} You must solve
\[
\begin{cases} 2x + y = 4 \\ x + y = 3 \end{cases}
\]
which gives the unique solution $x = 1$, $y = 2$. Because $\det A = 2(1) - 1(1) = 1 \neq 0$, this system \emph{always} has exactly one solution, whatever $Q$ is. The transformation can be reversed. Contrast this with $A = \begin{pmatrix} 1 & 2 \\ 2 & 4 \end{pmatrix}$, where $\det A = 0$: the point $(2,1)$ and the point $(0,0)$ both have image $(2,4)$, so there is no way to know where you came from. Reversibility is impossible.
\end{experience}

\begin{definition}[Bijective linear transformation and its inverse]
A linear transformation $T : \mathbb{R}^n \to \mathbb{R}^n$ is said to be \cle{bijective} (or \cle{invertible}) when it is both \textbf{injective} ($\ker T = \{\vec{0}\}$) and \textbf{surjective} ($\operatorname{Im} T = \mathbb{R}^n$). In that case every vector $\vec{v}$ has a unique pre-image, and the \cle{inverse transformation} $T^{-1}$ is defined by
\[
T^{-1}(\vec{v}) = \vec{u} \iff T(\vec{u}) = \vec{v}, \qquad \text{so that} \qquad T^{-1}\bigl(T(\vec{u})\bigr) = \vec{u} \quad \text{and} \quad T\bigl(T^{-1}(\vec{v})\bigr) = \vec{v}.
\]
For plane transformations ($n = 2$), injectivity alone already implies surjectivity (by the rank–nullity theorem), so $T$ is invertible as soon as $\ker T = \{\vec{0}\}$.
\end{definition}

\begin{propriete}[Invertibility criterion and matrix of the inverse]
Let $T$ be a linear transformation of the plane with matrix $A$. Then
\[
T \text{ is invertible} \iff \det A \neq 0,
\]
and in that case $T^{-1}$ is also a \textbf{linear} transformation, whose matrix is the inverse matrix:
\[
A^{-1} = \frac{1}{\det A}\begin{pmatrix} d & -b \\ -c & a \end{pmatrix} \qquad \text{for} \qquad A = \begin{pmatrix} a & b \\ c & d \end{pmatrix}.
\]
The proof of the criterion (via $\ker T$ and the rank of $A$) is examinable.
\end{propriete}

\textbf{Justification.} $T(\vec{u}) = A\vec{u} = \vec{0}$ has a non-zero solution if and only if the columns of $A$ are linearly dependent, i.e. $\det A = 0$. Hence $\ker T = \{\vec{0}\} \iff \det A \neq 0$. By the rank–nullity theorem, $\dim(\ker T) + \dim(\operatorname{Im} T) = 2$, so $\ker T = \{\vec{0}\}$ forces $\dim(\operatorname{Im} T) = 2$, i.e. $\operatorname{Im} T = \mathbb{R}^2$ (surjectivity). Linearity of $T^{-1}$ follows because for any $\vec{v}_1, \vec{v}_2$, the pre-image of $\vec{v}_1 + \vec{v}_2$ under $T$ is the sum of the pre-images (since $T$ is linear). Finally $AA^{-1} = A^{-1}A = I$, so $A^{-1}(A\vec{u}) = \vec{u}$, confirming that $A^{-1}$ undoes $A$. $\blacksquare$

\begin{exemplebox}[Computing an inverse transformation]
Let $T$ have matrix $A = \begin{pmatrix} 3 & 1 \\ 2 & 1 \end{pmatrix}$. Then $\det A = 3 - 2 = 1 \neq 0$, so $T$ is invertible and
\[
A^{-1} = \frac{1}{1}\begin{pmatrix} 1 & -1 \\ -2 & 3 \end{pmatrix} = \begin{pmatrix} 1 & -1 \\ -2 & 3 \end{pmatrix}.
\]
Check: $AA^{-1} = \begin{pmatrix} 3-2 & -3+3 \\ 2-2 & -2+3 \end{pmatrix} = \begin{pmatrix} 1 & 0 \\ 0 & 1 \end{pmatrix} = I$. \checkmark
\end{exemplebox}

\begin{popfigure}
\begin{center}
\begin{tikzpicture}[scale=0.9, >=stealth]
% Left plane
\draw[popline] (-4.2,-1.6) rectangle (-0.6,1.6);
\node[popdark] at (-2.4,1.9) {\small Domain};
\fill[popblue] (-2.4,-0.2) circle (2pt) node[below left, popblue] {\small $\vec{u}$};
% Right plane
\draw[popline] (0.6,-1.6) rectangle (4.2,1.6);
\node[popdark] at (2.4,1.9) {\small Codomain};
\fill[poporange] (2.4,0.3) circle (2pt) node[below right, poporange] {\small $\vec{v}=T(\vec{u})$};
% Arrows
\draw[->, very thick, popblue] (-2.0,0.6) .. controls (-0.8,1.3) and (0.8,1.3) .. (2.0,0.9);
\node[popblue] at (0,1.5) {\small $T$};
\draw[->, very thick, poporange] (2.0,-0.4) .. controls (0.8,-1.1) and (-0.8,-1.1) .. (-2.0,-0.8);
\node[poporange] at (0,-1.35) {\small $T^{-1}$};
\end{tikzpicture}
\end{center}
\legende{An invertible transformation $T$ and its inverse $T^{-1}$: each undoes the other.}
\end{popfigure}

\courssub{Using the Inverse to Find Pre-images}

\begin{methode}[Finding the pre-image of a point]
To find the point $P$ such that $T(P) = Q$:
\begin{enumerate}
\item Check that $\det A \neq 0$ (otherwise $Q$ may have no pre-image, or infinitely many).
\item Compute $A^{-1}$ using the determinant formula.
\item The pre-image is $P = A^{-1}Q$.
\item \textbf{Verify} by computing $T(P)$ and checking it equals $Q$.
\end{enumerate}
\end{methode}

\begin{exoresolu}[Finding a pre-image]
Let $T$ be the plane transformation with matrix $A = \begin{pmatrix} 2 & 1 \\ 1 & 1 \end{pmatrix}$. Find the point $P$ whose image is $Q(5,3)$.
\tcblower
$\det A = 2 - 1 = 1 \neq 0$, so $T$ is invertible with
$A^{-1} = \begin{pmatrix} 1 & -1 \\ -1 & 2 \end{pmatrix}$.
Then
\[
P = A^{-1}Q = \begin{pmatrix} 1 & -1 \\ -1 & 2 \end{pmatrix}\begin{pmatrix} 5 \\ 3 \end{pmatrix} = \begin{pmatrix} 5-3 \\ -5+6 \end{pmatrix} = \begin{pmatrix} 2 \\ 1 \end{pmatrix}.
\]
Verification: $T(2,1) = (2\cdot 2 + 1,\ 2 + 1) = (5,3) = Q$. \checkmark\ Hence $P(2,1)$.
\end{exoresolu}

\begin{attention}[GCE extra: check your inverse before using it]
In the exam, one arithmetic slip in $A^{-1}$ ruins every following answer. Spend ten seconds verifying $AA^{-1} = I$ \emph{before} using the inverse. Also remember: $(AB)^{-1} = B^{-1}A^{-1}$ — the order reverses, exactly like putting on and taking off shoes and socks.
\end{attention}

\begin{savaistu}[Determinant of the inverse]
If $T$ is invertible, $\det(A^{-1}) = \dfrac{1}{\det A}$. This follows from $\det(AA^{-1}) = \det I = 1$. Geometrically, if $T$ scales areas by factor $|\det A|$, then $T^{-1}$ scales them by the reciprocal factor.
\end{savaistu}

\courssub{Invariant (Fixed) Points}

\begin{experience}[Which points do not move?]
Stand in front of a mirror. Your nose moves to a new position in the reflection — but a point \emph{on the mirror surface itself} stays exactly where it is. Reflections, projections and rotations all behave differently: some points are fixed, others are not. Finding the fixed points of a transformation reveals its geometric nature.
\end{experience}

\begin{definition}[Invariant point]
A point $M$ with position vector $\vec{x}$ is an \cle{invariant point} (or \cle{fixed point}) of $T$ when
\[
T(M) = M, \qquad \text{i.e.} \qquad A\vec{x} = \vec{x}.
\]
Equivalently, $\vec{x}$ satisfies $(A - I)\vec{x} = \vec{0}$, where $I$ is the $2 \times 2$ identity matrix.
\end{definition}

\begin{methode}[Finding and interpreting invariant points]
\begin{enumerate}
\item Write the condition $A\vec{x} = \vec{x}$, i.e. $(A - I)\vec{x} = \vec{0}$.
\item Solve the resulting homogeneous system in $x$ and $y$.
\item Interpret geometrically:
\begin{itemize}
\item only $(0,0)$: the unique fixed point is the origin (e.g. a rotation about $O$);
\item a line of solutions: every point of that line is fixed (e.g. the mirror line of a reflection, the line onto which we project);
\item the whole plane: $T$ is the identity transformation.
\end{itemize}
\end{enumerate}
\end{methode}

\begin{attention}[The classic trap]
Fixed points satisfy $(A - I)\vec{x} = \vec{0}$, \textbf{not} $A\vec{x} = \vec{0}$. Solving $A\vec{x} = \vec{0}$ gives the \emph{kernel} (points sent to the origin), which is a completely different question. Subtract the identity first!
\end{attention}

\begin{exoresolu}[Fixed points of a reflection]
$T$ is the reflection in the line $y = x$, with matrix $A = \begin{pmatrix} 0 & 1 \\ 1 & 0 \end{pmatrix}$. Find its invariant points.
\tcblower
$A - I = \begin{pmatrix} -1 & 1 \\ 1 & -1 \end{pmatrix}$, so $(A-I)\vec{x} = \vec{0}$ gives $-x + y = 0$, i.e. $y = x$. Every point of the mirror line is fixed — exactly as expected: points on the mirror do not move under reflection. Any point off the line, e.g. $(3,1) \mapsto (1,3)$, is genuinely moved.
\end{exoresolu}

\begin{exoresolu}[Fixed points of a projection and of a rotation]
(a) Find the fixed points of the projection onto the $x$-axis, $A = \begin{pmatrix} 1 & 0 \\ 0 & 0 \end{pmatrix}$.\quad (b) Same question for the rotation of angle $\theta \neq 0$ about $O$.
\tcblower
(a) $A - I = \begin{pmatrix} 0 & 0 \\ 0 & -1 \end{pmatrix}$ gives $-y = 0$, i.e. $y = 0$: the fixed points are exactly the $x$-axis, the line onto which we project. Points already on the target line stay put.

(b) $A - I = \begin{pmatrix} \cos\theta - 1 & -\sin\theta \\ \sin\theta & \cos\theta - 1 \end{pmatrix}$ has determinant $(\cos\theta - 1)^2 + \sin^2\theta = 2 - 2\cos\theta \neq 0$ for $\theta \neq 0$, so the only solution is $(0,0)$: a rotation fixes only the origin.
\end{exoresolu}

\begin{popfigure}
\begin{center}
\begin{tikzpicture}[scale=1.1, >=stealth]
\draw[->, popline] (-2.2,0) -- (2.6,0) node[right] {\small $x$};
\draw[->, popline] (0,-2.2) -- (0,2.6) node[above] {\small $y$};
\draw[very thick, popblue] (-2,-2) -- (2.4,2.4) node[above right, popblue] {\small $y=x$: fixed line};
\fill[popblue] (1,1) circle (2pt) node[below right, popblue] {\small $M_1$ fixed};
\fill[poporange] (1.8,0.4) circle (2pt) node[below, poporange] {\small $M$};
\fill[poporange] (0.4,1.8) circle (2pt) node[left, poporange] {\small $T(M)$};
\draw[->, dashed, poporange] (1.8,0.4) -- (0.4,1.8);
\draw[dotted, popmuted] (1.8,0.4) -- (1.1,1.1) -- (0.4,1.8);
\end{tikzpicture}
\end{center}
\legende{Reflection in $y = x$: every point of the mirror line is invariant; $M$ swaps with its image.}
\end{popfigure}

\begin{popfigure}
\begin{center}
\begin{tikzpicture}[scale=1.1, >=stealth]
\draw[->, popline] (-0.5,0) -- (4.6,0) node[right] {\small $x$};
\draw[->, popline] (0,-0.5) -- (0,2.8) node[above] {\small $y$};
\draw[very thick, popgreen] (0,0) -- (4.3,0);
\node[popgreen] at (2.2,-0.35) {\small line of fixed points ($y=0$)};
\fill[popgreen] (3,0) circle (2pt) node[above right, popgreen] {\small $M_1$ fixed};
\fill[poporange] (1.5,2.2) circle (2pt) node[above, poporange] {\small $M$};
\fill[poporange] (1.5,0) circle (2pt) node[below, poporange] {\small $T(M)$};
\draw[->, dashed, poporange] (1.5,2.2) -- (1.5,0.12);
\end{tikzpicture}
\end{center}
\legende{Projection onto the $x$-axis: the axis is the set of invariant points; $M$ slides down to its image.}
\end{popfigure}

\courssub{Invariant Lines: a First Look at Eigen-directions}

A line through the origin is \cle{invariant} (as a set) when every vector on it is sent to a vector \emph{on the same line}: $T(\vec{u}) = \lambda\vec{u}$ for some scalar $\lambda$. The direction of such a line is called an \cle{eigen-direction}.

\begin{exemplebox}[Invariant lines of a reflection]
For the reflection in $y = x$: the mirror line itself is invariant with $\lambda = 1$ (pointwise fixed), and the perpendicular line $y = -x$ is also invariant, since $(t,-t) \mapsto (-t,t) = -(t,-t)$, with $\lambda = -1$. These two special directions — vectors simply \emph{stretched} by $T$ — are exactly the \textbf{eigenvectors} we shall study in Section 5, and $\lambda$ the corresponding \textbf{eigenvalues}.
\end{exemplebox}

\begin{experience}[Discovering invariant lines]
Take the transformation $T$ with matrix $A = \begin{pmatrix} 2 & 1 \\ 1 & 2 \end{pmatrix}$. Test the lines $y = x$ and $y = -x$:
\begin{itemize}
\item On $y = x$: vector $(t,t)$ gives $T(t,t) = (3t,3t) = 3(t,t)$ — invariant with $\lambda = 3$.
\item On $y = -x$: vector $(t,-t)$ gives $T(t,-t) = (t,-t) = 1(t,-t)$ — invariant with $\lambda = 1$.
\end{itemize}
These are the only invariant lines through the origin. This is a preview of the eigenvalue problem: find $\lambda$ and non-zero $\vec{u}$ such that $A\vec{u} = \lambda\vec{u}$.
\end{experience}

\begin{aretenir}[Key points]
\begin{itemize}
\item $T$ is invertible $\iff$ $\det A \neq 0$; then $T^{-1}$ is linear with matrix $A^{-1} = \dfrac{1}{\det A}\begin{pmatrix} d & -b \\ -c & a \end{pmatrix}$.
\item Pre-image of $Q$: compute $P = A^{-1}Q$, then verify $T(P) = Q$.
\item Fixed points: solve $(A - I)\vec{x} = \vec{0}$ — never $A\vec{x} = \vec{0}$.
\item Geometry: mirror line (reflection), target line/plane (projection) are fixed; a rotation fixes only $O$.
\item Invariant lines through $O$ give eigen-directions: $T(\vec{u}) = \lambda\vec{u}$.
\end{itemize}
\end{aretenir}

\begin{exercice}[Practice]
\begin{enumerate}
\item Let $T$ have matrix $A = \begin{pmatrix} 4 & 3 \\ 1 & 1 \end{pmatrix}$. Show $T$ is invertible, find $A^{-1}$, and determine the pre-image of $Q(7,2)$.
\item Find the invariant points of the transformation with matrix $\begin{pmatrix} 2 & 1 \\ 1 & 2 \end{pmatrix}$. What do you conclude?
\item Show that the lines $y = x$ and $y = -x$ are invariant under the transformation of question 2, and find the factor $\lambda$ on each line.
\end{enumerate}
\end{exercice}

\begin{corrige}[Exercise 1]
\textbf{1.} $\det A = 4 - 3 = 1 \neq 0$, so $T$ is invertible; $A^{-1} = \begin{pmatrix} 1 & -3 \\ -1 & 4 \end{pmatrix}$ (check: $AA^{-1} = I$). Pre-image: $P = A^{-1}Q = (7 - 6,\ -7 + 8) = (1,1)$. Verify: $T(1,1) = (7,2)$. \checkmark

\textbf{2.} $A - I = \begin{pmatrix} 1 & 1 \\ 1 & 1 \end{pmatrix}$ gives $x + y = 0$: the fixed points form the line $y = -x$.

\textbf{3.} On $y = x$: $(t,t) \mapsto (3t, 3t) = 3(t,t)$, invariant with $\lambda = 3$. On $y = -x$: $(t,-t) \mapsto (t,-t)$, invariant with $\lambda = 1$ — which is why this line is pointwise fixed.
\end{corrige}

\coursec{Images of Lines and Planes}

In the previous section we studied the inverse of a linear transformation and the points left invariant by it. We now ask a geometric question of a different kind: instead of tracking a single point, we track a whole \emph{set} of points. If $T$ is a linear transformation and $L$ is a line, what does the set of images $\{T(P) : P \in L\}$ look like? If $\Pi$ is a plane, what is its image? The answer is remarkably clean, and it is one of the most frequently examined ideas in this chapter: linear transformations send lines to lines (or points) and planes to planes (or lines or points). Understanding \emph{why}, and being able to compute the image quickly, is the goal of this section.

\courssub{Image of a Line Through the Origin}

Let us start with the simplest case. Take a line through the origin in the plane (or in space), say the line $L$ spanned by a non-zero vector $\vec{v}$. Every point of $L$ has position vector $t\vec{v}$ for some $t \in \mathbb{R}$. Apply a linear transformation $T$:
\[
T(t\vec{v}) = t\,T(\vec{v}).
\]
As $t$ runs through $\mathbb{R}$, the image $t\,T(\vec{v})$ runs through the line spanned by $T(\vec{v})$ — unless $T(\vec{v}) = \vec{0}$, in which case every image is the origin. This single line of algebra, which uses nothing but the linearity of $T$, proves the following fundamental result.

\begin{propriete}[Image of a line or a plane]
Let $T$ be a linear transformation.
\begin{itemize}
\item The image of a \cle{line} under $T$ is a \textbf{line or a point}.
\item The image of a \cle{plane} under $T$ is a \textbf{plane, a line or a point}.
\end{itemize}
More precisely, the line through the origin spanned by $\vec{v}$ maps onto the line spanned by $T(\vec{v})$; it collapses to the origin exactly when $\vec{v} \in \ker T$.
\end{propriete}

\begin{attention}[A frequent misconception]
The image of a line that does \textbf{not} pass through the origin is \emph{not} required to pass through the origin. Only \textbf{subspaces} (lines and planes through $O$) are guaranteed to map onto subspaces through $O$, because $T(\vec{0}) = \vec{0}$. A general line maps to a general line, which may or may not contain $O$. Keep this in mind whenever you read off the image line.
\end{attention}

\courssub{Image of a General Line: the Parametric Method}

A general line $L$ is described by one point $A$ on it and a direction vector $\vec{v}$:
\[
\vec{r} = \vec{a} + t\vec{v}, \qquad t \in \mathbb{R},
\]
where $\vec{a}$ is the position vector of $A$. Applying $T$ and using linearity:
\[
T(\vec{r}) = T(\vec{a} + t\vec{v}) = T(\vec{a}) + t\,T(\vec{v}).
\]
This is again the parametric equation of a line: it passes through the point $T(A)$ and has direction vector $T(\vec{v})$. If $T(\vec{v}) = \vec{0}$, the image is the single point $T(A)$.

\begin{methode}[Finding the image of a line]
To find the image of the line $\vec{r} = \vec{a} + t\vec{v}$ under the linear transformation with matrix $A$:
\begin{enumerate}
\item Compute the image of the point: $A\vec{a}$, giving a point $A'$ of the image line.
\item Compute the image of the direction vector: $A\vec{v}$, giving the direction $\vec{v}\,'$ of the image line.
\item Write the image line as $\vec{r}\,' = A\vec{a} + t\,(A\vec{v})$, and convert to Cartesian form if required.
\item Check the degenerate case: if $A\vec{v} = \vec{0}$, the image is the single point $A'$.
\end{enumerate}
\end{methode}

\begin{exoresolu}[Image of a line under a shear]
The transformation $T$ has matrix $M = \begin{pmatrix} 1 & 1 \\ 0 & 1 \end{pmatrix}$ (a shear parallel to the $x$-axis). Find the image of the line $L: y = x + 1$.
\tcblower
\textbf{Step 1: parametrise $L$.} A point on $L$ is $A(0,1)$ and a direction vector is $\vec{v} = \begin{pmatrix} 1 \\ 1 \end{pmatrix}$, so $\vec{r} = \begin{pmatrix} 0 \\ 1 \end{pmatrix} + t\begin{pmatrix} 1 \\ 1 \end{pmatrix}$.

\textbf{Step 2: transform point and direction.}
\[
M\vec{a} = \begin{pmatrix} 1 & 1 \\ 0 & 1 \end{pmatrix}\begin{pmatrix} 0 \\ 1 \end{pmatrix} = \begin{pmatrix} 1 \\ 1 \end{pmatrix}, \qquad
M\vec{v} = \begin{pmatrix} 1 & 1 \\ 0 & 1 \end{pmatrix}\begin{pmatrix} 1 \\ 1 \end{pmatrix} = \begin{pmatrix} 2 \\ 1 \end{pmatrix}.
\]

\textbf{Step 3: write the image line.} $\vec{r}\,' = \begin{pmatrix} 1 \\ 1 \end{pmatrix} + t\begin{pmatrix} 2 \\ 1 \end{pmatrix}$, i.e. $x = 1 + 2t$, $y = 1 + t$. Eliminating $t$: $t = y - 1$, so $x = 1 + 2(y-1) = 2y - 1$.

\textbf{Conclusion:} the image is the line $x = 2y - 1$, that is $y = \dfrac{x+1}{2}$. Note that it does \emph{not} pass through the origin — exactly as warned above.
\end{exoresolu}

\begin{popfigure}
\begin{center}
\begin{tikzpicture}[scale=1.6, >=stealth]
\draw[->, gray] (-1.4,0) -- (3.4,0) node[below left] {$x$};
\draw[->, gray] (0,-0.6) -- (0,3.2) node[below left] {$y$};
\draw[popblue, thick] (-1.2,-0.2) -- (2.2,3.2) node[pos=0.85, above left, popblue] {$L: y=x+1$};
\draw[poporange, thick] (-1.4,-0.2) -- (3.4,2.2) node[pos=0.9, below right, poporange] {$L': x=2y-1$};
\fill[popblue] (0,1) circle (1.3pt) node[above left, popblue] {$A(0,1)$};
\fill[poporange] (1,1) circle (1.3pt) node[below right, poporange] {$A'(1,1)$};
\draw[->, popblue, thick] (0,1) -- (1,2) node[midway, above left, popblue] {$\vec{v}$};
\draw[->, poporange, thick] (1,1) -- (3,2) node[midway, below right, poporange] {$\vec{v}\,'$};
\draw[->, popmuted, dashed] (0.15,1.15) to[bend left=25] (0.85,1.12);
\node[popmuted] at (0.5,1.55) {$T$};
\end{tikzpicture}
\end{center}
\legende{The shear of matrix $\begin{pmatrix} 1 & 1 \\ 0 & 1 \end{pmatrix}$ maps $L$ to $L'$; the point $A$ and the direction vector $\vec{v}$ are tracked.}
\end{popfigure}

\begin{exoresolu}[Image of a line under a rotation]
$T$ is the rotation about $O$ through $90^{\circ}$ anticlockwise, with matrix $R = \begin{pmatrix} 0 & -1 \\ 1 & 0 \end{pmatrix}$. Find the image of the line $L: 2x + y = 4$.
\tcblower
Parametrise: $A(2,0)$ lies on $L$, direction $\vec{v} = \begin{pmatrix} 1 \\ -2 \end{pmatrix}$ (since $2(1) + (-2) = 0$).
\[
R\vec{a} = \begin{pmatrix} 0 & -1 \\ 1 & 0 \end{pmatrix}\begin{pmatrix} 2 \\ 0 \end{pmatrix} = \begin{pmatrix} 0 \\ 2 \end{pmatrix}, \qquad
R\vec{v} = \begin{pmatrix} 0 & -1 \\ 1 & 0 \end{pmatrix}\begin{pmatrix} 1 \\ -2 \end{pmatrix} = \begin{pmatrix} 2 \\ 1 \end{pmatrix}.
\]
Image line: $\vec{r}\,' = \begin{pmatrix} 0 \\ 2 \end{pmatrix} + t\begin{pmatrix} 2 \\ 1 \end{pmatrix}$, i.e. $x = 2t$, $y = 2 + t$, giving $y = 2 + \dfrac{x}{2}$, or $x - 2y + 4 = 0$. Geometrically this is exactly the original line rotated through $90^{\circ}$, as expected.
\end{exoresolu}

\courssub{Image of a Plane in Space}

The same idea works in three dimensions. A plane $\Pi$ through the origin is spanned by two non-parallel vectors $\vec{u}$ and $\vec{v}$: its points are $s\vec{u} + t\vec{v}$. Applying $T$:
\[
T(s\vec{u} + t\vec{v}) = s\,T(\vec{u}) + t\,T(\vec{v}),
\]
so the image is the set of all linear combinations of $T(\vec{u})$ and $T(\vec{v})$:
\begin{itemize}
\item if $T(\vec{u})$ and $T(\vec{v})$ are \textbf{not parallel}, the image is the \textbf{plane} spanned by them;
\item if they are parallel (and not both zero), the image collapses to a \textbf{line};
\item if both are zero, the image is the single point $O$.
\end{itemize}
For a general plane with equation $\vec{r} = \vec{a} + s\vec{u} + t\vec{v}$, the image is $\vec{r}\,' = T(\vec{a}) + s\,T(\vec{u}) + t\,T(\vec{v})$.

\begin{methode}[Finding the image of a plane]
\begin{enumerate}
\item Identify one point $A$ of the plane and two spanning (direction) vectors $\vec{u}, \vec{v}$.
\item Compute $T(\vec{a})$, $T(\vec{u})$, $T(\vec{v})$ using the matrix of $T$.
\item If $T(\vec{u})$ and $T(\vec{v})$ are not parallel, the image is the plane through $T(A)$ spanned by them; find its Cartesian equation using the normal vector $T(\vec{u}) \times T(\vec{v})$.
\item If they are parallel, the image is a line; if both vanish, a point.
\end{enumerate}
\end{methode}

\begin{exoresolu}[Image of a plane in space]
$T$ has matrix $M = \begin{pmatrix} 1 & 0 & 1 \\ 0 & 1 & 1 \\ 0 & 0 & 0 \end{pmatrix}$. Find the image of the plane $\Pi: z = 0$.
\tcblower
$\Pi$ passes through $O$ and is spanned by $\vec{u} = \begin{pmatrix} 1 \\ 0 \\ 0 \end{pmatrix}$ and $\vec{v} = \begin{pmatrix} 0 \\ 1 \\ 0 \end{pmatrix}$.
\[
M\vec{u} = \begin{pmatrix} 1 \\ 0 \\ 0 \end{pmatrix}, \qquad M\vec{v} = \begin{pmatrix} 0 \\ 1 \\ 0 \end{pmatrix}.
\]
These are not parallel, so the image is the plane spanned by them: again the plane $z = 0$. Note that $T$ is singular ($\det M = 0$): the whole of space is projected onto the plane $z = 0$ (since the third column is $(1,1,0)^T$ and the third row is zero). This particular plane $\Pi$ is invariant as a set; in fact, every point of $\Pi$ is fixed because $M(x,y,0)^T = (x,y,0)^T$.
\end{exoresolu}

\begin{popfigure}
\begin{center}
\begin{tikzpicture}[scale=1.1, >=stealth]
\draw[fill=popblueL, draw=popblue, thick] (0,0) -- (3,0.6) -- (4,2.2) -- (1,1.6) -- cycle;
\node[popblue] at (2,1.1) {$\Pi$};
\draw[->, popblue] (1.2,0.7) -- (2.2,0.9) node[midway, below] {$\vec{u}$};
\draw[->, popblue] (1.2,0.7) -- (1.7,1.5) node[midway, left] {$\vec{v}$};
\draw[->, popmuted, thick] (4.4,1.2) -- (5.4,1.2) node[midway, above] {$T$};
\draw[fill=popgreenL, draw=popgreen, thick] (5.8,0.2) -- (8.6,0.2) -- (9.4,1.8) -- (6.6,1.8) -- cycle;
\node[popgreen] at (7.6,1.0) {$\Pi' = T(\Pi)$};
\draw[->, popgreen] (6.6,0.9) -- (7.9,0.9) node[midway, below] {$T(\vec{u})$};
\draw[->, popgreen] (6.6,0.9) -- (7.0,1.6) node[midway, left] {$T(\vec{v})$};
\end{tikzpicture}
\end{center}
\legende{A plane $\Pi$ spanned by $\vec{u}, \vec{v}$ maps to the plane $\Pi'$ spanned by $T(\vec{u}), T(\vec{v})$.}
\end{popfigure}

\courssub{Singular Transformations: Collapse of Lines, Grids and Planes}

When $T$ is \cle{singular} ($\det M = 0$), its image has dimension less than the ambient space, and whole families of lines and planes collapse. In the plane, a singular transformation of rank $1$ sends \emph{every} line onto the same fixed line (or a point on it): the entire plane is flattened onto $\operatorname{Im} T$.

\begin{exoresolu}[A singular transformation]
$T$ has matrix $M = \begin{pmatrix} 1 & 2 \\ 2 & 4 \end{pmatrix}$. Find the image of the line $L: y = x + 1$, and describe the image of the whole plane.
\tcblower
$\det M = 4 - 4 = 0$, so $T$ is singular. Every image vector has the form
\[
M\begin{pmatrix} x \\ y \end{pmatrix} = \begin{pmatrix} x + 2y \\ 2x + 4y \end{pmatrix} = (x+2y)\begin{pmatrix} 1 \\ 2 \end{pmatrix},
\]
so the whole plane maps onto the line $y = 2x$ (the line spanned by $(1,2)$).

For $L$: take $A(0,1)$, $\vec{v} = \begin{pmatrix} 1 \\ 1 \end{pmatrix}$. Then $M\vec{a} = \begin{pmatrix} 2 \\ 4 \end{pmatrix}$ and $M\vec{v} = \begin{pmatrix} 3 \\ 6 \end{pmatrix} = 3\begin{pmatrix} 1 \\ 2 \end{pmatrix}$. The image is the line through $(2,4)$ in direction $(1,2)$, i.e. $y = 2x$: the line $L$ maps \emph{onto the same line} into which the whole plane collapses. Note $M\vec{v} \neq \vec{0}$, so $L$ does not collapse to a point — but a line with direction $\begin{pmatrix} 2 \\ -1 \end{pmatrix}$ would, since $\begin{pmatrix} 2 \\ -1 \end{pmatrix} \in \ker M$.
\end{exoresolu}

\begin{popfigure}
\begin{center}
\begin{tikzpicture}[scale=0.85, >=stealth]
\draw[step=1, popline, very thin] (-0.5,-0.5) grid (3.5,3.5);
\draw[->, gray] (-0.6,0) -- (3.7,0);
\draw[->, gray] (0,-0.6) -- (0,3.7);
\draw[popblue, thick] (-0.3,-0.3) -- (3.3,3.3) node[pos=0.9, above left, popblue] {$y=x$};
\fill[popblue] (1,1) circle (1.5pt);
\fill[popblue] (2,2) circle (1.5pt);
\node at (1.5,-1.0) {Original plane};
\draw[->, popmuted, thick] (4.0,1.5) -- (5.0,1.5) node[midway, above] {$T$};
\draw[->, gray] (5.4,0) -- (9.4,0);
\draw[->, gray] (6,-0.6) -- (6,3.7);
\draw[poporange, very thick] (5.7,-0.3) -- (9.1,3.1) node[pos=0.85, below right, poporange] {$y=2x$};
\fill[poporange] (7,2) circle (1.5pt);
\fill[poporange] (7.5,3) circle (1.5pt);
\node[poporange] at (7.5,-1.0) {Grid collapsed onto a line};
\end{tikzpicture}
\end{center}
\legende{The singular matrix $\begin{pmatrix} 1 & 2 \\ 2 & 4 \end{pmatrix}$ collapses the whole plane (grid and line) onto the single line $y = 2x$.}
\end{popfigure}

In space, a singular transformation can collapse a plane onto a line: for instance with $M = \begin{pmatrix} 1 & 0 & 0 \\ 0 & 0 & 0 \\ 0 & 0 & 0 \end{pmatrix}$, the plane $z = 0$ spanned by $(1,0,0)$ and $(0,1,0)$ maps onto the $x$-axis, because $M(1,0,0)^T = (1,0,0)^T$ but $M(0,1,0)^T = \vec{0}$.

\courssub{What Is Preserved and What Is Not}

Because $T(\vec{a} + t\vec{v}) = T(\vec{a}) + t\,T(\vec{v})$ with the \emph{same} parameter $t$, the position of a point along a line, measured as a fraction of the way along, is unchanged.

\begin{propriete}[Preservation properties of linear transformations]
A linear transformation $T$ preserves:
\begin{itemize}
\item the \cle{origin}: $T(\vec{0}) = \vec{0}$;
\item \cle{parallelism}: parallel lines map to parallel lines (or points), since a common direction $\vec{v}$ maps to the common direction $T(\vec{v})$;
\item \cle{midpoints} and \cle{ratios of division}: if $M$ divides $AB$ in the ratio $\lambda : \mu$, then $T(M)$ divides $T(A)T(B)$ in the same ratio $\lambda : \mu$.
\end{itemize}
However, \textbf{distances} and \textbf{angles} are generally \emph{not} preserved (a shear stretches lengths and distorts angles); they are preserved only by special transformations such as rotations and reflections (orthogonal matrices).
\end{propriete}

The proof of ratio preservation is one line: if $\vec{m} = \dfrac{\mu\vec{a} + \lambda\vec{b}}{\lambda + \mu}$, then by linearity $T(\vec{m}) = \dfrac{\mu\,T(\vec{a}) + \lambda\,T(\vec{b})}{\lambda + \mu}$, which is exactly the point dividing $T(A)T(B)$ in the ratio $\lambda:\mu$.

\begin{savaistu}[Computer Graphics and Linear Transformations]
In 3D computer graphics, every object is represented by a mesh of vertices (points). To move, rotate, scale, or shear an object, the graphics pipeline applies a linear transformation (represented by a $4 \times 4$ matrix in homogeneous coordinates) to every vertex. The fact that linear transformations map straight lines to straight lines means that the edges of the mesh (line segments between vertices) remain straight — we only need to transform the endpoints! This property makes real-time rendering of complex scenes computationally feasible.
\end{savaistu}

\courssub{GCE Extra: Area and Volume Scaling}

\begin{propriete}[Area and volume scaling]
If $T$ has matrix $A$, then areas in the plane are multiplied by $|\det A|$, and volumes in space are multiplied by $|\det A|$. In particular, the unit square maps to a parallelogram of area $|\det A|$. This is a quick check on your work and a favourite one-mark question.
\end{propriete}

For the shear $M = \begin{pmatrix} 1 & 1 \\ 0 & 1 \end{pmatrix}$, $\det M = 1$: areas are preserved. For the singular matrix $\begin{pmatrix} 1 & 2 \\ 2 & 4 \end{pmatrix}$, $\det = 0$: every region collapses to zero area, confirming the flattening we saw above.

\begin{popfigure}
\begin{center}
\begin{tikzpicture}[scale=1.5, >=stealth]
\draw[fill=popblueL, draw=popblue, thick] (0,0) -- (1,0) -- (1,1) -- (0,1) -- cycle;
\node[popblue] at (0.5,0.5) {area $1$};
\draw[->, popmuted, thick] (1.4,0.5) -- (2.2,0.5) node[midway, above] {$T$};
\draw[fill=poporange!25, draw=poporange, thick] (2.6,0) -- (4.1,0) -- (4.6,1.2) -- (3.1,1.2) -- cycle;
\node[poporange] at (3.7,0.6) {area $|\det A|$};
\draw[->, poporange] (2.6,0) -- (4.1,0) node[midway, below] {$T(\vec{e}_1)$};
\draw[->, poporange] (2.6,0) -- (3.1,1.2) node[midway, left] {$T(\vec{e}_2)$};
\end{tikzpicture}
\end{center}
\legende{The unit square maps to the parallelogram spanned by the columns of $A$; its area is $|\det A|$.}
\end{popfigure}

\begin{exercice}[GCE practice]
\begin{enumerate}
\item $T$ has matrix $\begin{pmatrix} 2 & 1 \\ 0 & 3 \end{pmatrix}$. Find the image of the line $x + y = 2$, and state the factor by which areas are scaled.
\item $T$ has matrix $\begin{pmatrix} 1 & -1 \\ 2 & -2 \end{pmatrix}$. Show that the whole plane maps onto a line, find that line, and find the image of the line $y = x$.
\item In space, $T$ has matrix $\begin{pmatrix} 1 & 0 & 0 \\ 0 & 1 & 0 \\ 1 & 1 & 0 \end{pmatrix}$. Find the image of the plane $x + y + z = 1$.
\end{enumerate}
\end{exercice}

\begin{corrige}[GCE practice]
\textbf{1.} Point $A(2,0)$, direction $\vec{v} = (1,-1)^T$. Images: $A\vec{a} = (4,0)^T$, $A\vec{v} = (1,-3)^T$. Image line: $\vec{r} = (4,0)^T + t(1,-3)^T$, i.e. $3x + y = 12$. Areas scale by $|\det| = 6$.

\textbf{2.} Every image is $(x-y)(1,2)^T$, so the plane maps onto the line $y = 2x$. For $y = x$: direction $(1,1)^T$ maps to $\vec{0}$, so the whole line maps to the single point $O$ (indeed $y = x$ is the kernel direction, and the line passes through $O$).

\textbf{3.} The plane has point $A(1,0,0)$ and spanning vectors $\vec{u} = (1,-1,0)^T$, $\vec{v} = (1,0,-1)^T$. Images: $M\vec{a} = (1,0,1)^T$, $M\vec{u} = (1,-1,0)^T$, $M\vec{v} = (1,0,-1)^T$. Normal: $(1,-1,0) \times (1,0,-1) = (1,1,1)$. Image plane: $x + y + z = 1 + 0 + 1 = 2$, i.e. $x + y + z = 2$.
\end{corrige}

\begin{aretenir}[Key points]
\begin{itemize}
\item Image of $\vec{r} = \vec{a} + t\vec{v}$: transform the point and the direction separately; the image is $\vec{r}\,' = T(\vec{a}) + t\,T(\vec{v})$.
\item A line maps to a line or a point; a plane maps to a plane, a line or a point.
\item Lines through $O$ map to lines through $O$; general lines need not.
\item Parallelism, midpoints and division ratios are preserved; distances and angles are not.
\item Areas scale by $|\det A|$; $\det A = 0$ means total collapse of area.
\end{itemize}
\end{aretenir}

\coursec{Eigenvalues and Eigenvectors}

In the previous section we studied the images of lines and planes under a linear transformation $T$. We saw that a generic line through the origin is sent to a \emph{different} line through the origin: the transformation rotates and stretches most directions. But a natural question arises: are there special directions that $T$ leaves \emph{invariant} — directions along which a vector is simply stretched, compressed or reversed, but never rotated away? These special directions, when they exist, are given by the \cle{eigenvectors} of $T$, and the stretch factors are the \cle{eigenvalues}. They are the key to understanding the deep structure of a linear transformation, and they appear every year in the GCE Advanced Level paper.

\courssub{Intuition and Definition}

\begin{experience}[Looking for invariant directions]
Consider the linear transformation of the plane given by
\[ A=\begin{pmatrix} 2 & 0 \\ 0 & \tfrac12 \end{pmatrix}. \]
Take the vector $\vec{u}=\begin{pmatrix}1\\0\end{pmatrix}$. Then $A\vec{u}=\begin{pmatrix}2\\0\end{pmatrix}=2\vec{u}$: the vector keeps its direction and is doubled. Take $\vec{v}=\begin{pmatrix}0\\1\end{pmatrix}$. Then $A\vec{v}=\begin{pmatrix}0\\1/2\end{pmatrix}=\tfrac12\vec{v}$: same direction, halved length. But take $\vec{w}=\begin{pmatrix}1\\1\end{pmatrix}$: $A\vec{w}=\begin{pmatrix}2\\1/2\end{pmatrix}$, which is \emph{not} a multiple of $\vec{w}$ — the direction has changed. The two coordinate axes are invariant directions; the diagonal is not.
\end{experience}

\begin{popfigure}
\begin{center}
\begin{tikzpicture}[scale=1.05]
% original grid
\draw[step=0.5, popline, very thin] (-1.5,-1.5) grid (1.5,1.5);
\draw[->, thick] (-1.7,0) -- (1.7,0);
\draw[->, thick] (0,-1.7) -- (0,1.7);
\draw[popblue, very thick, ->] (0,0) -- (1,0) node[midway, below, text=popblue] {$\vec{u}$};
\draw[poporange, very thick, ->] (0,0) -- (0,1) node[midway, right, text=poporange] {$\vec{v}$};
\draw[poppurple, very thick, ->] (0,0) -- (1,1) node[midway, above left, text=poppurple] {$\vec{w}$};
\node at (0,-2.1) {\small Original grid};
\end{tikzpicture}
\hspace{1cm}
\begin{tikzpicture}[scale=1.05]
% transformed grid: x doubled, y halved
\draw[popline, very thin] (-3,-0.75) -- (3,-0.75) (3,-0.75) -- (3,0.75) (3,0.75) -- (-3,0.75) (-3,0.75) -- (-3,-0.75);
\foreach \x in {-3,-2.5,...,3} \draw[popline, very thin] (\x,-0.75) -- (\x,0.75);
\foreach \y in {-0.75,-0.25,0.25,0.75} \draw[popline, very thin] (-3,\y) -- (3,\y);
\draw[->, thick] (-3.4,0) -- (3.4,0);
\draw[->, thick] (0,-1.1) -- (0,1.1);
\draw[popblue, very thick, ->] (0,0) -- (2,0) node[midway, below, text=popblue] {$2\vec{u}$};
\draw[poporange, very thick, ->] (0,0) -- (0,0.5) node[midway, right, text=poporange] {$\frac12\vec{v}$};
\draw[poppurple, very thick, ->] (0,0) -- (2,0.5) node[midway, above, text=poppurple] {$A\vec{w}$};
\node at (0,-1.6) {\small Image under $A$: horizontal stretch $\times 2$, vertical compression $\times \frac12$};
\end{tikzpicture}
\end{center}
\legende{The two eigen-directions of $A=\begin{pmatrix}2&0\\0&1/2\end{pmatrix}$: the horizontal axis is stretched by $2$, the vertical axis compressed by $\tfrac12$. A generic vector $\vec{w}$ changes direction.}
\end{popfigure}

\begin{definition}[Eigenvalue and eigenvector]
Let $T$ be a linear transformation with matrix $A$. A \textbf{non-zero} vector $\vec{v}$ is called an \cle{eigenvector} of $T$ (or of $A$) if there exists a scalar $\lambda$ such that
\[ T(\vec{v}) = A\vec{v} = \lambda \vec{v}. \]
The scalar $\lambda$ is called the \cle{eigenvalue} corresponding to $\vec{v}$. The set of all solutions of $A\vec{v}=\lambda\vec{v}$ (including $\vec{0}$) is a subspace called the \cle{eigenspace} of $\lambda$.
\end{definition}

Geometrically, an eigenvector spans an \textbf{invariant line through the origin}: every vector on that line stays on that line after transformation. Along it, $T$ acts as a simple scaling:
\begin{itemize}
\item $\lambda > 1$: a \textbf{stretch};
\item $0 < \lambda < 1$: a \textbf{compression};
\item $\lambda < 0$: a stretch combined with a \textbf{reversal} of sense;
\item $\lambda = 1$: the line is fixed \emph{pointwise};
\item $\lambda = 0$: the direction is annihilated — the vector is sent to $\vec{0}$.
\end{itemize}

\begin{popfigure}
\begin{center}
\begin{tikzpicture}[scale=1.6]
\draw[->, thick, gray] (-1.6,0) -- (2.6,0);
\draw[->, thick, gray] (0,-1.4) -- (0,2.2);
% eigenvector v and image
\draw[popblue, very thick, ->] (0,0) -- (1,0.6) node[above left, text=popblue] {$\vec{v}$};
\draw[popblue!60, very thick, ->, dashed] (0,0) -- (2,1.2) node[right, text=popblue] {$T(\vec{v})=\lambda\vec{v}$};
% non-eigenvector w
\draw[poporange, very thick, ->] (0,0) -- (0.4,1) node[above, text=poporange] {$\vec{w}$};
\draw[poporange!60, very thick, ->, dashed] (0,0) -- (-0.9,0.9) node[left, text=poporange] {$T(\vec{w})$};
\draw[popline, thin] (-1.2,-0.72) -- (2.4,1.44);
\node[text=popmuted] at (2.3,1.0) {\small eigen-line};
\end{tikzpicture}
\end{center}
\legende{An eigenvector $\vec{v}$ keeps its line (its image is collinear with it); a non-eigenvector $\vec{w}$ changes direction.}
\end{popfigure}

\begin{attention}[The zero vector is never an eigenvector]
The equation $A\vec{v}=\lambda\vec{v}$ is satisfied by $\vec{v}=\vec{0}$ for \emph{every} $\lambda$, which would make the notion meaningless. By definition, an eigenvector must be \textbf{non-zero}. Saying ``$\vec{0}$ is an eigenvector'' is a classic exam error.
\end{attention}

\courssub{The Characteristic Equation}

How do we find the eigenvalues of a given matrix $A$? Rewrite the defining equation:
\[ A\vec{v}=\lambda\vec{v} \iff A\vec{v}-\lambda\vec{v}=\vec{0} \iff (A-\lambda I)\vec{v}=\vec{0}, \]
where $I$ is the identity matrix. This is a homogeneous system. Since we demand $\vec{v}\neq\vec{0}$, the system must have non-trivial solutions, which happens exactly when the matrix $A-\lambda I$ is \textbf{singular}:
\[ \det(A-\lambda I)=0. \]

\begin{aretenir}[Characteristic equation]
The eigenvalues of $A$ are the solutions of the \cle{characteristic equation}
\[ \det(A-\lambda I)=0. \]
For a $2\times 2$ matrix $A=\begin{pmatrix}a&b\\c&d\end{pmatrix}$ this is the quadratic
\[ \lambda^2-(a+d)\lambda+(ad-bc)=0, \qquad\text{i.e.}\qquad \lambda^2-(\operatorname{tr}A)\lambda+\det A=0. \]
\end{aretenir}

\begin{attention}[Subtract $\lambda$ on the diagonal only]
A very common sign error: students write $A-\lambda I$ incorrectly. Remember that $\lambda I=\begin{pmatrix}\lambda&0\\0&\lambda\end{pmatrix}$, so you subtract $\lambda$ \textbf{only from the diagonal entries}:
\[ A-\lambda I=\begin{pmatrix}a-\lambda & b\\ c & d-\lambda\end{pmatrix}. \]
Subtracting $\lambda$ from every entry, or adding it, destroys the whole computation.
\end{attention}

\begin{methode}[Step 1 — Find the eigenvalues]
\begin{enumerate}
\item Form the matrix $A-\lambda I$ (subtract $\lambda$ on the diagonal only).
\item Compute $\det(A-\lambda I)$ and expand: you obtain the characteristic polynomial.
\item Solve $\det(A-\lambda I)=0$. For $2\times2$, this is a quadratic; for $3\times3$, a cubic (look for an integer root by trial).
\end{enumerate}
\end{methode}

\begin{methode}[Step 2 — Find the eigenvectors (eigenspaces)]
For \emph{each} eigenvalue $\lambda$ found in Step 1:
\begin{enumerate}
\item Substitute $\lambda$ into $A-\lambda I$.
\item Solve the homogeneous system $(A-\lambda I)X=0$. The two equations are always proportional (check this — it confirms your eigenvalue!); keep one.
\item Write the general solution: it is a line through the origin, e.g.\ $\begin{pmatrix}x\\y\end{pmatrix}=t\begin{pmatrix}p\\q\end{pmatrix}$, $t\in\mathbb{R}$. Any non-zero multiple of $\begin{pmatrix}p\\q\end{pmatrix}$ is an eigenvector.
\end{enumerate}
\end{methode}

\courssub{Worked Examples for $2\times2$ Matrices}

\begin{exoresolu}[Two distinct real eigenvalues]
Find the eigenvalues and eigenvectors of $A=\begin{pmatrix} 4 & 1 \\ 2 & 3 \end{pmatrix}$.
\tcblower
\textbf{Step 1.} $A-\lambda I=\begin{pmatrix}4-\lambda & 1\\ 2 & 3-\lambda\end{pmatrix}$, so
\[ \det(A-\lambda I)=(4-\lambda)(3-\lambda)-2=\lambda^2-7\lambda+10=(\lambda-2)(\lambda-5). \]
Eigenvalues: $\lambda_1=2$, $\lambda_2=5$.

\textbf{Step 2.} For $\lambda_1=2$: $(A-2I)X=0$ gives $\begin{cases}2x+y=0\\2x+y=0\end{cases}$, so $y=-2x$: eigenvectors $t\begin{pmatrix}1\\-2\end{pmatrix}$.

For $\lambda_2=5$: $(A-5I)X=0$ gives $-x+y=0$, so $y=x$: eigenvectors $t\begin{pmatrix}1\\1\end{pmatrix}$.

\textbf{Check:} $A\begin{pmatrix}1\\-2\end{pmatrix}=\begin{pmatrix}2\\-4\end{pmatrix}=2\begin{pmatrix}1\\-2\end{pmatrix}$ \checkmark; $A\begin{pmatrix}1\\1\end{pmatrix}=\begin{pmatrix}5\\5\end{pmatrix}=5\begin{pmatrix}1\\1\end{pmatrix}$ \checkmark.
\end{exoresolu}

\begin{exoresolu}[A repeated eigenvalue]
Find the eigenvalues and eigenvectors of $A=\begin{pmatrix} 3 & 1 \\ 0 & 3 \end{pmatrix}$.
\tcblower
$\det(A-\lambda I)=(3-\lambda)^2=0$, so $\lambda=3$ is a \textbf{double} eigenvalue (algebraic multiplicity $2$).

$(A-3I)X=0$ gives $\begin{pmatrix}0&1\\0&0\end{pmatrix}\begin{pmatrix}x\\y\end{pmatrix}=\begin{pmatrix}0\\0\end{pmatrix}$, i.e.\ $y=0$. The eigenspace is only the line $t\begin{pmatrix}1\\0\end{pmatrix}$: there is only \emph{one} invariant direction (geometric multiplicity $1$), even though the eigenvalue is repeated. Such a matrix (a shear-like map) cannot be diagonalised.
\end{exoresolu}

\begin{exoresolu}[No real eigenvalues: a rotation]
Find the eigenvalues of the rotation matrix $A=\begin{pmatrix} 0 & -1 \\ 1 & 0 \end{pmatrix}$ (rotation through $90^\circ$ about $O$).
\tcblower
$\det(A-\lambda I)=\lambda^2+1=0$, which has \textbf{no real solution}. This makes perfect geometric sense: a $90^\circ$ rotation moves every non-zero vector off its own line, so no real invariant direction exists. (Over $\mathbb{C}$ the eigenvalues are $\lambda=\pm i$.)
\end{exoresolu}

\courssub{Links with Kernel, Image and Invariant Points}

\begin{propriete}[Eigenvalue $1$ and eigenvalue $0$]
Let $T$ be a linear transformation with matrix $A$.
\begin{itemize}
\item The eigenvectors with eigenvalue $\lambda=1$ are exactly the non-zero \textbf{invariant (fixed) points} of $T$: $T(\vec{v})=\vec{v}$. The eigenspace for $\lambda=1$ is the set of fixed points of Section 3.
\item The eigenvectors with eigenvalue $\lambda=0$ are exactly the non-zero vectors of the \textbf{kernel}: $T(\vec{v})=\vec{0}$. Hence $\ker T$ is the eigenspace for $\lambda=0$, and $T$ is singular (non-invertible) if and only if $0$ is an eigenvalue.
\end{itemize}
\end{propriete}

This unifies the whole chapter: fixed points, kernel, and invariant lines are all special cases of the single equation $A\vec{v}=\lambda\vec{v}$.

\begin{popfigure}
\begin{center}
\begin{tikzpicture}[scale=1.5]
% unit circle
\draw[popblue, thick] (0,0) circle (1);
\draw[popblue, very thick, ->] (0,0) -- (1,0);
\draw[popblue, very thick, ->] (0,0) -- (0,1);
\draw[->, thick, gray] (-1.4,0) -- (1.4,0);
\draw[->, thick, gray] (0,-1.4) -- (0,1.4);
\node at (0,-1.8) {\small Unit circle, eigen-directions};
\end{tikzpicture}
\hspace{0.6cm}
\begin{tikzpicture}[scale=1.5]
\draw[->, very thick, popdark] (0,0.4) -- (0.7,0.4) node[midway, above] {$A$};
\end{tikzpicture}
\hspace{0.6cm}
\begin{tikzpicture}[scale=1.5]
% ellipse with semi-axes 2 and 0.5
\draw[poporange, thick] (0,0) ellipse (2 and 0.5);
\draw[poporange, very thick, ->] (0,0) -- (2,0) node[right, text=poporange] {$\lambda_1=2$};
\draw[poporange, very thick, ->] (0,0) -- (0,0.5) node[above, text=poporange] {$\lambda_2=\tfrac12$};
\draw[->, thick, gray] (-2.4,0) -- (2.4,0);
\draw[->, thick, gray] (0,-1.4) -- (0,1.4);
\node at (0,-1.8) {\small Image ellipse, axes along eigen-directions};
\end{tikzpicture}
\end{center}
\legende{Summary: a linear transformation maps the unit circle to an ellipse whose principal axes lie along the eigen-directions, with semi-axis lengths $|\lambda_1|$ and $|\lambda_2|$.}
\end{popfigure}

\courssub{Extension to $3\times3$ Matrices}

For a $3\times3$ matrix the characteristic equation $\det(A-\lambda I)=0$ is a \textbf{cubic}. At GCE level the matrices are chosen so that the cubic has simple integer roots: expand the determinant along a row or column with zeros, then look for an integer root $\lambda_0$ by trial ($\pm1,\pm2,\pm3,\dots$), factor out $(\lambda-\lambda_0)$, and solve the remaining quadratic.

\begin{exoresolu}[A $3\times3$ example]
Find the eigenvalues and eigenvectors of $A=\begin{pmatrix} 2 & 1 & 0 \\ 1 & 2 & 0 \\ 0 & 0 & 3 \end{pmatrix}$.
\tcblower
Expanding along the third row:
\[ \det(A-\lambda I)=(3-\lambda)\big[(2-\lambda)^2-1\big]=(3-\lambda)(\lambda^2-4\lambda+3)=(3-\lambda)^2(1-\lambda). \]
Eigenvalues: $\lambda=1$ and $\lambda=3$ (double, algebraic multiplicity $2$).

For $\lambda=1$: $(A-I)X=0$ gives $x+y=0$ and $2z=0$, so eigenvectors $t\begin{pmatrix}1\\-1\\0\end{pmatrix}$ (geometric multiplicity $1$).

For $\lambda=3$: $(A-3I)X=0$ gives $-x+y=0$, $z$ free, so eigenvectors $\begin{pmatrix}x\\x\\z\end{pmatrix}=s\begin{pmatrix}1\\1\\0\end{pmatrix}+t\begin{pmatrix}0\\0\\1\end{pmatrix}$: here the eigenspace is a \emph{plane} through the origin (geometric multiplicity $2$).
\end{exoresolu}

\courssub{GCE Extra: Powers of a Matrix by Diagonalisation}

\textbf{(Beyond the core syllabus — a recurring advanced question type.)} Suppose a $2\times2$ matrix $A$ has two distinct eigenvalues $\lambda_1,\lambda_2$ with eigenvectors $\vec{v}_1,\vec{v}_2$. Form the matrix $P=(\vec{v}_1\ \vec{v}_2)$ whose columns are the eigenvectors. Then
\[ P^{-1}AP=D=\begin{pmatrix}\lambda_1&0\\0&\lambda_2\end{pmatrix}, \qquad\text{so}\qquad A=PDP^{-1}. \]
The power of diagonalisation: $A^2=(PDP^{-1})(PDP^{-1})=PD^2P^{-1}$, and in general
\[ A^n=PD^nP^{-1}=P\begin{pmatrix}\lambda_1^{\,n}&0\\0&\lambda_2^{\,n}\end{pmatrix}P^{-1}, \]
because $P^{-1}P=I$ cancels at every junction. This turns the hard problem of computing $A^n$ into trivial powers of scalars.

\begin{exoresolu}[Computing $A^n$]
Given $A=\begin{pmatrix} 4 & 1 \\ 2 & 3 \end{pmatrix}$, express $A^n$ in closed form.
\tcblower
From the first worked example: $\lambda_1=2$, $\vec{v}_1=\begin{pmatrix}1\\-2\end{pmatrix}$; $\lambda_2=5$, $\vec{v}_2=\begin{pmatrix}1\\1\end{pmatrix}$. Take $P=\begin{pmatrix}1&1\\-2&1\end{pmatrix}$; then $\det P=3$ and $P^{-1}=\dfrac13\begin{pmatrix}1&-1\\2&1\end{pmatrix}$. Hence
\begin{align*}
A^n &= P\begin{pmatrix}2^n&0\\0&5^n\end{pmatrix}P^{-1}
=\frac13\begin{pmatrix}1&1\\-2&1\end{pmatrix}\begin{pmatrix}2^n&0\\0&5^n\end{pmatrix}\begin{pmatrix}1&-1\\2&1\end{pmatrix}\\
&=\frac13\begin{pmatrix}2^n+2\cdot5^n & -2^n+5^n\\ -2^{n+1}+2\cdot5^n & 2^{n+1}+5^n\end{pmatrix}.
\end{align*}
Check for $n=1$: $\frac13\begin{pmatrix}2+10&-2+5\\-4+10&4+5\end{pmatrix}=\begin{pmatrix}4&1\\2&3\end{pmatrix}=A$ \checkmark.
\end{exoresolu}

\begin{attention}[Pitfalls in diagonalisation]
\begin{itemize}
\item The columns of $P$ must be eigenvectors \textbf{in the same order} as the eigenvalues in $D$.
\item Diagonalisation requires \emph{enough} independent eigenvectors: the repeated-eigenvalue shear $\begin{pmatrix}3&1\\0&3\end{pmatrix}$ above cannot be diagonalised.
\item Do not forget the factor $\frac{1}{\det P}$ when computing $P^{-1}$.
\end{itemize}
\end{attention}

\begin{savaistu}[Why ``eigen''?]
The prefix comes from the German word \emph{eigen}, meaning ``own'' or ``characteristic'': eigenvectors are the transformation's ``own'' directions. Eigenvalues are everywhere in modern science: they describe the natural vibration frequencies of bridges and engines, the energy levels of atoms in quantum mechanics, and they are at the heart of the PageRank algorithm that orders web search results.
\end{savaistu}

\begin{exercice}[Practice]
\begin{enumerate}
\item Find the eigenvalues and eigenvectors of $A=\begin{pmatrix}5&-2\\4&-1\end{pmatrix}$. Describe geometrically the action of $A$ along each eigen-direction.
\item Show that $A=\begin{pmatrix}2&2\\1&3\end{pmatrix}$ has eigenvalues $1$ and $4$, and compute $A^n$ by diagonalisation.
\item Find the eigenvalues of $A=\begin{pmatrix}1&1&0\\0&2&0\\0&1&1\end{pmatrix}$ and an eigenvector for each.
\item Explain why a reflection in the line $y=x$ has eigenvalues $1$ and $-1$, and identify the corresponding eigenvectors.
\end{enumerate}
\end{exercice}

\begin{corrige}[Exercise 1]
$\det(A-\lambda I)=(5-\lambda)(-1-\lambda)+8=\lambda^2-4\lambda+3=(\lambda-1)(\lambda-3)$. Eigenvalues $1$ and $3$. For $\lambda=1$: $4x-2y=0$, eigenvectors $t\begin{pmatrix}1\\2\end{pmatrix}$ — a fixed line. For $\lambda=3$: $2x-2y=0$, eigenvectors $t\begin{pmatrix}1\\1\end{pmatrix}$ — a direction stretched by $3$.
\end{corrige}

\begin{corrige}[Exercise 2]
$\det(A-\lambda I)=(2-\lambda)(3-\lambda)-2=\lambda^2-5\lambda+4=(\lambda-1)(\lambda-4)$. Eigenvectors: for $\lambda=1$, $t\begin{pmatrix}2\\-1\end{pmatrix}$; for $\lambda=4$, $t\begin{pmatrix}1\\1\end{pmatrix}$. With $P=\begin{pmatrix}2&1\\-1&1\end{pmatrix}$, $P^{-1}=\frac13\begin{pmatrix}1&-1\\1&2\end{pmatrix}$:
\[ A^n=\frac13\begin{pmatrix}2+4^n & -2+2\cdot4^n\\ -1+4^n & 1+2\cdot4^n\end{pmatrix}. \]
\end{corrige}

\begin{corrige}[Exercise 3]
The matrix is block-triangular; expanding along the first column, $\det(A-\lambda I)=(1-\lambda)^2(2-\lambda)$. Eigenvalues $\lambda=1$ (double, algebraic multiplicity $2$) and $\lambda=2$. For $\lambda=1$: $(A-I)X=0$ gives $y=0$, so $x$ and $z$ are free. The eigenspace is the plane spanned by $\begin{pmatrix}1\\0\\0\end{pmatrix}$ and $\begin{pmatrix}0\\0\\1\end{pmatrix}$ (geometric multiplicity $2$). For $\lambda=2$: $-x+y=0$ and $y-z=0$, eigenvectors $t\begin{pmatrix}1\\1\\1\end{pmatrix}$.
\end{corrige}

\begin{corrige}[Exercise 4]
Reflection in $y=x$ fixes every vector on the line $y=x$, so $\begin{pmatrix}1\\1\end{pmatrix}$ is an eigenvector with $\lambda=1$. It reverses every vector perpendicular to that line, so $\begin{pmatrix}1\\-1\end{pmatrix}$ is an eigenvector with $\lambda=-1$. Its matrix $\begin{pmatrix}0&1\\1&0\end{pmatrix}$ confirms: $\det(A-\lambda I)=\lambda^2-1=0$.
\end{corrige}

\newpage

\coursec{Integration activity}

\courssub{Aims of the activity}

This integration activity mobilises, in a single authentic modelling task, all the skills developed in this chapter:

\begin{itemize}
\item recognise a linear transformation and write its \cle{matrix} in the standard basis;
\item verify the \cle{linearity} axioms $T(\vec{u}+\vec{v})=T(\vec{u})+T(\vec{v})$ and $T(\lambda\vec{u})=\lambda T(\vec{u})$;
\item determine the \cle{kernel} and the \cle{image} of a linear transformation and interpret them geometrically;
\item decide whether a transformation is \cle{bijective} and compute its \cle{inverse};
\item find the \cle{invariant points} of a transformation;
\item compute the \cle{image of a line} under a linear transformation;
\item compute \cle{eigenvalues} and \cle{eigenvectors} and interpret them as privileged directions of stretching;
\item relate the \cle{determinant} of the matrix to the \cle{area scale factor} of the map.
\end{itemize}

Beyond the technical skills, the activity develops modelling, teamwork and scientific communication: each group must present its conclusions on a poster, with a diagram, as a real cartography team would.

\courssub{Situation}

\textbf{The Cartographer of Yaoundé: Designing a Map Projection.}\par

A student in the Department of Geography of the University of Yaoundé I is preparing a simplified digital map of the city centre. She takes the \emph{Poste Centrale} as the origin $O$ of a coordinate system, the unit on each axis being one kilometre, the $x$-axis pointing East and the $y$-axis pointing North. A point of the city is then represented by a vector $\begin{pmatrix} x \\ y \end{pmatrix}$.

To produce the map on a screen, the computer applies to every point of the city the transformation
\[ T(x,y) = (2x+y,\; x+2y). \]
In other words, the point with city coordinates $(x,y)$ is drawn on the screen at the point with \emph{map coordinates} $(x',y')$ where
\[ x' = 2x+y, \qquad y' = x+2y. \]

The student must understand this transformation completely before trusting the map:

\begin{itemize}
\item Is $T$ genuinely a \emph{linear} transformation, so that the whole theory of this chapter applies?
\item Can two different streets of the city collapse onto the same point of the map (which would make the map ambiguous)? Equivalently, is $T$ bijective, and can one recover the city coordinates from the map coordinates?
\item Which points of the city, if any, are drawn exactly where they are (invariant points)?
\item The main avenue is modelled by the straight line of equation $y=x+1$. What does it become on the map?
\item Are there privileged directions of the city which are simply \emph{stretched} by the projection, without being rotated? By what factor?
\item The map will be used to estimate land areas. By what factor are areas multiplied?
\end{itemize}

Working in groups of three or four, you play the role of the cartography team. Answer the tasks below, then prepare a poster summarising your findings, including a diagram showing the eigen-directions and the image of the avenue.

\courssub{Tasks}

\textbf{Task 1 — Matrix and linearity.}
\begin{enumerate}
\item Write the matrix $A$ of $T$ in the standard basis $\left(\begin{pmatrix}1\\0\end{pmatrix},\begin{pmatrix}0\\1\end{pmatrix}\right)$.
\item Prove that $T$ is linear: for all vectors $\vec{u}=(x_1,y_1)$, $\vec{v}=(x_2,y_2)$ and every scalar $\lambda$, show that $T(\vec{u}+\vec{v})=T(\vec{u})+T(\vec{v})$ and $T(\lambda\vec{u})=\lambda T(\vec{u})$.
\end{enumerate}

\textbf{Task 2 — Kernel, image and bijectivity.}
\begin{enumerate}
\item Determine $\ker T$, the set of points of the city sent to the origin of the map.
\item Deduce that $T$ is injective, then determine $\operatorname{Im} T$ and conclude that $T$ is bijective. Interpret: can two distinct points of the city have the same image on the map?
\item Compute $\det A$ and verify that $\det A \neq 0$, confirming bijectivity.
\end{enumerate}

\textbf{Task 3 — The inverse transformation.}
\begin{enumerate}
\item Compute $A^{-1}$.
\item Deduce the formula giving the city coordinates $(x,y)$ from the map coordinates $(x',y')$.
\item A landmark appears on the map at $(x',y')=(4,5)$. Where is it in the city?
\end{enumerate}

\textbf{Task 4 — Invariant points.}

Determine all points $(x,y)$ such that $T(x,y)=(x,y)$. Interpret the result for the map.

\textbf{Task 5 — Image of the avenue.}

The avenue is the line $\mathcal{D}$ of equation $y=x+1$.
\begin{enumerate}
\item Parametrise $\mathcal{D}$ and compute the image of a general point of $\mathcal{D}$.
\item Show that the image is a straight line and give its Cartesian equation in the coordinates $(x',y')$.
\end{enumerate}

\textbf{Task 6 — Eigenvalues and eigenvectors.}
\begin{enumerate}
\item Compute the characteristic polynomial of $A$ and deduce the eigenvalues $\lambda_1$ and $\lambda_2$.
\item For each eigenvalue, find the corresponding eigenspace.
\item Interpret: which directions of the city are simply stretched by the projection, and by what factor? Draw a diagram showing the two eigen-directions.
\end{enumerate}

\textbf{Task 7 — Area scale factor.}

Using $\det A$, state by how much areas are multiplied on the map. A plot of land of area $2\ \mathrm{km^2}$ in the city occupies what area on the map?

\textbf{Task 8 — Poster.}

Prepare a group poster summarising all results, with a TikZ-style diagram of the eigen-directions and of the image of the avenue.

\courssub{Model answer}

\textbf{Task 1.} Since $T(x,y)=(2x+y,\;x+2y)$, the matrix of $T$ in the standard basis is
\[ A=\begin{pmatrix} 2 & 1 \\ 1 & 2 \end{pmatrix}, \qquad \text{so that } \begin{pmatrix} x' \\ y' \end{pmatrix} = A\begin{pmatrix} x \\ y \end{pmatrix}. \]

Linearity. Let $\vec{u}=(x_1,y_1)$, $\vec{v}=(x_2,y_2)$ and $\lambda\in\mathbb{R}$.
\begin{align*}
T(\vec{u}+\vec{v}) &= T(x_1+x_2,\;y_1+y_2)\\
&= \big(2(x_1+x_2)+(y_1+y_2),\;(x_1+x_2)+2(y_1+y_2)\big)\\
&= (2x_1+y_1,\;x_1+2y_1)+(2x_2+y_2,\;x_2+2y_2) = T(\vec{u})+T(\vec{v}).
\end{align*}
Similarly,
\[ T(\lambda\vec{u}) = (2\lambda x_1+\lambda y_1,\;\lambda x_1+2\lambda y_1) = \lambda(2x_1+y_1,\;x_1+2y_1)=\lambda T(\vec{u}). \]
Hence $T$ is a linear transformation. (One could also quote the theorem: every transformation given by a matrix is linear.)

\textbf{Task 2.} The kernel is the set of $(x,y)$ with $T(x,y)=(0,0)$:
\[ \begin{cases} 2x+y=0 \\ x+2y=0 \end{cases} \]
From the first equation $y=-2x$; substituting into the second, $x-4x=0$, so $x=0$ and $y=0$. Therefore
\[ \ker T = \{(0,0)\}. \]
Since the kernel is reduced to the zero vector, $T$ is \cle{injective}. By the rank--nullity theorem, $\dim\operatorname{Im} T = 2-0 = 2$, so $\operatorname{Im} T = \mathbb{R}^2$: $T$ is \cle{surjective}, hence \cle{bijective}. Also $\det A = 2\times 2 - 1\times 1 = 3 \neq 0$, which confirms bijectivity. \emph{Interpretation:} no two distinct points of the city ever collapse onto the same point of the map, and every point of the map corresponds to a real point of the city.

\textbf{Task 3.} Using the formula for the inverse of a $2\times 2$ matrix,
\[ A^{-1} = \frac{1}{\det A}\begin{pmatrix} 2 & -1 \\ -1 & 2 \end{pmatrix} = \frac{1}{3}\begin{pmatrix} 2 & -1 \\ -1 & 2 \end{pmatrix}. \]
Hence the inverse transformation is
\[ T^{-1}(x',y') = \left(\tfrac{2x'-y'}{3},\; \tfrac{-x'+2y'}{3}\right). \]
For the landmark at $(x',y')=(4,5)$:
\[ x = \frac{2(4)-5}{3} = 1, \qquad y = \frac{-4+2(5)}{3} = 2. \]
The landmark is at $(1,2)$ in the city, i.e. $1$ km East and $2$ km North of the Poste Centrale. Check: $T(1,2)=(2+2,\,1+4)=(4,5)$. \checkmark

\textbf{Task 4.} Invariant points satisfy $T(x,y)=(x,y)$:
\[ \begin{cases} 2x+y=x \\ x+2y=y \end{cases} \iff \begin{cases} x+y=0 \\ x+y=0 \end{cases} \iff y=-x. \]
The set of invariant points is the line $y=-x$. \emph{Interpretation:} every point of the city lying on the South-East/North-West diagonal through the Poste Centrale is drawn exactly at its true position on the map.

\textbf{Task 5.} Parametrise the avenue: a point of $\mathcal{D}$ is $(t,\;t+1)$, $t\in\mathbb{R}$. Its image is
\[ T(t,t+1) = \big(2t+(t+1),\; t+2(t+1)\big) = (3t+1,\; 3t+2). \]
So $x'=3t+1$ and $y'=3t+2$. Eliminating $t$: $t=\dfrac{x'-1}{3}$, hence
\[ y' = 3\cdot\frac{x'-1}{3}+2 = x'+1. \]
The image of the avenue is the straight line $y'=x'+1$. Remarkably, this line is \emph{globally invariant}: the avenue is mapped onto itself (though each of its points moves along it, except the invariant point where $\mathcal{D}$ meets $y=-x$, namely $\left(-\tfrac12,\tfrac12\right)$).

\textbf{Task 6.} The characteristic polynomial is
\[ \det(A-\lambda I) = \begin{vmatrix} 2-\lambda & 1 \\ 1 & 2-\lambda \end{vmatrix} = (2-\lambda)^2-1 = \lambda^2-4\lambda+3 = (\lambda-1)(\lambda-3). \]
The eigenvalues are $\lambda_1=1$ and $\lambda_2=3$.

For $\lambda_1=1$: $(A-I)\vec{v}=\vec{0}$ gives $x+y=0$, so the eigenspace is
\[ E_1 = \operatorname{span}\begin{pmatrix} 1 \\ -1 \end{pmatrix}, \quad \text{the line } y=-x. \]
For $\lambda_2=3$: $(A-3I)\vec{v}=\vec{0}$ gives $-x+y=0$, so the eigenspace is
\[ E_3 = \operatorname{span}\begin{pmatrix} 1 \\ 1 \end{pmatrix}, \quad \text{the line } y=x. \]

\emph{Interpretation:} the direction $y=-x$ (the diagonal through the Poste Centrale) is left \emph{fixed} (stretch factor $1$) — this is exactly the line of invariant points found in Task 4. The perpendicular direction $y=x$ is \emph{stretched by a factor $3$}, without rotation. Every other direction is a mixture of these two effects.

\begin{popfigure}
\begin{center}
\begin{tikzpicture}[scale=1.1]
\draw[->,gray] (-2.6,0) -- (2.6,0) node[right] {\small $x$};
\draw[->,gray] (0,-2.6) -- (0,2.6) node[above] {\small $y$};
\draw[popblue,very thick] (-2.2,-2.2) -- (2.2,2.2) node[above right, text width=2.5cm, align=center] {\small $y=x$\\ \small stretch $\times 3$};
\draw[poporange,very thick] (-1.9,1.9) -- (1.9,-1.9) node[below right, text width=2.5cm, align=center] {\small $y=-x$\\ \small fixed};
\draw[popgreen,thick,dashed] (-2.3,-1.3) -- (2.3,3.3) node[above left] {\small avenue $y=x+1$};
\fill[popdark] (0,0) circle (1.5pt) node[below left] {\small $O$};
\fill[popdark] (-0.5,0.5) circle (1.5pt) node[above left] {\small invariant};
\end{tikzpicture}
\end{center}
\legende{Eigen-directions of the projection and image of the avenue.}
\end{popfigure}

\textbf{Task 7.} A linear transformation multiplies all areas by $|\det A| = 3$. Hence every area on the map is $3$ times the true area. The plot of $2\ \mathrm{km^2}$ occupies $3\times 2 = 6\ \mathrm{km^2}$ on the map. Note the coherence with Task 6: the area factor equals the product of the eigenvalues, $1\times 3 = 3$.

\textbf{Task 8.} The poster should display: the matrix $A$ and its inverse; $\ker T=\{\vec 0\}$ and $\operatorname{Im}T=\mathbb{R}^2$ with their interpretation; the invariant line $y=-x$; the image line $y'=x'+1$; the eigenvalues $1$ and $3$ with their eigen-directions; the area factor $3$; and the diagram above.

\begin{attention}[Pitfalls to avoid in this activity]
\begin{itemize}
\item Do not confuse the invariant \emph{points} (solutions of $T(\vec v)=\vec v$, eigenvalue $1$) with a globally invariant \emph{line} (a line mapped onto itself, like the avenue here).
\item The area factor is $|\det A|$, not $\det A$ in general — here $\det A=3>0$ so both coincide.
\item When eliminating the parameter in Task 5, keep the coordinates $(x',y')$ of the \emph{image} plane; do not mix them with the city coordinates.
\end{itemize}
\end{attention}

\newpage

\coursec{Methods and worked examples}

\coursec{Linear Transformations and Their Properties}

\begin{definition}[Linear Transformation]
A map $T:\mathbb{R}^n\to\mathbb{R}^m$ is a \cle{linear transformation} if for all vectors $\vec{u},\vec{v}\in\mathbb{R}^n$ and all scalars $\lambda,\mu\in\mathbb{R}$:
\[
T(\lambda\vec{u}+\mu\vec{v})=\lambda T(\vec{u})+\mu T(\vec{v}).
\]
Equivalently, $T$ satisfies:
\begin{enumerate}
\item \textbf{Additivity:} $T(\vec{u}+\vec{v})=T(\vec{u})+T(\vec{v})$,
\item \textbf{Homogeneity:} $T(\lambda\vec{u})=\lambda T(\vec{u})$.
\end{enumerate}
\end{definition}

\begin{aretenir}[Immediate consequences]
If $T$ is linear then:
\begin{itemize}
\item $T(\vec{0})=\vec{0}$ (take $\lambda=\mu=0$),
\item $T(-\vec{u})=-T(\vec{u})$,
\item $T(\vec{u}-\vec{v})=T(\vec{u})-T(\vec{v})$.
\end{itemize}
\end{aretenir}

\begin{methode}[Proving that $T:\mathbb{R}^n\to\mathbb{R}^m$ is linear]
\begin{enumerate}
\item \textbf{Identify the formula.} Write $T(\vec{x})$ explicitly in coordinates.
\item \textbf{Quick test (necessary condition).} Check whether $T(\vec{0})=\vec{0}$. If not, $T$ is \emph{not} linear: stop.
\item \textbf{Prove additivity.} Take arbitrary $\vec{u},\vec{v}$ and show $T(\vec{u}+\vec{v})=T(\vec{u})+T(\vec{v})$.
\item \textbf{Prove homogeneity.} Take arbitrary $\lambda\in\mathbb{R}$ and show $T(\lambda\vec{u})=\lambda T(\vec{u})$.
\item \textbf{Conclude.} Both conditions hold, so $T$ is linear. (Equivalently, verify in one step that $T(\lambda\vec{u}+\mu\vec{v})=\lambda T(\vec{u})+\mu T(\vec{v})$.)
\end{enumerate}
\textbf{Reflex:} any map given by \emph{linear} expressions in the coordinates (no constant term, no squares, no products of variables) is linear. Terms like $x^2$, $xy$, $x+3$ are danger signals.
\end{methode}

\begin{exemplebox}[Testing linearity]
Determine which of the following maps are linear transformations. Justify your answers.
\begin{enumerate}
\item $T_1:\mathbb{R}^2\to\mathbb{R}^2$, \quad $T_1(x,y)=(2x-3y,\;x+y)$.
\item $T_2:\mathbb{R}^2\to\mathbb{R}^2$, \quad $T_2(x,y)=(x+1,\;y)$.
\item $T_3:\mathbb{R}^2\to\mathbb{R}$, \quad $T_3(x,y)=xy$.
\item $T_4:\mathbb{R}^3\to\mathbb{R}^2$, \quad $T_4(x,y,z)=(x+2y-z,\;3x-y+4z)$.
\end{enumerate}
\end{exemplebox}

\begin{exoresolu}[Testing linearity]
\textbf{Solution.}
\begin{enumerate}
\item \textbf{For $T_1$.} First, $T_1(0,0)=(0,0)$, so the necessary condition passes. Let $\vec{u}=(x_1,y_1)$, $\vec{v}=(x_2,y_2)$ and $\lambda\in\mathbb{R}$.
\begin{align*}
T_1(\vec{u}+\vec{v})&=T_1(x_1+x_2,\;y_1+y_2)\\
&=\bigl(2(x_1+x_2)-3(y_1+y_2),\;(x_1+x_2)+(y_1+y_2)\bigr)\\
&=(2x_1-3y_1,\;x_1+y_1)+(2x_2-3y_2,\;x_2+y_2)=T_1(\vec{u})+T_1(\vec{v}).
\end{align*}
Also,
\[
T_1(\lambda\vec{u})=\bigl(2\lambda x_1-3\lambda y_1,\;\lambda x_1+\lambda y_1\bigr)=\lambda T_1(\vec{u}).
\]
Both conditions hold, hence $T_1$ \textbf{is a linear transformation}.

\item \textbf{For $T_2$.} $T_2(0,0)=(1,0)\neq(0,0)$. A linear transformation must send $\vec{0}$ to $\vec{0}$, so $T_2$ \textbf{is not linear}. (It is a translation, an \emph{affine} map.)

\item \textbf{For $T_3$.} Although $T_3(0,0)=0$, take $\vec{u}=(1,1)$ and $\lambda=2$:
\[
T_3(2\vec{u})=T_3(2,2)=4,\qquad 2T_3(\vec{u})=2\times 1=2.
\]
Since $T_3(2\vec{u})\neq 2T_3(\vec{u})$, homogeneity fails and $T_3$ \textbf{is not linear}.

\item \textbf{For $T_4$.} $T_4(0,0,0)=(0,0)$. Let $\vec{u}=(x_1,y_1,z_1)$, $\vec{v}=(x_2,y_2,z_2)$, $\lambda,\mu\in\mathbb{R}$.
\begin{align*}
T_4(\lambda\vec{u}+\mu\vec{v})&=T_4(\lambda x_1+\mu x_2,\;\lambda y_1+\mu y_2,\;\lambda z_1+\mu z_2)\\
&=\bigl((\lambda x_1+\mu x_2)+2(\lambda y_1+\mu y_2)-(\lambda z_1+\mu z_2),\\
&\qquad 3(\lambda x_1+\mu x_2)-(\lambda y_1+\mu y_2)+4(\lambda z_1+\mu z_2)\bigr)\\
&=\lambda(x_1+2y_1-z_1,\;3x_1-y_1+4z_1)+\mu(x_2+2y_2-z_2,\;3x_2-y_2+4z_2)\\
&=\lambda T_4(\vec{u})+\mu T_4(\vec{v}).
\end{align*}
Hence $T_4$ \textbf{is a linear transformation}.
\end{enumerate}
\end{exoresolu}

\begin{attention}[Common GCE traps]
\begin{itemize}
\item Forgetting to test $T(\vec{0})=\vec{0}$ first — it can save time.
\item Confusing \emph{linear} with \emph{affine}: $T(\vec{x})=A\vec{x}+\vec{b}$ with $\vec{b}\neq\vec{0}$ is \textbf{not} linear.
\item Assuming $T(\vec{u}\cdot\vec{v})=T(\vec{u})\cdot T(\vec{v})$ or $T(\vec{u}\times\vec{v})=T(\vec{u})\times T(\vec{v})$ — these are false in general.
\end{itemize}
\end{attention}

\courssub{Matrix of a Linear Transformation}

\begin{propriete}[Matrix representation]
Every linear transformation $T:\mathbb{R}^n\to\mathbb{R}^m$ can be represented by a unique $m\times n$ matrix $A$ such that $T(\vec{x})=A\vec{x}$ for all $\vec{x}\in\mathbb{R}^n$. The columns of $A$ are the images of the standard basis vectors:
\[
A=\begin{pmatrix} | & | & & | \\ T(\vec{e}_1) & T(\vec{e}_2) & \cdots & T(\vec{e}_n) \\ | & | & & | \end{pmatrix}.
\]
\end{propriete}

\begin{methode}[Finding the matrix of $T$ relative to the standard bases]
\begin{enumerate}
\item \textbf{Compute} $T(\vec{e}_1), T(\vec{e}_2), \dots, T(\vec{e}_n)$ using the definition of $T$.
\item \textbf{Write} each image as a column vector.
\item \textbf{Assemble} the matrix $A$ by placing these columns side by side.
\item \textbf{Check:} verify on one arbitrary vector that $T(\vec{x})=A\vec{x}$.
\end{enumerate}
\end{methode}

\begin{exoresolu}[Matrix from the definition]
A linear transformation $T:\mathbb{R}^3\to\mathbb{R}^2$ is defined by
\[
T(x,y,z)=(x+2y-z,\;3x+y+z).
\]
Find the matrix $A$ of $T$ with respect to the standard bases, and compute $T(1,-1,2)$.
\tcblower
\textbf{Solution.} The standard basis of $\mathbb{R}^3$ is $\vec{e}_1=(1,0,0)$, $\vec{e}_2=(0,1,0)$, $\vec{e}_3=(0,0,1)$. Compute the images:
\[
T(\vec{e}_1)=(1,3),\qquad T(\vec{e}_2)=(2,1),\qquad T(\vec{e}_3)=(-1,1).
\]
Placing these as columns:
\[
A=\begin{pmatrix} 1 & 2 & -1 \\ 3 & 1 & 1 \end{pmatrix}.
\]
\textbf{Check and computation:}
\[
T(1,-1,2)=A\begin{pmatrix}1\\-1\\2\end{pmatrix}
=\begin{pmatrix} 1(1)+2(-1)+(-1)(2) \\ 3(1)+1(-1)+1(2)\end{pmatrix}
=\begin{pmatrix} -3 \\ 4 \end{pmatrix}.
\]
Direct verification: $T(1,-1,2)=(1-2-2,\;3-1+2)=(-3,4)$. \checkmark
\end{exoresolu}

\begin{exemplebox}[Matrix from images of a basis]
Let $T:\mathbb{R}^2\to\mathbb{R}^2$ be a linear transformation such that
$T(1,1)=(3,2)$ and $T(1,-1)=(1,4)$. Find the matrix of $T$ relative to the standard basis.
\end{exemplebox}

\begin{exoresolu}[Matrix from images of a basis]
\textbf{Solution.} We need $T(1,0)$ and $T(0,1)$. Express the standard basis in terms of the given vectors:
\[
(1,0)=\tfrac12(1,1)+\tfrac12(1,-1),\qquad (0,1)=\tfrac12(1,1)-\tfrac12(1,-1).
\]
By linearity:
\begin{align*}
T(1,0)&=\tfrac12 T(1,1)+\tfrac12 T(1,-1)=\tfrac12(3,2)+\tfrac12(1,4)=(2,3),\\
T(0,1)&=\tfrac12 T(1,1)-\tfrac12 T(1,-1)=\tfrac12(3,2)-\tfrac12(1,4)=(1,-1).
\end{align*}
Hence the matrix is $A=\begin{pmatrix} 2 & 1 \\ 3 & -1 \end{pmatrix}$.
\end{exoresolu}

\begin{savaistu}[Application: Computer graphics]
Linear transformations are the backbone of 2D/3D computer graphics. Rotations, scaling, shearing, and reflections are all linear. In Cameroon's growing tech sector (e.g., animation studios in Douala, game developers in Yaoundé), matrices like $\begin{pmatrix}\cos\theta & -\sin\theta\\ \sin\theta & \cos\theta\end{pmatrix}$ rotate objects by angle $\theta$ — a direct application of the theory you are mastering.
\end{savaistu}

\coursec{Kernel and Image of a Linear Transformation}

\begin{definition}[Kernel and Image]
Let $T:\mathbb{R}^n\to\mathbb{R}^m$ be a linear transformation with matrix $A$.
\begin{itemize}
\item The \cle{kernel} of $T$ is $\ker T=\{\vec{x}\in\mathbb{R}^n : T(\vec{x})=\vec{0}\}=\{\vec{x}:A\vec{x}=\vec{0}\}$.
\item The \cle{image} of $T$ is $\operatorname{Im} T=\{T(\vec{x}):\vec{x}\in\mathbb{R}^n\}=\{A\vec{x}:\vec{x}\in\mathbb{R}^n\}$, the column space of $A$.
\end{itemize}
\end{definition}

\begin{propriete}[Subspace properties]
$\ker T$ is a subspace of $\mathbb{R}^n$ (the domain). $\operatorname{Im} T$ is a subspace of $\mathbb{R}^m$ (the codomain).
\end{propriete}

\begin{methode}[Finding $\ker T$ and $\operatorname{Im} T$]
Let $T:\mathbb{R}^n\to\mathbb{R}^m$ have matrix $A$.
\begin{enumerate}
\item \textbf{Kernel.} Solve $A\vec{x}=\vec{0}$ (a homogeneous system) by row reduction (Gaussian elimination). Express the solution set with free parameters. This set is $\ker T$. Its dimension is the \cle{nullity} $\nu(T)$.
\item \textbf{Image.} $\operatorname{Im} T$ is spanned by the \emph{columns} of $A$: $\operatorname{Im} T=\operatorname{span}\{T(\vec{e}_1),\dots,T(\vec{e}_n)\}$. Extract a basis by keeping only linearly independent columns (pivot columns in the row-echelon form of $A$). Its dimension is the \cle{rank} $\rho(T)$.
\item \textbf{Verify the Rank--Nullity Theorem:} $\dim\ker T+\dim\operatorname{Im} T=n$.
\item \textbf{Conclusions:} $T$ is injective (one-to-one) $\iff \ker T=\{\vec{0}\}$; $T$ is surjective (onto) $\iff \operatorname{Im} T=\mathbb{R}^m$.
\end{enumerate}
\end{methode}

\begin{exoresolu}[Kernel and image]
Let $T:\mathbb{R}^3\to\mathbb{R}^3$ be the linear transformation with matrix
\[
A=\begin{pmatrix} 1 & 2 & 1 \\ 2 & 4 & 3 \\ 1 & 2 & 2 \end{pmatrix}.
\]
Find a basis and the dimension of $\ker T$ and of $\operatorname{Im} T$. Is $T$ injective? Surjective?
\tcblower
\textbf{Solution.}

\textbf{Step 1: kernel.} Solve $A\vec{x}=\vec{0}$. Row-reduce:
\[
\begin{pmatrix} 1 & 2 & 1 \\ 2 & 4 & 3 \\ 1 & 2 & 2 \end{pmatrix}
\xrightarrow[R_3\to R_3-R_1]{R_2\to R_2-2R_1}
\begin{pmatrix} 1 & 2 & 1 \\ 0 & 0 & 1 \\ 0 & 0 & 1 \end{pmatrix}
\xrightarrow{R_3\to R_3-R_2}
\begin{pmatrix} 1 & 2 & 1 \\ 0 & 0 & 1 \\ 0 & 0 & 0 \end{pmatrix}.
\]
The system gives $z=0$ and $x+2y=0$, i.e. $x=-2y$ with $y$ free. Setting $y=t$:
\[
\ker T=\left\{t\begin{pmatrix}-2\\1\\0\end{pmatrix}:t\in\mathbb{R}\right\}
=\operatorname{span}\left\{\begin{pmatrix}-2\\1\\0\end{pmatrix}\right\},\qquad \dim\ker T=1.
\]

\textbf{Step 2: image.} The columns are $\vec{c}_1=(1,2,1)$, $\vec{c}_2=(2,4,2)=2\vec{c}_1$, $\vec{c}_3=(1,3,2)$. Since $\vec{c}_2$ is a multiple of $\vec{c}_1$ while $\vec{c}_1$ and $\vec{c}_3$ are not proportional, a basis is $\{\vec{c}_1,\vec{c}_3\}$:
\[
\operatorname{Im} T=\operatorname{span}\left\{\begin{pmatrix}1\\2\\1\end{pmatrix},\begin{pmatrix}1\\3\\2\end{pmatrix}\right\},\qquad \dim\operatorname{Im} T=2.
\]

\textbf{Step 3: verification.} $\dim\ker T+\dim\operatorname{Im} T=1+2=3=\dim\mathbb{R}^3$. \checkmark

\textbf{Conclusion.} Since $\ker T\neq\{\vec{0}\}$, $T$ is \textbf{not injective}; since $\operatorname{Im} T\neq\mathbb{R}^3$ (dimension $2<3$), $T$ is \textbf{not surjective}. Geometrically, $T$ collapses $\mathbb{R}^3$ onto a plane.
\end{exoresolu}

\begin{popfigure}
\begin{tikzpicture}[scale=0.8]
% Domain R^3
\draw[popblue, thick] (0,0) circle (1.5);
\node[popblue] at (0,-2) {Domain $\mathbb{R}^3$};
\node[popblue] at (0,0.5) {$\ker T$ (line)};
\draw[->, popblue] (0,0) -- (0.8,0.5) node[right] {$\vec{v}$};
\draw[->, popblue] (0,0) -- (-0.8,-0.5) node[left] {$-\vec{v}$};
% Arrow
\draw[poporange, thick, ->] (2.5,0) -- (4.5,0) node[midway, above] {$T$};
% Codomain R^3
\draw[popgreen, thick] (6,0) circle (1.5);
\node[popgreen] at (6,-2) {Codomain $\mathbb{R}^3$};
\node[popgreen] at (6,0.5) {$\operatorname{Im} T$ (plane)};
\draw[->, popgreen] (6,0) -- (7,0.5) node[right] {$T(\vec{v})$};
\draw[->, popgreen] (6,0) -- (5,-0.5) node[left] {$T(-\vec{v})$};
% Kernel vectors map to 0
\draw[dashed, popmuted] (0.8,0.5) -- (6,0);
\draw[dashed, popmuted] (-0.8,-0.5) -- (6,0);
\node[popmuted] at (3,1.2) {$\vec{v}\mapsto\vec{0}$};
\end{tikzpicture}
\legende{Geometric interpretation: kernel vectors collapse to $\vec{0}$; the image is a plane.}
\end{popfigure}

\begin{exercice}[Practice: kernel and image]
Let $T:\mathbb{R}^4\to\mathbb{R}^3$ have matrix $A=\begin{pmatrix} 1 & 2 & 0 & 1 \\ 0 & 1 & 1 & 1 \\ 1 & 3 & 1 & 2 \end{pmatrix}$.
\begin{enumerate}
\item Find a basis for $\ker T$ and state its dimension.
\item Find a basis for $\operatorname{Im} T$ and state its dimension.
\item Verify the rank--nullity theorem.
\item Is $T$ injective? Surjective?
\end{enumerate}
\end{exercice}

\begin{corrige}[Exercise 1]
\textbf{Solution.}
Row-reduce $A$:
\[
\begin{pmatrix} 1 & 2 & 0 & 1 \\ 0 & 1 & 1 & 1 \\ 1 & 3 & 1 & 2 \end{pmatrix}
\xrightarrow{R_3\to R_3-R_1}
\begin{pmatrix} 1 & 2 & 0 & 1 \\ 0 & 1 & 1 & 1 \\ 0 & 1 & 1 & 1 \end{pmatrix}
\xrightarrow{R_3\to R_3-R_2}
\begin{pmatrix} 1 & 2 & 0 & 1 \\ 0 & 1 & 1 & 1 \\ 0 & 0 & 0 & 0 \end{pmatrix}
\xrightarrow{R_1\to R_1-2R_2}
\begin{pmatrix} 1 & 0 & -2 & -1 \\ 0 & 1 & 1 & 1 \\ 0 & 0 & 0 & 0 \end{pmatrix}.
\]
Pivot columns: 1 and 2. Free variables: $x_3=s$, $x_4=t$.
System: $x_1-2x_3-x_4=0 \Rightarrow x_1=2s+t$; $x_2+x_3+x_4=0 \Rightarrow x_2=-s-t$.
\[
\ker T=\left\{s\begin{pmatrix}2\\-1\\1\\0\end{pmatrix}+t\begin{pmatrix}1\\-1\\0\\1\end{pmatrix}:s,t\in\mathbb{R}\right\},\quad \dim\ker T=2.
\]
Basis of $\operatorname{Im} T$: pivot columns of original $A$: $\begin{pmatrix}1\\0\\1\end{pmatrix},\begin{pmatrix}2\\1\\3\end{pmatrix}$. $\dim\operatorname{Im} T=2$.
Rank--nullity: $2+2=4=\dim\mathbb{R}^4$. \checkmark
Not injective ($\ker T\neq\{\vec{0}\}$), not surjective ($\dim\operatorname{Im} T=2<3$).
\end{corrige}

\coursec{Inverse of a Linear Transformation and Invariant Points}

\begin{propriete}[Invertibility criteria]
For a linear transformation $T:\mathbb{R}^n\to\mathbb{R}^n$ with matrix $A$, the following are equivalent:
\begin{enumerate}
\item $T$ is invertible (bijective).
\item $A$ is invertible.
\item $\det A\neq 0$.
\item $\ker T=\{\vec{0}\}$.
\item $\operatorname{Im} T=\mathbb{R}^n$.
\item $0$ is not an eigenvalue of $A$.
\end{enumerate}
\end{propriete}

\begin{methode}[Inverse transformation and invariant points]
Let $T$ have matrix $A$.
\begin{enumerate}
\item \textbf{Invertibility.} Check $\det A\neq 0$ (or $\ker T=\{\vec{0}\}$).
\item \textbf{Inverse.} Then $T^{-1}$ is linear with matrix $A^{-1}$. For a $2\times 2$ matrix, $A^{-1}=\dfrac{1}{\det A}\begin{pmatrix} d & -b \\ -c & a \end{pmatrix}$.
\item \textbf{Invariant (fixed) points.} Solve $T(\vec{x})=\vec{x}$, i.e. $(A-I)\vec{x}=\vec{0}$. The fixed points form a subspace (at least $\vec{0}$).
\item \textbf{Interpretation.} A line of fixed points means $T$ acts as the identity on that line (e.g. a reflection fixes its axis).
\end{enumerate}
\begin{attention}[Trap]
Do not confuse \emph{invariant points} ($T(\vec{x})=\vec{x}$, pointwise fixed) with \emph{invariant lines} ($T$ maps the line into itself, points may move along it). Read the question carefully!
\end{attention}
\end{methode}

\begin{exoresolu}[Inverse and fixed points]
The linear transformation $T:\mathbb{R}^2\to\mathbb{R}^2$ has matrix $A=\begin{pmatrix} 3 & -1 \\ 5 & -2 \end{pmatrix}$.
\begin{enumerate}
\item Show that $T$ is invertible and find the matrix of $T^{-1}$.
\item Determine all invariant points of $T$.
\end{enumerate}
\tcblower
\textbf{Solution.}
\begin{enumerate}
\item $\det A=(3)(-2)-(-1)(5)=-6+5=-1\neq 0$, so $A$ is invertible and $T$ is invertible.
\[
A^{-1}=\frac{1}{-1}\begin{pmatrix} -2 & 1 \\ -5 & 3 \end{pmatrix}=\begin{pmatrix} 2 & -1 \\ 5 & -3 \end{pmatrix}.
\]
Hence $T^{-1}(x,y)=(2x-y,\;5x-3y)$. \textbf{Check:} $A^{-1}A=\begin{pmatrix}2&-1\\5&-3\end{pmatrix}\begin{pmatrix}3&-1\\5&-2\end{pmatrix}=\begin{pmatrix}1&0\\0&1\end{pmatrix}$. \checkmark

\item Invariant points satisfy $(A-I)\vec{x}=\vec{0}$:
\[
A-I=\begin{pmatrix} 2 & -1 \\ 5 & -3 \end{pmatrix},\qquad
\begin{cases} 2x-y=0 \\ 5x-3y=0.\end{cases}
\]
From the first equation $y=2x$; substituting: $5x-6x=-x=0$, so $x=0$, $y=0$. The \textbf{only invariant point is the origin} $(0,0)$.
\end{enumerate}
\end{exoresolu}

\begin{exemplebox}[Invariant points of a projection]
Let $T:\mathbb{R}^2\to\mathbb{R}^2$ be the orthogonal projection onto the line $y=x$. Its matrix is $A=\frac12\begin{pmatrix}1&1\\1&1\end{pmatrix}$. Find the invariant points.
\end{exemplebox}

\begin{exoresolu}[Invariant points of a projection]
\textbf{Solution.} Solve $(A-I)\vec{x}=\vec{0}$:
\[
A-I=\frac12\begin{pmatrix}-1&1\\1&-1\end{pmatrix}\sim\begin{pmatrix}1&-1\\0&0\end{pmatrix}\Rightarrow x-y=0.
\]
The invariant points are exactly the line $y=x$ (the axis of projection). Every point on this line is fixed; points off the line move onto it.
\end{exoresolu}

\coursec{Images of Lines and Planes}

\begin{methode}[Image of a line or plane under $T$]
\begin{enumerate}
\item \textbf{Parametrise} the object: a line through $\vec{a}$ with direction $\vec{d}$ is $\vec{r}=\vec{a}+t\vec{d}$; a plane is $\vec{r}=\vec{a}+s\vec{u}+t\vec{v}$.
\item \textbf{Apply linearity:} $T(\vec{r})=T(\vec{a})+t\,T(\vec{d})$ (respectively $T(\vec{a})+s\,T(\vec{u})+t\,T(\vec{v})$).
\item \textbf{Read off the image:} the image is a line through $T(\vec{a})$ with direction $T(\vec{d})$ — \emph{unless} $T(\vec{d})=\vec{0}$, in which case the image is the single point $T(\vec{a})$.
\item \textbf{For a plane:} if $T(\vec{u})$ and $T(\vec{v})$ are independent, the image is a plane; if they are parallel but non-zero, a line; if both zero, a point.
\end{enumerate}
\textbf{Reflex:} linear transformations map lines to lines (or points) and planes to planes (or lines or points): they can only \emph{lower} the dimension, never raise it.
\end{methode}

\begin{exoresolu}[Image of a line and of a plane]
Let $T:\mathbb{R}^3\to\mathbb{R}^3$ be defined by $T(x,y,z)=(x+y,\;y+z,\;x+2y+z)$.
\begin{enumerate}
\item Find the image of the line $L:\ \vec{r}=(1,0,2)+t(1,-1,1)$.
\item Find the image of the plane $P:\ \vec{r}=(0,0,0)+s(1,0,0)+t(0,1,0)$.
\end{enumerate}
\tcblower
\textbf{Solution.} Note first that the matrix of $T$ is $A=\begin{pmatrix}1&1&0\\0&1&1\\1&2&1\end{pmatrix}$, and the third row is the sum of the first two, so $\operatorname{rank} A=2$: images will live in a plane.
\begin{enumerate}
\item Compute $T(1,0,2)=(1+0,\;0+2,\;1+0+2)=(1,2,3)$ and $T(1,-1,1)=(1-1,\;-1+1,\;1-2+1)=(0,0,0)$. The direction vector lies in $\ker T$! Hence
\[
T(\vec{r})=(1,2,3)+t(0,0,0)=(1,2,3).
\]
The image of $L$ is the \textbf{single point} $(1,2,3)$: the whole line collapses.

\item $T(1,0,0)=(1,0,1)$ and $T(0,1,0)=(1,1,2)$. These are not proportional, so they are independent. The image is the \textbf{plane}
\[
\vec{r}=s(1,0,1)+t(1,1,2),
\]
i.e. the plane through the origin spanned by $(1,0,1)$ and $(1,1,2)$. (Cartesian form: a normal is $(1,0,1)\times(1,1,2)=(-1,-1,1)$, giving $x+y-z=0$.)
\end{enumerate}
\end{exoresolu}

\begin{popfigure}
\begin{tikzpicture}[scale=0.7]
% Original line and plane
\draw[popblue, thick] (-2,0,0) -- (2,0,0) node[right] {$x$};
\draw[popblue, thick] (0,-2,0) -- (0,2,0) node[above] {$y$};
\draw[popblue, thick] (0,0,-2) -- (0,0,2) node[above] {$z$};
% Line L: (1,0,2) + t(1,-1,1)
\draw[poporange, thick, ->] (1,0,2) -- (2,-1,3) node[above right] {$L$};
\draw[poporange, thick, ->] (1,0,2) -- (0,1,1);
% Plane P: xy-plane
\fill[popgreen, opacity=0.1] (-1.5,-1.5,0) -- (1.5,-1.5,0) -- (1.5,1.5,0) -- (-1.5,1.5,0) -- cycle;
\draw[popgreen, thick] (-1.5,-1.5,0) -- (1.5,-1.5,0) -- (1.5,1.5,0) -- (-1.5,1.5,0) -- cycle;
\node[popgreen] at (0,0,0.5) {$P$};
% Transformation arrow
\draw[poppurple, thick, ->] (3,0,0) -- (5,0,0) node[midway, above] {$T$};
% Image: point and plane
\begin{scope}[xshift=6cm]
\draw[popblue, thick] (-2,0,0) -- (2,0,0) node[right] {$x$};
\draw[popblue, thick] (0,-2,0) -- (0,2,0) node[above] {$y$};
\draw[popblue, thick] (0,0,-2) -- (0,0,2) node[above] {$z$};
% Image of L: point (1,2,3)
\fill[poporange] (1,2,3) circle (3pt);
\node[poporange] at (1,2,3.5) {$T(L)$};
% Image of P: plane x+y-z=0
\fill[popgreen, opacity=0.1] (-1.5,-1.5,-3) -- (1.5,-1.5,0) -- (1.5,1.5,3) -- (-1.5,1.5,0) -- cycle;
\draw[popgreen, thick] (-1.5,-1.5,-3) -- (1.5,-1.5,0) -- (1.5,1.5,3) -- (-1.5,1.5,0) -- cycle;
\node[popgreen] at (0,0,0.5) {$T(P)$};
\end{scope}
\end{tikzpicture}
\legende{Line $L$ collapses to a point because its direction is in $\ker T$; plane $P$ maps to another plane.}
\end{popfigure}

\begin{exercice}[Practice: images of lines and planes]
Let $T:\mathbb{R}^3\to\mathbb{R}^3$ have matrix $A=\begin{pmatrix} 1 & 0 & 1 \\ 0 & 1 & 1 \\ 0 & 0 & 0 \end{pmatrix}$.
\begin{enumerate}
\item Find the image of the line $\vec{r}=(1,2,3)+t(1,1,-2)$.
\item Find the image of the plane $x+2y-z=4$.
\item Describe geometrically what $T$ does to $\mathbb{R}^3$.
\end{enumerate}
\end{exercice}

\begin{corrige}[Exercise 2]
\textbf{Solution.}
$T(x,y,z)=(x+z,\;y+z,\;0)$.
\begin{enumerate}
\item $T(1,2,3)=(4,5,0)$, $T(1,1,-2)=(1-2,\;1-2,\;0)=(-1,-1,0)\neq\vec{0}$. Image: line $\vec{r}=(4,5,0)+t(-1,-1,0)$ in the $xy$-plane.
\item Plane: $\vec{r}=(4,0,0)+s(1,0,1)+t(0,1,2)$ (since $z=x+2y-4$). $T(4,0,0)=(4,0,0)$, $T(1,0,1)=(2,1,0)$, $T(0,1,2)=(2,3,0)$. These two direction vectors are independent, so image is the whole $xy$-plane ($z=0$).
\item $T$ projects $\mathbb{R}^3$ onto the $xy$-plane along the direction $(1,1,-1)$ (kernel).
\end{enumerate}
\end{corrige}

\coursec{Eigenvalues and Eigenvectors}

\begin{definition}[Eigenvalues and Eigenvectors]
Let $A$ be an $n\times n$ matrix. A scalar $\lambda$ is an \cle{eigenvalue} of $A$ if there exists a non-zero vector $\vec{v}$ such that $A\vec{v}=\lambda\vec{v}$. Such a vector $\vec{v}$ is called an \cle{eigenvector} corresponding to $\lambda$. The set of all eigenvectors for $\lambda$ together with $\vec{0}$ is the \cle{eigenspace} $E_\lambda=\ker(A-\lambda I)$.
\end{definition}

\begin{propriete}[Characteristic equation]
The eigenvalues of $A$ are the roots of the \cle{characteristic polynomial} $p(\lambda)=\det(A-\lambda I)=0$.
\end{propriete}

\begin{methode}[Finding eigenvalues and eigenvectors]
Given a square matrix $A$:
\begin{enumerate}
\item \textbf{Characteristic equation.} Compute $\det(A-\lambda I)=0$ and solve for $\lambda$: the roots are the \cle{eigenvalues}.
\item \textbf{Eigenvectors.} For each eigenvalue $\lambda_i$, solve $(A-\lambda_i I)\vec{x}=\vec{0}$. The non-zero solutions are the \cle{eigenvectors}; the solution space is the \cle{eigenspace} $E_{\lambda_i}$.
\item \textbf{Checks.} (i) Sum of eigenvalues $=\operatorname{tr}(A)$; (ii) product of eigenvalues $=\det A$; (iii) verify $A\vec{v}=\lambda\vec{v}$ on one eigenvector.
\item \textbf{Geometric meaning.} An eigenvector gives an \emph{invariant line} through the origin: $T$ stretches it by the factor $\lambda$ (reversing direction if $\lambda<0$).
\end{enumerate}
\end{methode}

\begin{exoresolu}[Eigenvalues and eigenvectors of a $2\times2$ matrix]
Find the eigenvalues and corresponding eigenvectors of $A=\begin{pmatrix} 4 & 2 \\ 1 & 3 \end{pmatrix}$, and describe the geometric action of the associated transformation on each eigenspace.
\tcblower
\textbf{Solution.}

\textbf{Step 1: eigenvalues.}
\[
\det(A-\lambda I)=\begin{vmatrix} 4-\lambda & 2 \\ 1 & 3-\lambda \end{vmatrix}
=(4-\lambda)(3-\lambda)-2=\lambda^2-7\lambda+10=0.
\]
Factorising: $(\lambda-2)(\lambda-5)=0$, so $\lambda_1=2$, $\lambda_2=5$.
\textbf{Check:} $\operatorname{tr}(A)=7=2+5$ \checkmark; $\det A=12-2=10=2\times5$ \checkmark.

\textbf{Step 2: eigenvectors for $\lambda_1=2$.} Solve $(A-2I)\vec{x}=\vec{0}$:
\[
\begin{pmatrix} 2 & 2 \\ 1 & 1 \end{pmatrix}\begin{pmatrix}x\\y\end{pmatrix}=\vec{0}
\implies x+y=0.
\]
Eigenvectors: $\vec{v}=t\begin{pmatrix}1\\-1\end{pmatrix}$, $t\neq0$. Eigenspace $E_2=\operatorname{span}\{(1,-1)\}$.

\textbf{Step 3: eigenvectors for $\lambda_2=5$.} Solve $(A-5I)\vec{x}=\vec{0}$:
\[
\begin{pmatrix} -1 & 2 \\ 1 & -2 \end{pmatrix}\begin{pmatrix}x\\y\end{pmatrix}=\vec{0}
\implies x=2y.
\]
Eigenvectors: $\vec{v}=t\begin{pmatrix}2\\1\end{pmatrix}$, $t\neq0$. Eigenspace $E_5=\operatorname{span}\{(2,1)\}$.

\textbf{Step 4: verification.} $A\begin{pmatrix}1\\-1\end{pmatrix}=\begin{pmatrix}2\\-2\end{pmatrix}=2\begin{pmatrix}1\\-1\end{pmatrix}$ \checkmark; \quad $A\begin{pmatrix}2\\1\end{pmatrix}=\begin{pmatrix}10\\5\end{pmatrix}=5\begin{pmatrix}2\\1\end{pmatrix}$ \checkmark.

\textbf{Geometry.} The line $y=-x$ is invariant, stretched by factor $2$; the line $y=\tfrac{x}{2}$ is invariant, stretched by factor $5$.
\end{exoresolu}

\begin{popfigure}
\begin{tikzpicture}[scale=1.2]
\draw[popline] (-3,0) -- (3,0);
\draw[popline] (0,-3) -- (0,3);
% Eigenspace for lambda=2: line y=-x
\draw[poporange, thick, ->] (0,0) -- (2,-2) node[below right] {$\vec{v}_1=(1,-1)$};
\draw[poporange, thick, ->] (0,0) -- (-2,2) node[above left] {$-\vec{v}_1$};
\node[poporange] at (2.5,-1.5) {$E_2: y=-x$};
\node[poporange] at (2.5,-2) {stretch $\times 2$};
% Eigenspace for lambda=5: line y=x/2
\draw[popgreen, thick, ->] (0,0) -- (2,1) node[above right] {$\vec{v}_2=(2,1)$};
\draw[popgreen, thick, ->] (0,0) -- (-2,-1) node[below left] {$-\vec{v}_2$};
\node[popgreen] at (-2.5,1.5) {$E_5: y=x/2$};
\node[popgreen] at (-2.5,1) {stretch $\times 5$};
% Unit circle transformed to ellipse
\draw[poppurple, dashed] (0,0) ellipse (2 and 1);
\node[poppurple] at (0,-2.5) {Unit circle $\to$ ellipse};
\end{tikzpicture}
\legende{Eigenvectors define invariant lines; the transformation stretches by eigenvalues.}
\end{popfigure}

\courssub{Diagonalisation of $3\times3$ Matrices}

\begin{methode}[$3\times3$ case and diagonalisation]
\begin{enumerate}
\item Expand $\det(A-\lambda I)$ along a row or column (choose one with zeros). Look for an obvious root ($\lambda=0,1,2,-1,\dots$) and factorise.
\item Find a basis of each eigenspace by solving $(A-\lambda_i I)\vec{x}=\vec{0}$.
\item If three independent eigenvectors $\vec{v}_1,\vec{v}_2,\vec{v}_3$ exist, form $P=(\vec{v}_1\ \vec{v}_2\ \vec{v}_3)$; then $P^{-1}AP=D$ is diagonal with the eigenvalues on the diagonal, \emph{in the same order} as the columns of $P$.
\item \textbf{Application:} $A^n=PD^nP^{-1}$, which makes powers easy to compute.
\end{enumerate}
\begin{attention}[Order matters]
The eigenvalues in $D$ must appear in the same order as the corresponding eigenvectors in $P$. Swapping columns of $P$ without swapping entries of $D$ is a classic GCE error.
\end{attention}
\end{methode}

\begin{exoresolu}[Diagonalising a $3\times3$ matrix]
Let $A=\begin{pmatrix} 1 & 0 & 0 \\ 0 & 2 & 1 \\ 0 & 1 & 2 \end{pmatrix}$.
\begin{enumerate}
\item Find the eigenvalues of $A$.
\item Find a basis of each eigenspace.
\item Write down matrices $P$ and $D$ such that $A=PDP^{-1}$, and deduce $A^4$.
\end{enumerate}
\tcblower
\textbf{Solution.}
\begin{enumerate}
\item Expanding along the first row:
\[
\det(A-\lambda I)=(1-\lambda)\begin{vmatrix} 2-\lambda & 1 \\ 1 & 2-\lambda \end{vmatrix}
=(1-\lambda)\bigl[(2-\lambda)^2-1\bigr]=(1-\lambda)(\lambda^2-4\lambda+3).
\]
Since $\lambda^2-4\lambda+3=(\lambda-1)(\lambda-3)$, we get $\det(A-\lambda I)=(1-\lambda)^2(3-\lambda)\cdot(-1)^0$; precisely, the roots are $\lambda=1$ (double) and $\lambda=3$.
\textbf{Check:} $\operatorname{tr}(A)=5=1+1+3$ \checkmark.

\item \textbf{For $\lambda=1$:} $(A-I)\vec{x}=\vec{0}$:
\[
\begin{pmatrix} 0 & 0 & 0 \\ 0 & 1 & 1 \\ 0 & 1 & 1 \end{pmatrix}\vec{x}=\vec{0}\implies y+z=0,\ x\ \text{free}.
\]
Two free parameters ($x$ and $y$, with $z=-y$): $E_1=\operatorname{span}\{(1,0,0),\ (0,1,-1)\}$.

\textbf{For $\lambda=3$:} $(A-3I)\vec{x}=\vec{0}$:
\[
\begin{pmatrix} -2 & 0 & 0 \\ 0 & -1 & 1 \\ 0 & 1 & -1 \end{pmatrix}\vec{x}=\vec{0}\implies x=0,\ y=z.
\]
$E_3=\operatorname{span}\{(0,1,1)\}$.

\item We have three independent eigenvectors, so $A$ is diagonalisable:
\[
P=\begin{pmatrix} 1 & 0 & 0 \\ 0 & 1 & 1 \\ 0 & -1 & 1 \end{pmatrix},\qquad
D=\begin{pmatrix} 1 & 0 & 0 \\ 0 & 1 & 0 \\ 0 & 0 & 3 \end{pmatrix}.
\]
Since the first two columns of $P$ lie in $E_1$ (eigenvalue $1$), $A^4=PD^4P^{-1}$ with $D^4=\operatorname{diag}(1,1,81)$. The block structure gives directly: the $2\times2$ block $B=\begin{pmatrix}2&1\\1&2\end{pmatrix}$ has eigenvalues $1,3$, and one computes $B^4=\begin{pmatrix}41&40\\40&41\end{pmatrix}$ (using $B^n=\tfrac12\begin{pmatrix}3^n+1 & 3^n-1\\ 3^n-1 & 3^n+1\end{pmatrix}$ with $3^4=81$). Hence
\[
A^4=\begin{pmatrix} 1 & 0 & 0 \\ 0 & 41 & 40 \\ 0 & 40 & 41 \end{pmatrix}.
\]
\end{enumerate}
\end{exoresolu}

\begin{savaistu}[Diagonalisation in practice]
Diagonalisation is used to solve systems of linear differential equations, analyse stability in dynamical systems, and compute matrix powers efficiently. In structural engineering (e.g., designing bridges in Cameroon), the eigenvalues of stiffness matrices determine natural vibration frequencies — avoiding resonance is a matter of safety!
\end{savaistu}

\coursec{Integrated GCE-Style Problems}

\begin{exoresolu}[Comprehensive problem: kernel, image, eigenvalues]
The linear transformation $T:\mathbb{R}^3\to\mathbb{R}^3$ has matrix
\[
A=\begin{pmatrix} 2 & 1 & 1 \\ 1 & 2 & 1 \\ 1 & 1 & 2 \end{pmatrix}.
\]
\begin{enumerate}
\item Show that $\vec{v}_1=(1,1,1)$ is an eigenvector and find its eigenvalue.
\item Show that every vector of the plane $x+y+z=0$ is an eigenvector with eigenvalue $1$.
\item Deduce the eigenvalues of $A$ and a basis of eigenvectors.
\item Find $\ker(T-I)$ and state the set of invariant points of $T$.
\item Is $T$ invertible? Justify.
\end{enumerate}
\tcblower
\textbf{Solution.}
\begin{enumerate}
\item $A\vec{v}_1=\begin{pmatrix}2+1+1\\1+2+1\\1+1+2\end{pmatrix}=\begin{pmatrix}4\\4\\4\end{pmatrix}=4\vec{v}_1$. So $\vec{v}_1$ is an eigenvector with eigenvalue $\lambda=4$.

\item Let $\vec{x}=(x,y,z)$ with $x+y+z=0$. Then
\[
A\vec{x}=\begin{pmatrix}2x+y+z\\x+2y+z\\x+y+2z\end{pmatrix}
=\begin{pmatrix}x+(x+y+z)\\y+(x+y+z)\\z+(x+y+z)\end{pmatrix}
=\begin{pmatrix}x\\y\\z\end{pmatrix}=\vec{x}.
\]
Every vector of the plane $x+y+z=0$ is fixed by $A$, i.e. is an eigenvector with eigenvalue $\lambda=1$.

\item The plane $x+y+z=0$ has dimension $2$ (basis e.g. $(1,-1,0)$, $(1,0,-1)$) and $\vec{v}_1$ is not in it, so $\{(1,1,1),(1,-1,0),(1,0,-1)\}$ is a basis of $\mathbb{R}^3$ of eigenvectors. Eigenvalues: $1$ (double) and $4$. \textbf{Check:} $\operatorname{tr}(A)=6=1+1+4$ \checkmark.

\item $\ker(T-I)=E_1=\{(x,y,z):x+y+z=0\}$, the plane through the origin with normal $(1,1,1)$. Since $(T-I)\vec{x}=\vec{0}\iff T(\vec{x})=\vec{x}$, the \textbf{set of invariant points of $T$ is exactly this plane}.

\item $\det A=$ product of eigenvalues $=1\times1\times4=4\neq0$, so $A$ is invertible and $T$ \textbf{is invertible}. (Equivalently, $\ker T=\{\vec{0}\}$ since $0$ is not an eigenvalue.)
\end{enumerate}
\end{exoresolu}

\begin{exoresolu}[GCE 2022 style: Invariant lines and diagonalisation]
A linear transformation $T:\mathbb{R}^3\to\mathbb{R}^3$ is represented by the matrix
\[
M=\begin{pmatrix} 4 & -2 & 2 \\ -2 & 1 & -1 \\ 2 & -1 & 1 \end{pmatrix}.
\]
\begin{enumerate}
\item Find the eigenvalues of $M$.
\item For each eigenvalue, find a basis of the corresponding eigenspace.
\item Write down a matrix $P$ and a diagonal matrix $D$ such that $M=PDP^{-1}$.
\item Hence find $M^n$ for any positive integer $n$.
\item Determine the Cartesian equations of all invariant lines of $T$ through the origin.
\end{enumerate}
\tcblower
\textbf{Solution.}
\begin{enumerate}
\item $\det(M-\lambda I)=\begin{vmatrix} 4-\lambda & -2 & 2 \\ -2 & 1-\lambda & -1 \\ 2 & -1 & 1-\lambda \end{vmatrix}$.
Expand along first row:
\begin{align*}
&(4-\lambda)\bigl[(1-\lambda)^2-1\bigr]+2\bigl[-2(1-\lambda)+2\bigr]+2\bigl[2-2(1-\lambda)\bigr]\\
&=(4-\lambda)(\lambda^2-2\lambda)+2(2\lambda)+2(2\lambda)\\
&=(4-\lambda)\lambda(\lambda-2)+8\lambda=\lambda\bigl[(4-\lambda)(\lambda-2)+8\bigr]\\
&=\lambda\bigl[-\lambda^2+6\lambda-8+8\bigr]=\lambda(-\lambda^2+6\lambda)=-\lambda^2(\lambda-6).
\end{align*}
Eigenvalues: $\lambda=0$ (double), $\lambda=6$.

\item \textbf{For $\lambda=0$:} Solve $M\vec{x}=\vec{0}$:
\[
\begin{pmatrix} 4 & -2 & 2 \\ -2 & 1 & -1 \\ 2 & -1 & 1 \end{pmatrix}\sim\begin{pmatrix} 2 & -1 & 1 \\ 0 & 0 & 0 \\ 0 & 0 & 0 \end{pmatrix}\implies 2x-y+z=0.
\]
Two free parameters: $E_0=\operatorname{span}\{(1,2,0),\ (1,0,-2)\}$.

\textbf{For $\lambda=6$:} Solve $(M-6I)\vec{x}=\vec{0}$:
\[
\begin{pmatrix} -2 & -2 & 2 \\ -2 & -5 & -1 \\ 2 & -1 & -5 \end{pmatrix}\sim\begin{pmatrix} 1 & 1 & -1 \\ 0 & -3 & -3 \\ 0 & 0 & 0 \end{pmatrix}\implies x+y-z=0,\ y+z=0.
\]
Solution: $y=-z$, $x=2z$. $E_6=\operatorname{span}\{(2,-1,1)\}$.

\item $P=\begin{pmatrix} 1 & 1 & 2 \\ 2 & 0 & -1 \\ 0 & -2 & 1 \end{pmatrix}$, $D=\begin{pmatrix} 0 & 0 & 0 \\ 0 & 0 & 0 \\ 0 & 0 & 6 \end{pmatrix}$ (order matches columns of $P$).

\item $M^n=PD^nP^{-1}$ with $D^n=\operatorname{diag}(0,0,6^n)$. Compute $P^{-1}$:
\[
P^{-1}=\frac{1}{9}\begin{pmatrix} -2 & 5 & -1 \\ -2 & -1 & 5 \\ 4 & 2 & 2 \end{pmatrix}.
\]
Then $M^n=6^n\cdot\frac{1}{9}\begin{pmatrix} 2 \\ -1 \\ 1 \end{pmatrix}\begin{pmatrix} 4 & 2 & 2 \end{pmatrix}=\frac{6^n}{9}\begin{pmatrix} 8 & 4 & 4 \\ -4 & -2 & -2 \\ 4 & 2 & 2 \end{pmatrix}=\frac{2\cdot6^{n-1}}{3}\begin{pmatrix} 4 & 2 & 2 \\ -2 & -1 & -1 \\ 2 & 1 & 1 \end{pmatrix}$.

\item Invariant lines through origin = eigenspaces.
$\lambda=0$: plane $2x-y+z=0$ (every line in this plane is invariant, mapped to $\vec{0}$).
$\lambda=6$: line $\frac{x}{2}=\frac{y}{-1}=\frac{z}{1}$ (stretched by factor $6$).
\end{enumerate}
\end{exoresolu}

\begin{exercice}[Final practice: mixed topics]
Let $T:\mathbb{R}^3\to\mathbb{R}^3$ be defined by $T(x,y,z)=(2x+y,\;x+2y,\;3z)$.
\begin{enumerate}
\item Find the matrix $A$ of $T$.
\item Find $\ker T$ and $\operatorname{Im} T$. State their dimensions.
\item Find all eigenvalues and a basis for each eigenspace.
\item Is $T$ diagonalisable? If yes, find $P$ and $D$ such that $A=PDP^{-1}$.
\item Find the image of the line $\vec{r}=(1,1,1)+t(1,-1,0)$.
\item Find the invariant points of $T$.
\end{enumerate}
\end{exercice}

\begin{corrige}[Exercise 3]
\textbf{Solution.}
\begin{enumerate}
\item $A=\begin{pmatrix} 2 & 1 & 0 \\ 1 & 2 & 0 \\ 0 & 0 & 3 \end{pmatrix}$.
\item $A\vec{x}=\vec{0}\Rightarrow 2x+y=0,\;x+2y=0,\;3z=0\Rightarrow x=y=z=0$. $\ker T=\{\vec{0}\}$, $\dim=0$. $\operatorname{Im} T=\mathbb{R}^3$, $\dim=3$. Rank-nullity: $0+3=3$.
\item $\det(A-\lambda I)=(3-\lambda)\begin{vmatrix}2-\lambda&1\\1&2-\lambda\end{vmatrix}=(3-\lambda)(\lambda^2-4\lambda+3)=(3-\lambda)^2(1-\lambda)$.
Eigenvalues: $\lambda=3$ (double), $\lambda=1$.
$\lambda=3$: $(A-3I)\vec{x}=\vec{0}\Rightarrow -x+y=0,\;x-y=0,\;0=0\Rightarrow y=x,\;z$ free. $E_3=\operatorname{span}\{(1,1,0),\ (0,0,1)\}$.
$\lambda=1$: $(A-I)\vec{x}=\vec{0}\Rightarrow x+y=0,\;x+y=0,\;2z=0\Rightarrow y=-x,\;z=0$. $E_1=\operatorname{span}\{(1,-1,0)\}$.
\item Yes, three independent eigenvectors. $P=\begin{pmatrix}1&0&1\\1&0&-1\\0&1&0\end{pmatrix}$, $D=\begin{pmatrix}3&0&0\\0&3&0\\0&0&1\end{pmatrix}$.
\item $T(1,1,1)=(3,3,3)$, $T(1,-1,0)=(1,-1,0)$. Image: line $\vec{r}=(3,3,3)+t(1,-1,0)$.
\item Invariant points: $(A-I)\vec{x}=\vec{0}\Rightarrow x+y=0,\;2z=0\Rightarrow y=-x,\;z=0$. Line $\frac{x}{1}=\frac{y}{-1},\;z=0$ (the eigenspace $E_1$).
\end{enumerate}
\end{corrige}

\begin{aretenir}[Chapter summary: Linear Transformations]
\begin{itemize}
\item \textbf{Linearity:} $T(\lambda\vec{u}+\mu\vec{v})=\lambda T(\vec{u})+\mu T(\vec{v})$. Test $T(\vec{0})=\vec{0}$ first.
\item \textbf{Matrix:} Columns are $T(\vec{e}_i)$. $T(\vec{x})=A\vec{x}$.
\item \textbf{Kernel:} Solve $A\vec{x}=\vec{0}$. $\dim\ker T$ = nullity.
\item \textbf{Image:} Column space of $A$. $\dim\operatorname{Im} T$ = rank.
\item \textbf{Rank--Nullity:} $\rho(T)+\nu(T)=n$.
\item \textbf{Inverse:} Exists $\iff \det A\neq0 \iff \ker T=\{\vec{0}\}$. Matrix $A^{-1}$.
\item \textbf{Invariant points:} Solve $(A-I)\vec{x}=\vec{0}$.
\item \textbf{Images of lines/planes:} Apply $T$ to parametric form. Dimension never increases.
\item \textbf{Eigenvalues:} Roots of $\det(A-\lambda I)=0$.
\item \textbf{Eigenvectors:} Non-zero solutions of $(A-\lambda I)\vec{x}=\vec{0}$.
\item \textbf{Diagonalisation:} $A=PDP^{-1}$ if $n$ independent eigenvectors exist. $A^n=PD^nP^{-1}$.
\end{itemize}
\end{aretenir}

\newpage

\coursec{Exercises}

\courssub{I check my knowledge}

\begin{exercice}[Multiple choice questions]
For each question, choose the single correct answer, justifying briefly.
\begin{enumerate}
\item Which of the following mappings $T:\mathbb{R}^{2}\to\mathbb{R}^{2}$ is a linear transformation?
\begin{enumerate}
\item $T(x,y)=(2x-y,\;x+3y)$
\item $T(x,y)=(x^{2},\;y)$
\item $T(x,y)=(x+y,\;1)$
\item $T(x,y)=(xy,\;x-y)$
\end{enumerate}
\item The kernel of the linear transformation $T(x,y)=(x-y,\;2x-2y)$ is:
\begin{enumerate}
\item $\{(0,0)\}$
\item the line $y=x$
\item the line $y=2x$
\item the whole plane $\mathbb{R}^{2}$
\end{enumerate}
\item A linear transformation $T:\mathbb{R}^{2}\to\mathbb{R}^{2}$ is bijective if and only if:
\begin{enumerate}
\item $\ker T=\mathbb{R}^{2}$
\item $\operatorname{Im} T=\{(0,0)\}$
\item $\ker T=\{(0,0)\}$
\item $\det T=0$
\end{enumerate}
\item The eigenvalues of $A=\begin{pmatrix} 3 & 0 \\ 0 & 5 \end{pmatrix}$ are:
\begin{enumerate}
\item $3$ and $-5$
\item $3$ and $5$
\item $8$ and $0$
\item $15$ and $0$
\end{enumerate}
\end{enumerate}
\end{exercice}

\begin{exercice}[True or false]
State whether each assertion is true or false, justifying your answer.
\begin{enumerate}
\item Every linear transformation maps the zero vector to the zero vector.
\item If $T(u+v)=T(u)+T(v)$ for all vectors $u,v$, then $T$ is necessarily linear.
\item The kernel of a linear transformation $T:\mathbb{R}^{2}\to\mathbb{R}^{2}$ is always a line through the origin.
\item If $T$ is a linear transformation of the plane with $\det T\neq 0$, then the image of any line is a line.
\item An eigenvector of a linear transformation can never be the zero vector.
\end{enumerate}
\end{exercice}

\begin{exercice}[Key definitions]
\begin{enumerate}
\item Define a \cle{linear transformation} from $\mathbb{R}^{2}$ to $\mathbb{R}^{2}$.
\item Define the \cle{kernel} and the \cle{image} of a linear transformation $T$.
\item Define an \cle{eigenvalue} and an \cle{eigenvector} of a linear transformation.
\item What is an \cle{invariant point} (fixed point) of a transformation?
\end{enumerate}
\end{exercice}

\begin{exercice}[Direct calculations]
\begin{enumerate}
\item Given $T(x,y)=(3x+y,\;x-2y)$, compute $T(1,2)$, $T(-1,3)$ and $T(0,0)$.
\item Write down the matrix of $T(x,y)=(3x+y,\;x-2y)$ and compute its determinant.
\item A linear transformation of the plane maps $(1,0)$ to $(3,-1)$ and $(0,1)$ to $(2,4)$. Write down its matrix and find the image of $(5,2)$.
\item Compute the characteristic polynomial of $A=\begin{pmatrix} 2 & 1 \\ 1 & 2 \end{pmatrix}$.
\end{enumerate}
\end{exercice}

\courssub{I practise}

\begin{exercice}[Testing linearity]
Determine whether each of the following mappings is a linear transformation, justifying your answer carefully:
\begin{enumerate}
\item $T(x,y)=(2x-y,\;x+3y)$;
\item $S(x,y)=(x^{2},\;y)$;
\item $U(x,y)=(x+y,\;1)$;
\item $V(x,y)=(x-2y,\;0,\;3x+y)$ from $\mathbb{R}^{2}$ to $\mathbb{R}^{3}$.
\end{enumerate}
\end{exercice}

\begin{exercice}[Kernel and image]
Let $T(x,y)=(x+2y,\;2x+4y)$.
\begin{enumerate}
\item Find $\ker T$ and describe it geometrically.
\item Find $\operatorname{Im} T$ and describe it geometrically.
\item State, with reasons, whether $T$ is injective and whether $T$ is surjective.
\item Verify that $\dim(\ker T)+\dim(\operatorname{Im} T)=2$.
\end{enumerate}
\end{exercice}

\begin{exercice}[Inverse and invariant points]
Consider the transformation with matrix $A=\begin{pmatrix} 2 & 1 \\ 1 & 2 \end{pmatrix}$.
\begin{enumerate}
\item Show that the transformation is bijective.
\item Find its inverse transformation, i.e. express $T^{-1}(x',y')$ in terms of $x'$ and $y'$.
\item Determine all the invariant points of $T$.
\end{enumerate}
\end{exercice}

\begin{exercice}[Image of a line]
Find the image of the line $2x-y+3=0$ under the transformation $T(x,y)=(x-y,\;2x+y)$, giving a Cartesian equation of the image line. Verify your answer by checking that the images of two points of the original line satisfy your equation.
\end{exercice}

\courssub{I prepare for the GCE}

\begin{exercice}[GCE style: a shear transformation] \hfill (12 marks)\\
$T$ is the linear transformation of the plane with matrix $M=\begin{pmatrix} 1 & 2 \\ 0 & 1 \end{pmatrix}$.
\begin{enumerate}
\item Identify $T$ geometrically, stating precisely the direction and the factor of the transformation. \hfill (3 marks)
\item Find the images of the vertices of the unit square with vertices $(0,0)$, $(1,0)$, $(1,1)$, $(0,1)$, sketch the image figure, and compute its area. \hfill (4 marks)
\item Find all the fixed points of $T$. \hfill (2 marks)
\item Find all the eigenvalues of $M$ and the corresponding eigenvectors. Comment on the link with part (c). \hfill (3 marks)
\end{enumerate}
\end{exercice}

\begin{exercice}[GCE style: transformation of space] \hfill (14 marks)\\
In space, the linear transformation $T$ is defined by
\[ T(x,y,z)=(x+y,\;y+z,\;x+2y+z). \]
\begin{enumerate}
\item Write down the matrix of $T$ and compute its determinant. \hfill (3 marks)
\item Find the kernel of $T$ and describe it geometrically. \hfill (3 marks)
\item Find the image of $T$ and show that it is a plane through the origin, giving a Cartesian equation of this plane. \hfill (4 marks)
\item Determine the image of the plane $z=0$ under $T$. \hfill (2 marks)
\item Is $T$ bijective? Justify your answer in two different ways. \hfill (2 marks)
\end{enumerate}
\end{exercice}

\begin{exercice}[GCE style: eigenvalues, powers and iteration] \hfill (16 marks)\\
Let $A=\begin{pmatrix} 2 & 1 \\ 1 & 2 \end{pmatrix}$ and let $T$ be the associated linear transformation of the plane.
\begin{enumerate}
\item Find the eigenvalues $\lambda_{1},\lambda_{2}$ of $A$ and a corresponding eigenvector for each. \hfill (4 marks)
\item A transformation of the plane has eigenvalues $3$ and $-1$ with eigenvectors $(1,1)$ and $(1,-1)$ respectively. By first expressing $(4,0)$ as a linear combination of the eigenvectors, find the image of $(4,0)$. \hfill (4 marks)
\item Returning to the matrix $A$ above, conjecture a formula for $A^{n}$ of the form
\[ A^{n}=\frac{1}{2}\begin{pmatrix} 3^{n}+1 & 3^{n}-1 \\ 3^{n}-1 & 3^{n}+1 \end{pmatrix}, \]
verify it for $n=1$ and $n=2$, and prove it by induction on $n$. \hfill (5 marks)
\item Hence compute the image of the vector $(1,0)$ when $T$ is applied ten times successively. \hfill (3 marks)
\end{enumerate}
\end{exercice}

\newpage

\coursec{Solutions to the exercises}

\begin{corrige}[Exercise 1 — Multiple choice questions]
\begin{enumerate}
\item \textbf{Answer (a).} $T(x,y)=(2x-y,\;x+3y)$ has the form $(ax+by,\;cx+dy)$, hence it is linear (its matrix is $\begin{pmatrix} 2 & -1 \\ 1 & 3 \end{pmatrix}$). The others fail: (b) contains $x^{2}$ (e.g. $T(2,0)=(4,0)\neq 2T(1,0)=(2,0)$); (c) gives $T(0,0)=(0,1)\neq(0,0)$; (d) contains the product $xy$ (e.g. $T(2,2)=(4,0)\neq 2T(1,1)=(2,0)$).
\item \textbf{Answer (b).} $T(x,y)=(0,0)\iff x-y=0\iff y=x$, a line through the origin.
\item \textbf{Answer (c).} For a linear transformation between spaces of the same finite dimension, bijectivity is equivalent to injectivity, which is equivalent to $\ker T=\{(0,0)\}$ (equivalently $\det T\neq 0$, not $0$).
\item \textbf{Answer (b).} A diagonal matrix has its diagonal entries as eigenvalues: $\lambda=3$ and $\lambda=5$ (the characteristic polynomial is $(3-\lambda)(5-\lambda)$).
\end{enumerate}
\end{corrige}

\begin{corrige}[Exercise 2 — True or false]
\begin{enumerate}
\item \textbf{True.} $T(0)=T(0+0)=T(0)+T(0)$, and subtracting $T(0)$ from both sides gives $T(0)=0$.
\item \textbf{False.} Additivity alone does not imply homogeneity $T(\lambda u)=\lambda T(u)$. For example, on $\mathbb{R}$ viewed as a vector space over $\mathbb{Q}$, additive non-linear maps exist; more simply at this level, one must \emph{also} check $T(\lambda u)=\lambda T(u)$ — the two conditions together define linearity.
\item \textbf{False.} The kernel is always a \emph{vector subspace}, but it can be $\{(0,0)\}$ (e.g. the identity), a line through the origin, or the whole plane $\mathbb{R}^{2}$ (the zero transformation). It is not always a line.
\item \textbf{True.} If $\det T\neq 0$, $T$ is bijective and maps lines to lines: the image of a line with direction vector $u$ passing through $A$ is the line through $T(A)$ with direction vector $T(u)\neq 0$.
\item \textbf{True.} By definition an eigenvector is a \emph{non-zero} vector $u$ such that $T(u)=\lambda u$. The zero vector is excluded (otherwise every scalar would be an eigenvalue).
\end{enumerate}
\end{corrige}

\begin{corrige}[Exercise 3 — Key definitions]
\begin{enumerate}
\item A mapping $T:\mathbb{R}^{2}\to\mathbb{R}^{2}$ is a \cle{linear transformation} if for all vectors $u,v\in\mathbb{R}^{2}$ and every scalar $\lambda\in\mathbb{R}$:
\[ T(u+v)=T(u)+T(v)\qquad\text{and}\qquad T(\lambda u)=\lambda\,T(u). \]
Equivalently, $T(\lambda u+\mu v)=\lambda T(u)+\mu T(v)$ for all $u,v$ and all scalars $\lambda,\mu$.
\item The \cle{kernel} of $T$ is $\ker T=\{\,u\in\mathbb{R}^{2}\;:\;T(u)=0\,\}$, the set of vectors mapped to the zero vector. The \cle{image} of $T$ is $\operatorname{Im}T=\{\,T(u)\;:\;u\in\mathbb{R}^{2}\,\}$, the set of all values taken by $T$. Both are vector subspaces.
\item A scalar $\lambda$ is an \cle{eigenvalue} of $T$ if there exists a non-zero vector $u$ such that $T(u)=\lambda u$; such a vector $u$ is then called an \cle{eigenvector} associated with $\lambda$.
\item An \cle{invariant point} (fixed point) of a transformation $T$ is a point $M$ such that $T(M)=M$: the point is mapped onto itself.
\end{enumerate}
\end{corrige}

\begin{corrige}[Exercise 4 — Direct calculations]
\begin{enumerate}
\item $T(1,2)=(3(1)+2,\;1-2(2))=(5,-3)$; \quad $T(-1,3)=(3(-1)+3,\;-1-2(3))=(0,-7)$; \quad $T(0,0)=(0,0)$, as expected for a linear transformation.
\item The matrix is read from the columns $T(1,0)=(3,1)$ and $T(0,1)=(1,-2)$:
\[ A=\begin{pmatrix} 3 & 1 \\ 1 & -2 \end{pmatrix},\qquad \det A=(3)(-2)-(1)(1)=-7. \]
\item The columns of the matrix are the images of the basis vectors:
\[ A=\begin{pmatrix} 3 & 2 \\ -1 & 4 \end{pmatrix},\qquad T(5,2)=\begin{pmatrix} 3 & 2 \\ -1 & 4 \end{pmatrix}\begin{pmatrix} 5 \\ 2 \end{pmatrix}=\begin{pmatrix} 15+4 \\ -5+8 \end{pmatrix}=\begin{pmatrix} 19 \\ 3 \end{pmatrix}. \]
So $T(5,2)=(19,3)$.
\item The characteristic polynomial is
\[ \det(A-\lambda I)=\begin{vmatrix} 2-\lambda & 1 \\ 1 & 2-\lambda \end{vmatrix}=(2-\lambda)^{2}-1=\lambda^{2}-4\lambda+3. \]
\end{enumerate}
\end{corrige}

\begin{corrige}[Exercise 5 — Testing linearity]
\begin{enumerate}
\item \textbf{$T$ is linear.} Let $u=(x_{1},y_{1})$, $v=(x_{2},y_{2})$ and $\lambda\in\mathbb{R}$. Then
\[ T(u+v)=(2(x_{1}+x_{2})-(y_{1}+y_{2}),\;(x_{1}+x_{2})+3(y_{1}+y_{2}))=T(u)+T(v), \]
and $T(\lambda u)=(2\lambda x_{1}-\lambda y_{1},\;\lambda x_{1}+3\lambda y_{1})=\lambda T(u)$. Both axioms hold (equivalently, $T$ has the form $(ax+by,cx+dy)$).
\item \textbf{$S$ is not linear.} Counter-example to homogeneity: $S(2,0)=(4,0)$ but $2S(1,0)=2(1,0)=(2,0)$. Since $S(2u)\neq 2S(u)$ for $u=(1,0)$, $S$ fails.
\item \textbf{$U$ is not linear.} $U(0,0)=(0,1)\neq(0,0)$, whereas every linear transformation maps the zero vector to the zero vector.
\item \textbf{$V$ is linear.} Each component is a linear form in $x$ and $y$: $V(x,y)=(x-2y,\;0,\;3x+y)$. Checking: $V(u+v)=V(u)+V(v)$ and $V(\lambda u)=\lambda V(u)$ follow from distributivity of scalar multiplication in each coordinate. Its matrix is $\begin{pmatrix} 1 & -2 \\ 0 & 0 \\ 3 & 1 \end{pmatrix}$.
\end{enumerate}
\end{corrige}

\begin{corrige}[Exercise 6 — Kernel and image]
\begin{enumerate}
\item $(x,y)\in\ker T\iff \begin{cases} x+2y=0 \\ 2x+4y=0 \end{cases}\iff x=-2y$. Hence
\[ \ker T=\{\,(-2t,\;t)\;:\;t\in\mathbb{R}\,\}, \]
the \textbf{line through the origin} with direction vector $(-2,1)$ (equation $x+2y=0$).
\item $T(x,y)=(x+2y,\;2(x+2y))$: the second coordinate is always twice the first. Conversely, for any $s$, $(s,2s)=T(s,0)$. Hence
\[ \operatorname{Im}T=\{\,(s,2s)\;:\;s\in\mathbb{R}\,\}, \]
the \textbf{line through the origin} of equation $y=2x$, spanned by $(1,2)$.
\item $T$ is \textbf{not injective} because $\ker T\neq\{(0,0)\}$ (for instance $T(-2,1)=(0,0)=T(0,0)$). $T$ is \textbf{not surjective} because $\operatorname{Im}T\neq\mathbb{R}^{2}$ (for instance $(1,0)$ has no pre-image, since it would require $2x+4y=0$ and $x+2y=1$ simultaneously).
\item $\ker T$ is a line, so $\dim(\ker T)=1$; $\operatorname{Im}T$ is a line, so $\dim(\operatorname{Im}T)=1$. Then
\[ \dim(\ker T)+\dim(\operatorname{Im}T)=1+1=2=\dim\mathbb{R}^{2}, \]
confirming the rank–nullity theorem.
\end{enumerate}
\end{corrige}

\begin{corrige}[Exercise 7 — Inverse and invariant points]
\begin{enumerate}
\item $\det A=(2)(2)-(1)(1)=3\neq 0$, so $A$ is invertible and the transformation $T$ is bijective.
\item $A^{-1}=\dfrac{1}{3}\begin{pmatrix} 2 & -1 \\ -1 & 2 \end{pmatrix}$. Therefore, writing $(x',y')=T(x,y)$,
\[ T^{-1}(x',y')=\frac{1}{3}\begin{pmatrix} 2 & -1 \\ -1 & 2 \end{pmatrix}\begin{pmatrix} x' \\ y' \end{pmatrix}=\left(\frac{2x'-y'}{3},\;\frac{-x'+2y'}{3}\right). \]
Check: $T^{-1}(T(x,y))=\frac{1}{3}\big(2(2x+y)-(x+2y),\;-(2x+y)+2(x+2y)\big)=(x,y)$. \checkmark
\item Invariant points satisfy $T(x,y)=(x,y)$, i.e.
\[ \begin{cases} 2x+y=x \\ x+2y=y \end{cases}\iff \begin{cases} x+y=0 \\ x+y=0 \end{cases}\iff y=-x. \]
The set of invariant points is the \textbf{line $y=-x$} through the origin: every point of this line is fixed by $T$.
\end{enumerate}
\end{corrige}

\begin{corrige}[Exercise 8 — Image of a line]
Let $(x',y')=T(x,y)=(x-y,\;2x+y)$. Since $\det T=1\cdot 1-(-1)\cdot 2=3\neq 0$, $T$ is invertible and we can express $x,y$ in terms of $x',y'$:
\[ \begin{cases} x'=x-y \\ y'=2x+y \end{cases}\implies x=\frac{x'+y'}{3},\qquad y=\frac{y'-2x'}{3}. \]
Substituting into $2x-y+3=0$:
\[ 2\left(\frac{x'+y'}{3}\right)-\left(\frac{y'-2x'}{3}\right)+3=0 \iff \frac{2x'+2y'-y'+2x'}{3}+3=0 \iff 4x'+y'+9=0. \]
Hence the image is the line $\boxed{4x+y+9=0}$.

\textbf{Verification.} Choose two points of $2x-y+3=0$: $A=(0,3)$ and $B=(-1,1)$.
\[ T(A)=(0-3,\;0+3)=(-3,3):\quad 4(-3)+3+9=0\ \checkmark \]
\[ T(B)=(-1-1,\;-2+1)=(-2,-1):\quad 4(-2)-1+9=0\ \checkmark \]
Both images satisfy $4x+y+9=0$, confirming the result.
\end{corrige}

\begin{corrige}[Exercise 9 — GCE style: a shear transformation]
\begin{enumerate}
\item $T(x,y)=(x+2y,\;y)$. This is a \textbf{shear parallel to the $x$-axis} (direction: the horizontal axis, along the vector $(1,0)$), with \textbf{factor (shear ratio) $k=2$}: each point is displaced horizontally by $2y$, proportional to its distance from the $x$-axis. Points of the $x$-axis are fixed. \hfill [3 marks: direction 1, factor 1, description 1]
\item Images of the vertices:
\[ (0,0)\mapsto(0,0),\quad (1,0)\mapsto(1,0),\quad (1,1)\mapsto(3,1),\quad (0,1)\mapsto(2,1). \]
The image is a parallelogram with base $1$ (from $(0,0)$ to $(1,0)$) and height $1$, so its area is $1\times 1=1$ square unit. (Equivalently, $|\det M|=1$, so areas are preserved.)
\begin{popfigure}
\begin{tikzpicture}[scale=1.1]
\draw[->,gray] (-0.5,0) -- (3.8,0) node[right] {$x$};
\draw[->,gray] (0,-0.5) -- (0,1.8) node[above] {$y$};
\draw[popblue,thick] (0,0) -- (1,0) -- (1,1) -- (0,1) -- cycle;
\draw[poporange,very thick] (0,0) -- (1,0) -- (3,1) -- (2,1) -- cycle;
\node[popblue] at (0.5,0.5) {square};
\node[poporange] at (2.1,0.5) {image};
\foreach \p/\n in {(0,0)/O,(1,0)/A,(3,1)/B',(2,1)/C'} \fill \p circle (1.2pt);
\node[below left] at (0,0) {$O$}; \node[below] at (1,0) {$(1,0)$};
\node[above right] at (3,1) {$(3,1)$}; \node[above left] at (2,1) {$(2,1)$};
\end{tikzpicture}
\legende{The unit square (blue) and its image under the shear (orange).}
\end{popfigure}
\hfill [4 marks: images 2, sketch 1, area 1]
\item Fixed points: $T(x,y)=(x,y)\iff x+2y=x\iff y=0$. The set of fixed points is the \textbf{$x$-axis} (the line $y=0$). \hfill [2 marks]
\item Characteristic polynomial: $\det(M-\lambda I)=(1-\lambda)^{2}=0$, so the only eigenvalue is $\lambda=1$ (double). Eigenvectors: $(M-I)\binom{x}{y}=0\iff 2y=0\iff y=0$: the eigenvectors are the non-zero multiples of $(1,0)$.

\textbf{Link with (c):} eigenvectors with eigenvalue $1$ are exactly the non-zero fixed directions; since every point of the $x$-axis is fixed, the $x$-axis is both the line of fixed points and the eigenspace for $\lambda=1$. \hfill [3 marks: eigenvalue 1, eigenvectors 1, comment 1]
\end{enumerate}
\end{corrige}

\begin{corrige}[Exercise 10 — GCE style: transformation of space]
\begin{enumerate}
\item Reading the columns $T(1,0,0)=(1,0,1)$, $T(0,1,0)=(1,1,2)$, $T(0,0,1)=(0,1,1)$:
\[ A=\begin{pmatrix} 1 & 1 & 0 \\ 0 & 1 & 1 \\ 1 & 2 & 1 \end{pmatrix},\qquad \det A=1\begin{vmatrix}1&1\\2&1\end{vmatrix}-1\begin{vmatrix}0&1\\1&1\end{vmatrix}=1(1-2)-1(0-1)=-1+1=0. \]
\hfill [3 marks: matrix 2, determinant 1]
\item $(x,y,z)\in\ker T\iff \begin{cases} x+y=0 \\ y+z=0 \\ x+2y+z=0 \end{cases}$. The first two give $x=-y$, $z=-y$; the third is then $-y+2y-y=0$, automatically satisfied. Hence
\[ \ker T=\{\,(-t,\;t,\;-t)\;:\;t\in\mathbb{R}\,\}, \]
the \textbf{line through the origin} with direction vector $(-1,1,-1)$. \hfill [3 marks]
\item $T(x,y,z)=(x+y,\;y+z,\;x+2y+z)$. Note that the third coordinate equals the sum of the first two: $(x+y)+(y+z)=x+2y+z$. Hence every image vector $(X,Y,Z)$ satisfies $Z=X+Y$. Conversely, given $(X,Y,Z)$ with $Z=X+Y$, take $(x,y,z)=(X,0,Y)$: then $T(X,0,Y)=(X,Y,X+Y)=(X,Y,Z)$. Therefore
\[ \operatorname{Im}T=\{\,(X,Y,Z)\;:\;X+Y-Z=0\,\}, \]
the \textbf{plane through the origin} with Cartesian equation $x+y-z=0$ (normal vector $(1,1,-1)$). \hfill [4 marks]
\item Points of the plane $z=0$ are $(x,y,0)$; their images are $T(x,y,0)=(x+y,\;y,\;x+2y)$. As $(x,y)$ ranges over $\mathbb{R}^{2}$, the parameters $u=x+y$ and $v=y$ range over all of $\mathbb{R}^{2}$ (the map $(x,y)\mapsto(u,v)$ is invertible: $y=v$, $x=u-v$), and the third coordinate is $x+2y=u+v$. Hence the image is
\[ \{\,(u,v,u+v)\;:\;u,v\in\mathbb{R}\,\}, \]
i.e. the \textbf{plane $z=x+y$}, which is exactly $\operatorname{Im}T$ itself. \hfill [2 marks]
\item \textbf{$T$ is not bijective.}
\begin{itemize}
\item \emph{First way:} $\det A=0$, so $A$ is not invertible, so $T$ is not bijective.
\item \emph{Second way:} $\ker T\neq\{(0,0,0)\}$ (it is a whole line), so $T$ is not injective — for instance $T(-1,1,-1)=(0,0,0)=T(0,0,0)$ — hence not bijective. (Equivalently, $\operatorname{Im}T\neq\mathbb{R}^{3}$, so $T$ is not surjective.)
\end{itemize}
\hfill [2 marks: 1 per justification]
\end{enumerate}
\end{corrige}

\begin{corrige}[Exercise 11 — GCE style: eigenvalues, powers and iteration]
\begin{enumerate}
\item Characteristic polynomial: $\det(A-\lambda I)=(2-\lambda)^{2}-1=\lambda^{2}-4\lambda+3=(\lambda-1)(\lambda-3)$, so $\lambda_{1}=1$, $\lambda_{2}=3$.

For $\lambda=1$: $(A-I)\binom{x}{y}=0\iff x+y=0$, eigenvector $u_{1}=(1,-1)$.

For $\lambda=3$: $(A-3I)\binom{x}{y}=0\iff -x+y=0$, eigenvector $u_{2}=(1,1)$. \hfill [4 marks: polynomial 1, eigenvalues 1, eigenvectors 2]
\item Write $(4,0)=\alpha(1,1)+\beta(1,-1)$: then $\alpha+\beta=4$ and $\alpha-\beta=0$, giving $\alpha=\beta=2$. So $(4,0)=2(1,1)+2(1,-1)$. By linearity and the eigenvector property $T(u)=\lambda u$:
\[ T(4,0)=2\,T(1,1)+2\,T(1,-1)=2\cdot 3(1,1)+2\cdot 1(1,-1)=(6,6)+(2,-2)=(8,4). \]
\hfill [4 marks: decomposition 2, use of eigenvalues 1, answer 1]
\item \textbf{Verification.} For $n=1$: $\frac{1}{2}\begin{pmatrix} 3+1 & 3-1 \\ 3-1 & 3+1 \end{pmatrix}=\begin{pmatrix} 2 & 1 \\ 1 & 2 \end{pmatrix}=A$. \checkmark

For $n=2$: $A^{2}=\begin{pmatrix} 2 & 1 \\ 1 & 2 \end{pmatrix}\begin{pmatrix} 2 & 1 \\ 1 & 2 \end{pmatrix}=\begin{pmatrix} 5 & 4 \\ 4 & 5 \end{pmatrix}$, and the formula gives $\frac{1}{2}\begin{pmatrix} 9+1 & 9-1 \\ 9-1 & 9+1 \end{pmatrix}=\begin{pmatrix} 5 & 4 \\ 4 & 5 \end{pmatrix}$. \checkmark

\textbf{Induction.} Let $P(n)$ be the stated formula. $P(1)$ holds (shown above). Assume $P(n)$ for some $n\geq 1$. Then
\begin{align*}
A^{n+1}&=A^{n}A=\frac{1}{2}\begin{pmatrix} 3^{n}+1 & 3^{n}-1 \\ 3^{n}-1 & 3^{n}+1 \end{pmatrix}\begin{pmatrix} 2 & 1 \\ 1 & 2 \end{pmatrix}\\
&=\frac{1}{2}\begin{pmatrix} 2(3^{n}+1)+(3^{n}-1) & (3^{n}+1)+2(3^{n}-1) \\ 2(3^{n}-1)+(3^{n}+1) & (3^{n}-1)+2(3^{n}+1) \end{pmatrix}\\
&=\frac{1}{2}\begin{pmatrix} 3\cdot 3^{n}+1 & 3\cdot 3^{n}-1 \\ 3\cdot 3^{n}-1 & 3\cdot 3^{n}+1 \end{pmatrix}
=\frac{1}{2}\begin{pmatrix} 3^{n+1}+1 & 3^{n+1}-1 \\ 3^{n+1}-1 & 3^{n+1}+1 \end{pmatrix},
\end{align*}
which is $P(n+1)$. By the principle of mathematical induction, the formula holds for all $n\geq 1$. \hfill [5 marks: checks 2, induction step 2, conclusion 1]
\item Applying $T$ ten times to $(1,0)$ gives $A^{10}\binom{1}{0}$, the first column of $A^{10}$:
\[ A^{10}\begin{pmatrix} 1 \\ 0 \end{pmatrix}=\frac{1}{2}\begin{pmatrix} 3^{10}+1 \\ 3^{10}-1 \end{pmatrix}. \]
Since $3^{10}=59049$:
\[ T^{10}(1,0)=\left(\frac{59050}{2},\;\frac{59048}{2}\right)=(29525,\;29524). \]
\hfill [3 marks: method 1, value of $3^{10}$ 1, answer 1]
\end{enumerate}
\end{corrige}

\newpage

\coursec{Summary sheet}

\begin{popfigure}
\begin{center}\tcbox[colback=popdark,colframe=popdark,coltext=white,arc=9pt,boxsep=4pt,left=6pt,right=6pt]{\bfseries\small Linear Transformations}\end{center}
\vspace{-2pt}
\begin{tcbraster}[raster columns=3,raster equal height=rows,raster column skip=6pt,raster row skip=6pt,enhanced,blankest]
\begin{tcolorbox}[enhanced,colback=popblue,colframe=popblue,coltext=white,boxrule=0.9pt,arc=3pt,left=4pt,right=4pt,top=3pt,bottom=3pt,boxsep=1pt,halign=left]\bfseries\scriptsize Definition \& Properties\end{tcolorbox}
\begin{tcolorbox}[enhanced,colback=popgreen,colframe=popgreen,coltext=white,boxrule=0.9pt,arc=3pt,left=4pt,right=4pt,top=3pt,bottom=3pt,boxsep=1pt,halign=left]\bfseries\scriptsize Kernel \& Image\end{tcolorbox}
\begin{tcolorbox}[enhanced,colback=poporange,colframe=poporange,coltext=white,boxrule=0.9pt,arc=3pt,left=4pt,right=4pt,top=3pt,bottom=3pt,boxsep=1pt,halign=left]\bfseries\scriptsize Inverse \& Invariant Points\end{tcolorbox}
\begin{tcolorbox}[enhanced,colback=poppurple,colframe=poppurple,coltext=white,boxrule=0.9pt,arc=3pt,left=4pt,right=4pt,top=3pt,bottom=3pt,boxsep=1pt,halign=left]\bfseries\scriptsize Images of Lines \& Planes\end{tcolorbox}
\begin{tcolorbox}[enhanced,colback=poppink,colframe=poppink,coltext=white,boxrule=0.9pt,arc=3pt,left=4pt,right=4pt,top=3pt,bottom=3pt,boxsep=1pt,halign=left]\bfseries\scriptsize Eigenvalues \& Eigenvectors\end{tcolorbox}
\end{tcbraster}

\legende{Mind map of the chapter: the five big ideas around linear transformations.}
\end{popfigure}

\begin{aretenir}[Summary Sheet — Linear Transformations]

\textbf{Key definitions}
\begin{itemize}
\item A map $T:\mathbb{R}^{n}\to\mathbb{R}^{m}$ is a \cle{linear transformation} if for all vectors $\vec{u},\vec{v}\in\mathbb{R}^{n}$ and all scalars $\lambda,\mu\in\mathbb{R}$: $T(\lambda\vec{u}+\mu\vec{v})=\lambda T(\vec{u})+\mu T(\vec{v})$. Equivalently, $T(\vec{u}+\vec{v})=T(\vec{u})+T(\vec{v})$ and $T(\lambda\vec{u})=\lambda T(\vec{u})$.
\item Every linear transformation has a unique \cle{matrix} $A$ (relative to the standard bases) such that $T(\vec{x})=A\vec{x}$; the $j$-th column of $A$ is $T(\vec{e}_j)$, the image of the $j$-th standard basis vector.
\item \cle{Kernel}: $\ker T=\{\vec{x}\in\mathbb{R}^{n}:T(\vec{x})=\vec{0}\}$. \cle{Image}: $\operatorname{Im}T=\{T(\vec{x}):\vec{x}\in\mathbb{R}^{n}\}=\{A\vec{x}:\vec{x}\in\mathbb{R}^{n}\}$ (the column space of $A$).
\item $T$ is \cle{invertible} (bijective) iff $\ker T=\{\vec{0}\}$ iff $\det A\neq 0$; then $T^{-1}$ exists and has matrix $A^{-1}$.
\item \cle{Invariant (fixed) points}: points $\vec{x}$ such that $T(\vec{x})=\vec{x}$, i.e. $(A-I)\vec{x}=\vec{0}$.
\item \cle{Eigenvalue} $\lambda$ and \cle{eigenvector} $\vec{v}\neq\vec{0}$: $A\vec{v}=\lambda\vec{v}$, found from the \cle{characteristic equation} $\det(A-\lambda I)=0$.
\end{itemize}

\textbf{Formulas and theorems to know}
\begin{itemize}
\item $T(\vec{0})=\vec{0}$ always; if $T(\vec{0})\neq\vec{0}$ the map is \emph{not} linear.
\item Composition: if $T:\mathbb{R}^{n}\to\mathbb{R}^{m}$ has matrix $A$ and $S:\mathbb{R}^{m}\to\mathbb{R}^{p}$ has matrix $B$, then $S\circ T$ has matrix $BA$ (matrix of $S$ times matrix of $T$) — order matters.
\item Kernel and image are \cle{subspaces} of $\mathbb{R}^{n}$ and $\mathbb{R}^{m}$ respectively; $\ker T=\{\vec{0}\}\iff T$ is one-to-one (injective).
\item Rank--nullity theorem: $\dim(\ker T)+\dim(\operatorname{Im}T)=n$ (dimension of the domain).
\item For $A=\begin{pmatrix}a&b\\c&d\end{pmatrix}$: $\det A=ad-bc$, $A^{-1}=\dfrac{1}{\det A}\begin{pmatrix}d&-b\\-c&a\end{pmatrix}$ (provided $\det A\neq0$).
\item Characteristic polynomial ($2\times2$): $\lambda^{2}-(\operatorname{tr}A)\lambda+\det A=0$, where $\operatorname{tr}A=a+d$.
\item Eigenvectors for $\lambda$: solve $(A-\lambda I)\vec{v}=\vec{0}$; the set of solutions (including $\vec{0}$) forms the \cle{eigenspace} $E_{\lambda}$.
\item Image of a line through $\vec{a}$ with direction $\vec{d}$: line through $T(\vec{a})$ with direction $T(\vec{d})$ (if $T(\vec{d})\neq\vec{0}$); if $T(\vec{d})=\vec{0}$ the line collapses to the point $T(\vec{a})$.
\item Image of a plane $\vec{r}=\vec{a}+\lambda\vec{u}+\mu\vec{v}$: plane through $T(\vec{a})$ spanned by $T(\vec{u}),T(\vec{v})$ (if they are linearly independent); otherwise a line or a point.
\end{itemize}

\textbf{Methods}
\begin{enumerate}
\item \textbf{Test linearity:} check additivity and homogeneity, or the combined condition $T(\lambda\vec{u}+\mu\vec{v})=\lambda T(\vec{u})+\mu T(\vec{v})$; one counterexample disproves linearity.
\item \textbf{Find the matrix:} compute $T(\vec{e}_1),T(\vec{e}_2),\dots$ and write them as columns of $A$.
\item \textbf{Kernel:} solve $A\vec{x}=\vec{0}$ by Gaussian elimination; give a basis. \textbf{Image:} the column space of $A$; reduce columns to a linearly independent set (pivot columns) to obtain a basis.
\item \textbf{Inverse:} verify $\det A\neq 0$, compute $A^{-1}$ (adjugate/determinant or row reduction), then $T^{-1}(\vec{x})=A^{-1}\vec{x}$.
\item \textbf{Invariant points:} solve $(A-I)\vec{x}=\vec{0}$; interpret geometrically (fixed point, fixed line, whole space fixed if $A=I$).
\item \textbf{Eigenvalues/eigenvectors:} write $\det(A-\lambda I)=0$, solve for $\lambda$, then for each $\lambda$ solve $(A-\lambda I)\vec{v}=\vec{0}$ to find a basis of the eigenspace.
\end{enumerate}

\textbf{Common mistakes}
\begin{itemize}
\item Forgetting to check $T(\vec{0})=\vec{0}$ — translations are not linear.
\item Reversing the order in composition: the matrix of $S\circ T$ is $BA$, not $AB$.
\item Claiming invertibility without checking $\det A\neq 0$ (or $\ker T=\{\vec{0}\}$).
\item Solving $(A-\lambda I)\vec{v}=\vec{0}$ with $\vec{v}=\vec{0}$: the zero vector is \emph{never} an eigenvector.
\item Sign errors in $\det(A-\lambda I)$; expand carefully and factorise.
\item Confusing invariant \emph{points} ($T(\vec{x})=\vec{x}$) with invariant \emph{lines} (eigenspaces with $\lambda\neq 1$).
\item Forgetting that the image of a line may collapse to a point when the direction vector lies in the kernel.
\end{itemize}

\textbf{GCE tips}
\begin{itemize}
\item State definitions precisely — examiners award marks for $T(\lambda\vec{u}+\mu\vec{v})=\lambda T(\vec{u})+\mu T(\vec{v})$.
\item Show the characteristic equation explicitly before solving; quote $\lambda^{2}-(\operatorname{tr}A)\lambda+\det A=0$ for $2\times2$ speed.
\item Always give eigenvectors in a simple form, e.g. $\begin{pmatrix}1\\2\end{pmatrix}$, and state the eigenspace as a span.
\item For images of lines/planes, write parametric forms first, then apply $T$ to each piece.
\item Check answers: verify $A\vec{v}=\lambda\vec{v}$ numerically — a quick safeguard worth marks.
\end{itemize}
\end{aretenir}

\findecours

\end{document}
